% !TEX program = xelatex
% Θεμελιώδεις Αρχές Κυβερνοασφάλειας — Ελληνική έκδοση
\documentclass[11pt,openany]{book}
% ============================================================
%  Preamble — shared by main-en.tex and main-el.tex
%  Compile with XeLaTeX or LuaLaTeX (required for Greek + Unicode):
%     latexmk -xelatex main-el.tex
% ============================================================
\usepackage{fontspec}
\usepackage{polyglossia}
% Fonts with full Greek coverage (TeX Live ships Libertinus; swap freely:
% "GFS Didot", "Noto Serif", "CMU Serif", "Times New Roman" …)
\setmainfont{Libertinus Serif}
\setsansfont{Libertinus Sans}
\setmonofont[Scale=0.9]{DejaVu Sans Mono}
\newfontfamily\greekfont[Script=Greek]{Libertinus Serif}
\newfontfamily\greekfontsf[Script=Greek]{Libertinus Sans}
\newfontfamily\greekfonttt[Script=Greek,Scale=0.9]{DejaVu Sans Mono}

\usepackage[a4paper,margin=2.6cm,headheight=14pt]{geometry}
\usepackage{microtype}
\usepackage{xcolor}
\usepackage{booktabs,tabularx,array}
\usepackage{enumitem}
\usepackage[most]{tcolorbox}
\usepackage{titlesec}
\usepackage{fancyhdr}
\usepackage[font=small,labelfont=bf]{caption}
\usepackage[hidelinks]{hyperref}

\definecolor{accent}{HTML}{0E7490}
\definecolor{accent2}{HTML}{B45309}
\definecolor{ink}{HTML}{0F172A}

\setlength{\parskip}{0.45em}
\setlength{\parindent}{0pt}
\linespread{1.12}
\setlist{itemsep=0.2em,topsep=0.4em}

\titleformat{\chapter}[display]
  {\normalfont\sffamily\color{ink}}
  {\large\color{accent}\chaptertitlename\ \thechapter}
  {0.4em}{\Huge\bfseries}
\titlespacing*{\chapter}{0pt}{-10pt}{18pt}
\titleformat{\section}{\normalfont\sffamily\Large\bfseries\color{ink}}{\color{accent}\thesection}{0.8em}{}

\pagestyle{fancy}
\fancyhf{}
\fancyhead[LE,RO]{\small\thepage}
\fancyhead[RE]{\small\sffamily\leftmark}
\fancyhead[LO]{\small\sffamily\rightmark}
\renewcommand{\headrulewidth}{0.3pt}

% --- chapter metadata line -------------------------------------------------
\newcommand{\chaptermeta}[3]{%
  {\sffamily\small\color{accent}#1\par}\vspace{0.3em}
  {\sffamily\footnotesize\color{gray}#2 \quad·\quad #3\par}\vspace{1.2em}}

% --- boxes ------------------------------------------------------------------
\newtcolorbox{outcomesbox}[1]{enhanced,breakable,colback=accent!5,colframe=accent,
  boxrule=0.6pt,arc=2mm,fonttitle=\sffamily\bfseries,title=#1}
\newtcolorbox{examplebox}[1]{enhanced,breakable,colback=accent2!6,colframe=accent2,
  boxrule=0.5pt,arc=2mm,fonttitle=\sffamily\bfseries,title=#1}
\newtcolorbox{notebox}[1]{enhanced,breakable,colback=gray!6,colframe=gray!60,
  boxrule=0.5pt,arc=2mm,fonttitle=\sffamily\bfseries,coltitle=ink,title=#1}
\newtcolorbox{keybox}[1]{enhanced,breakable,colback=accent!8,colframe=accent!70!black,
  boxrule=0.8pt,arc=2mm,fonttitle=\sffamily\bfseries,title=#1}

\setdefaultlanguage{greek}
\setotherlanguage{english}

\title{Θεμελιώδεις Αρχές Κυβερνοασφάλειας}
\author{Dr. S. Karagiannis}

\begin{document}
\frontmatter
\begin{titlepage}
  \centering\sffamily
  \vspace*{3cm}
  {\color{accent}\rule{\linewidth}{1.2pt}}\par\vspace{1cm}
  {\Huge\bfseries Θεμελιώδεις Αρχές Κυβερνοασφάλειας\par}
  \vspace{0.6cm}
  {\Large Αρχές, Αρχιτεκτονικές και Αμυντικές Λειτουργίες\par}
  \vspace{1cm}
  {\color{accent}\rule{\linewidth}{1.2pt}}\par
  \vfill
  {\large Dr. S. Karagiannis\par}
  \vspace{0.3cm}
  {\small Αναθεωρημένη δίγλωσση έκδοση · 2026\par}
\end{titlepage}
\tableofcontents
\chapter*{Πρόλογος}
\addcontentsline{toc}{chapter}{Πρόλογος}

Το παρόν σύγγραμμα απευθύνεται σε προπτυχιακούς φοιτητές που παρακολουθούν ένα πρώτο πανεπιστημιακό μάθημα κυβερνοασφάλειας, καθώς και σε επαγγελματίες που επιθυμούν να εδραιώσουν την κατανόησή τους για το πεδίο σε στέρεη εννοιολογική βάση. Προϋποθέτει βασική εξοικείωση με υπολογιστές και δίκτυα, όχι όμως προηγούμενη μελέτη της ασφάλειας.

Το βιβλίο οργανώνεται σε δεκατρία κεφάλαια, ομαδοποιημένα σε τέσσερα μέρη. Η σειρά είναι συνειδητή. Ξεκινά από τις έννοιες και τις αρχές, περνά στο λειτουργικό σύστημα ως το πρώτο απτό όριο ασφάλειας και στη συνέχεια μελετά τον αντίπαλο: ποιος επιτίθεται, με ποια εργαλεία και μέσα από ποια κανάλια. Μόνο αφού ο αναγνώστης κατανοήσει τι πρέπει να προστατευθεί, το βιβλίο εισάγει τους μηχανισμούς προστασίας —την κρυπτογραφία, τη διαχείριση ταυτότητας, την ασφαλή μηχανική λογισμικού και τον έλεγχο ασφάλειας. Το τελευταίο μέρος στρέφεται στις λειτουργίες: την ανίχνευση και την απόκριση σε περιστατικά, την ανάλυση κακόβουλου κώδικα και αποδεικτικών στοιχείων και τη διακυβέρνηση της ασφάλειας ως οργανωσιακής λειτουργίας.

Κάθε κεφάλαιο ξεκινά με σύντομη εισαγωγή και με τα μαθησιακά αποτελέσματα, αναπτύσσει το αντικείμενό του σε αριθμημένες ενότητες και ολοκληρώνεται με γλωσσάριο βασικών όρων, σύνοψη και ερωτήσεις ανασκόπησης. Οι πίνακες συμπυκνώνουν συγκρίσεις που γίνονται ευκολότερα κατανοητές όταν παρουσιάζονται παράλληλα, ενώ τα πλαίσια με παραδείγματα συνδέουν τις αφηρημένες έννοιες με ρεαλιστικές καταστάσεις. Οι παραπομπές μεταξύ κεφαλαίων είναι συχνές, επειδή η ασφάλεια είναι ένα διασυνδεδεμένο πεδίο: μια αδυναμία σε ένα επίπεδο συνήθως αξιοποιείται μέσω κάποιου άλλου.

Το βιβλίο εκδίδεται παράλληλα στα αγγλικά και στα ελληνικά. Το ελληνικό κείμενο δεν αποτελεί κατά λέξη μετάφραση· είναι επιμελημένη απόδοση που ακολουθεί τις συμβάσεις του ελληνικού ακαδημαϊκού λόγου και χρησιμοποιεί συνεπή, καθιερωμένη ορολογία. Όπου δεν υπάρχει ευρέως αποδεκτός ελληνικός όρος, ο αγγλικός όρος δίνεται σε παρένθεση κατά την πρώτη εμφάνισή του.


\mainmatter
\part{Θεμέλια}
% ============================================================
% Chapter 1: Foundations of Information Security and the CIA Triad
% Language: Greek — edit freely; regenerate from src/content if preferred.
% ============================================================
\chapter{Θεμελιώδεις Έννοιες της Ασφάλειας Πληροφοριών και η Τριάδα CIA}\label{ch:1}
\chaptermeta{Ορισμοί · Η τριάδα CIA και οι επεκτάσεις της · Αρχές σχεδίασης · Ιστορική εξέλιξη · Μοντέλα ασφάλειας · Κριτήρια αξιολόγησης}{Επίπεδο: Εισαγωγικό επίπεδο · Έμφαση στη θεωρία}{Χρόνος μελέτης: 6–8 ώρες}

Κάθε επόμενο κεφάλαιο του βιβλίου —είτε αφορά λειτουργικά συστήματα, κρυπτογραφία, διαχείριση ταυτότητας είτε απόκριση σε περιστατικά— στηρίζεται σε ένα μικρό σύνολο κοινών εννοιών. Το παρόν κεφάλαιο εισάγει αυτό το λεξιλόγιο. Εξηγεί τι εννοούμε όταν μιλάμε για \emph{ασφάλεια}, ποιες ιδιότητες της πληροφορίας επιδιώκουμε να προστατεύσουμε και ποιες αρχές σχεδίασης εφαρμόζουν οι έμπειροι μηχανικοί πριν επιλέξουν οποιοδήποτε συγκεκριμένο προϊόν ή τεχνολογία.

Η παρουσίαση κινείται από τους ορισμούς στις αρχές και, στη συνέχεια, στα τυπικά μοντέλα. Αρχικά διακρίνουμε την κυβερνοασφάλεια από τις ευρύτερες έννοιες της διασφάλισης πληροφοριών και της κυβερνοανθεκτικότητας. Ακολούθως εξετάζουμε την τριάδα CIA και τις ιδιότητες που τη συμπληρώνουν, τις κλασικές αρχές σχεδίασης των Saltzer και Schroeder, καθώς και τα ιστορικά γεγονότα που διαμόρφωσαν τη σύγχρονη αμυντική σκέψη. Το κεφάλαιο ολοκληρώνεται με τα τυπικά μοντέλα (Bell–LaPadula, Biba, Clark–Wilson) και τα σχήματα αξιολόγησης που επιτρέπουν να διατυπώνονται οι ισχυρισμοί ασφάλειας με ακρίβεια και να επαληθεύονται ανεξάρτητα.

\begin{outcomesbox}{Μαθησιακά αποτελέσματα}
Μετά τη μελέτη του κεφαλαίου θα πρέπει να μπορείτε:
\begin{itemize}
  \item Να ορίζετε την κυβερνοασφάλεια, τη διασφάλιση πληροφοριών και την κυβερνοανθεκτικότητα και να εξηγείτε σε τι διαφέρει το εύρος καθεμιάς.
  \item Να αναλύετε μια απαίτηση ασφάλειας σε εμπιστευτικότητα, ακεραιότητα και διαθεσιμότητα και να αναγνωρίζετε πότε απαιτούνται πρόσθετες ιδιότητες (αυθεντικότητα, λογοδοσία, μη αποποίηση, κατοχή, χρησιμότητα).
  \item Να εφαρμόζετε τις αρχές του ελάχιστου προνομίου, της άμυνας σε βάθος, των ασφαλών προεπιλογών και της ασφάλειας εκ σχεδιασμού κατά την αξιολόγηση μιας αρχιτεκτονικής.
  \item Να διακρίνετε την απειλή, την ευπάθεια και τον κίνδυνο και να περιγράφετε πώς τα μέτρα ελέγχου μειώνουν την πιθανότητα ή τις συνέπειες.
  \item Να εξηγείτε τον σκοπό των τυπικών μοντέλων ασφάλειας και των ανεξάρτητων σχημάτων αξιολόγησης, όπως τα Common Criteria.
\end{itemize}
\end{outcomesbox}

\section{Ορισμοί: κυβερνοασφάλεια, διασφάλιση πληροφοριών και κυβερνοανθεκτικότητα}\label{sec:1.1}
Η \textbf{κυβερνοασφάλεια} είναι ο επιστημονικός και επαγγελματικός κλάδος που προστατεύει δικτυωμένα συστήματα, δεδομένα και υπηρεσίες από μη εξουσιοδοτημένη πρόσβαση, τροποποίηση, διακοπή και καταστροφή, διασφαλίζοντας ταυτόχρονα τη νόμιμη λειτουργία τους ακόμη και υπό εχθρικές συνθήκες. Η φράση \emph{υπό εχθρικές συνθήκες} έχει ιδιαίτερη σημασία: σε αντίθεση με τη μηχανική αξιοπιστίας, η οποία αντιμετωπίζει τυχαία σφάλματα, η ασφάλεια προϋποθέτει έναν ευφυή αντίπαλο που αναζητά συστηματικά το πιο αδύναμο σημείο.

Η \textbf{διασφάλιση πληροφοριών} (Information Assurance, IA) έχει ευρύτερο πεδίο. Συνδυάζει τεχνικά, διαδικαστικά και διοικητικά μέτρα, ώστε να διατηρείται \emph{τεκμηριωμένη εμπιστοσύνη} στην πληροφορία σε όλο τον κύκλο ζωής της —από τη δημιουργία και την αποθήκευση έως τη μετάδοση, τη χρήση και την τελική καταστροφή της. Επομένως, η διασφάλιση πληροφοριών δεν αφορά μόνο την αντίσταση σε επιθέσεις, αλλά και τη γνησιότητα, την προέλευση και τη διαρκή χρησιμότητα της πληροφορίας.

Η \textbf{κυβερνοανθεκτικότητα} διευρύνει και τις δύο έννοιες. Περιγράφει την ικανότητα ενός οργανισμού να προβλέπει, να αντέχει, να ανακάμπτει και να προσαρμόζεται σε δυσμενή κυβερνογεγονότα, συμπεριλαμβανομένων εκείνων που επιτυγχάνουν παρά τα προληπτικά μέτρα. Με απλά λόγια, η κυβερνοασφάλεια ρωτά \emph{πώς αποτρέπουμε και εντοπίζουμε μια παραβίαση;}, η διασφάλιση πληροφοριών \emph{μπορούμε να εμπιστευτούμε την πληροφορία μας;} και η ανθεκτικότητα \emph{μπορούμε να συνεχίσουμε να λειτουργούμε και πόσο γρήγορα θα ανακάμψουμε;}

\begin{table}[htbp]
  \centering\small
  \caption{Τρεις συγγενείς αλλά διακριτές έννοιες}
  \begin{tabularx}{\linewidth}{>{\raggedright\arraybackslash}X >{\raggedright\arraybackslash}X >{\raggedright\arraybackslash}X >{\raggedright\arraybackslash}X}
    \toprule
    \textbf{Έννοια} & \textbf{Βασικό ερώτημα} & \textbf{Τυπικά μέτρα} & \textbf{Συνέπεια αν παραμεληθεί} \\
    \midrule
    Κυβερνοασφάλεια & Μπορεί ένας αντίπαλος να παραβιάσει ή να διαταράξει το σύστημα; & Τείχη προστασίας, IDS/IPS, ενημερώσεις, MFA, κρυπτογράφηση & Εισβολή, διαρροή δεδομένων, διακοπή υπηρεσίας \\
    Διασφάλιση πληροφοριών & Μπορούμε να εμπιστευτούμε την πληροφορία; & Κατακερματισμός, ψηφιακές υπογραφές, προέλευση, διαβάθμιση & Αθόρυβη αλλοίωση, αποποίηση ευθύνης, παραπληροφόρηση \\
    Κυβερνοανθεκτικότητα & Μπορούμε να λειτουργούμε παρά την αστοχία; & Αντίγραφα ασφαλείας, εφεδρεία, σχέδια απόκρισης, BCP/DR & Παρατεταμένη διακοπή, μη ανακτήσιμη απώλεια \\
    \bottomrule
  \end{tabularx}
\end{table}
\section{Η τριάδα CIA και οι επεκτάσεις της}\label{sec:1.2}
Η \textbf{τριάδα CIA} —Εμπιστευτικότητα (Confidentiality), Ακεραιότητα (Integrity) και Διαθεσιμότητα (Availability)— αποτελεί τον καθιερωμένο τρόπο ανάλυσης μιας απαίτησης ασφάλειας στους βασικούς της στόχους. Η \textbf{εμπιστευτικότητα} περιορίζει την αποκάλυψη της πληροφορίας μόνο σε εξουσιοδοτημένα μέρη· επιβάλλεται συνήθως μέσω ελέγχου πρόσβασης και κρυπτογράφησης. Η \textbf{ακεραιότητα} προστατεύει την πληροφορία από μη εξουσιοδοτημένη ή μη ανιχνεύσιμη τροποποίηση, ώστε τα δεδομένα να παραμένουν ακριβή και πλήρη· σε αυτόν τον στόχο συμβάλλουν οι συναρτήσεις κατακερματισμού, οι ψηφιακές υπογραφές, ο έλεγχος εκδόσεων και οι επικυρωμένες συναλλαγές. Η \textbf{διαθεσιμότητα} εξασφαλίζει έγκαιρη και αξιόπιστη πρόσβαση στους εξουσιοδοτημένους χρήστες και υποστηρίζεται από πλεονασμό, σχεδιασμό χωρητικότητας, περιορισμό ρυθμού αιτημάτων και διαδικασίες ανάκαμψης.

Ορισμένες πρόσθετες ιδιότητες εξειδικεύουν την τριάδα χωρίς να την αντικαθιστούν. Η \textbf{αυθεντικότητα} βεβαιώνει ότι ένα δεδομένο ή ένα μέρος είναι πράγματι αυτό που ισχυρίζεται ότι είναι. Η \textbf{λογοδοσία} συνδέει κάθε ενέργεια με ένα αναγνωρίσιμο υποκείμενο μέσω αξιόπιστων αρχείων καταγραφής. Η \textbf{μη αποποίηση} παρέχει αποδεικτικά στοιχεία που καθιστούν δύσκολη τη μεταγενέστερη άρνηση μιας ενέργειας· η ισχύς της, ωστόσο, εξαρτάται τελικά από τη φύλαξη των κλειδιών, την ταυτοποίηση των προσώπων και το νομικό πλαίσιο.

Οι τρεις στόχοι αλληλεπιδρούν και συχνά ανταγωνίζονται μεταξύ τους. Η κρυπτογράφηση ενισχύει την εμπιστευτικότητα, όμως αν χαθούν τα κλειδιά τα δεδομένα καθίστανται μη διαθέσιμα. Ο πλεονασμός βελτιώνει τη διαθεσιμότητα, αλλά κάθε επιπλέον αντίγραφο είναι ένα ακόμη σύστημα που πρέπει να προστατεύεται και να διατηρείται συνεπές. Μια τεκμηριωμένη αρχιτεκτονική δηλώνει, επομένως, ποιους στόχους υπηρετεί κάθε μέτρο, εντοπίζει τυχόν νέα έκθεση που αυτό εισάγει και εξηγεί πώς θα αντιμετωπιστεί ο κίνδυνος που απομένει (υπολειπόμενος κίνδυνος).

\begin{examplebox}{Παράδειγμα εφαρμογής}
Ένα νοσοκομείο κρυπτογραφεί τη βάση δεδομένων ασθενών (εμπιστευτικότητα) και τη συγχρονίζει με δεύτερη εγκατάσταση (διαθεσιμότητα). Αν το μοναδικό αντίγραφο του κλειδιού αποκρυπτογράφησης βρίσκεται στον κύριο διακομιστή και αυτός καταστραφεί από λυτρισμικό, το αντίγραφο της βάσης είναι άχρηστο: η εμπιστευτικότητα διατηρήθηκε εις βάρος της διαθεσιμότητας. Το μέτρο που συμφιλιώνει τους δύο στόχους είναι η ορθή διαχείριση κλειδιών.
\end{examplebox}

\section{Αρχές σχεδίασης που διέπουν όλα τα επόμενα κεφάλαια}\label{sec:1.3}
Τέσσερις αρχές επανέρχονται σε όλο το βιβλίο ως αρχιτεκτονικοί περιορισμοί. Η \textbf{ασφάλεια εκ σχεδιασμού} (security by design) απαιτεί οι ιδιότητες ασφάλειας να προσδιορίζονται, να μοντελοποιούνται, να υλοποιούνται και να επαληθεύονται από την αρχή του κύκλου ζωής και όχι να προστίθενται μετά την παραγωγική λειτουργία, όταν οι αλλαγές είναι δαπανηρές και αποσπασματικές. Η \textbf{άμυνα σε βάθος} συνδυάζει ανεξάρτητα και ποικίλα προληπτικά, ανιχνευτικά και διορθωτικά μέτρα, ώστε η αστοχία ενός μεμονωμένου επιπέδου να μην κρίνει την έκβαση. Η \textbf{αρχή του ελάχιστου προνομίου} παραχωρεί σε κάθε χρήστη, διεργασία ή υπηρεσία μόνο την εξουσία που απαιτείται για συγκεκριμένη εργασία και για συγκεκριμένο χρονικό διάστημα, και την ανακαλεί όταν δεν είναι πλέον αναγκαία. Οι \textbf{ασφαλείς προεπιλογές} (fail-safe defaults) απορρίπτουν την πρόσβαση εκτός αν έχει ρητά επιτραπεί και απαιτούν από τα συστατικά ενός συστήματος να καταλήγουν, σε περίπτωση αστοχίας, σε γνωστή ασφαλή κατάσταση.

Οι αρχές αυτές προέρχονται από το θεμελιώδες άρθρο των Saltzer και Schroeder (1975), το οποίο εισήγαγε επίσης την \emph{οικονομία μηχανισμού} (απλοί και μικροί σχεδιασμοί, ώστε να μπορούν να ελεγχθούν), την \emph{πλήρη διαμεσολάβηση} (έλεγχος κάθε πρόσβασης και όχι μόνο της πρώτης), τον \emph{ανοιχτό σχεδιασμό} (η ασφάλεια δεν πρέπει να εξαρτάται από τη μυστικότητα του σχεδιασμού), τον \emph{διαχωρισμό προνομίων} (περισσότερες από μία ανεξάρτητες προϋποθέσεις για ευαίσθητες ενέργειες), τον \emph{ελάχιστο κοινό μηχανισμό} (περιορισμός των κοινόχρηστων στοιχείων μεταξύ χρηστών) και την \emph{ψυχολογική αποδοχή} (η ασφάλεια πρέπει να είναι εύχρηστη, διαφορετικά οι χρήστες θα την παρακάμψουν).

Οι αρχές είναι συμπληρωματικές και πρέπει να εφαρμόζονται από κοινού. Η πολυεπίπεδη άμυνα χωρίς ελάχιστο προνόμιο απλώς πολλαπλασιάζει τους ισχυρούς λογαριασμούς που μπορεί να υποκλέψει ένας επιτιθέμενος· μια αυστηρή πολιτική άρνησης εξ ορισμού χωρίς λειτουργικές διαδικασίες ωθεί τους χρήστες σε μη ασφαλείς παρακάμψεις. Μια χρήσιμη επαγγελματική πρακτική είναι να κατονομάζεται η αρχή που παραβιάστηκε σε κάθε ανασκόπηση μετά από περιστατικό· για παράδειγμα, μια τραπεζική μεταφορά που εγκρίνεται από ένα μόνο πρόσωπο παραβιάζει τον διαχωρισμό προνομίων.

\section{Ιστορική εξέλιξη: από το Creeper στον υπερσυνδεδεμένο οργανισμό}\label{sec:1.4}
Η σύγχρονη αρχιτεκτονική ασφάλειας γίνεται καλύτερα κατανοητή ως συσσώρευση διδαγμάτων και όχι ως αλληλουχία απειλών που η μία αντικατέστησε την άλλη. Τα πρώιμα πειράματα στο ARPANET, όπως το \emph{Creeper} και το αντίστοιχο \emph{Reaper} (αρχές της δεκαετίας του 1970), έδειξαν τόσο τη δυνατότητα αυτοδιαδιδόμενου κώδικα όσο και την αυτοματοποιημένη αφαίρεσή του. Το \emph{Morris Worm} του 1988 ανέδειξε πώς ένα μεμονωμένο σφάλμα λογισμικού μπορεί να εξαπλωθεί αλυσιδωτά σε διασυνδεδεμένους υπολογιστές και οδήγησε άμεσα στη δημιουργία της πρώτης Ομάδας Απόκρισης σε Υπολογιστικά Περιστατικά (CERT).

Στη συνέχεια, το εμπορικό Διαδίκτυο και ο Παγκόσμιος Ιστός διεύρυναν δραματικά την επιφάνεια επίθεσης: οι απομακρυσμένα προσβάσιμες υπηρεσίες, η έγχυση SQL, το cross-site scripting, τα botnet και το οργανωμένο κυβερνοέγκλημα έγιναν καθημερινά ζητήματα. Τις επόμενες δεκαετίες, το υπολογιστικό νέφος, οι κινητές πλατφόρμες, το λυτρισμικό ως υπηρεσία, οι επιθέσεις με επίκεντρο την ταυτότητα, οι παραβιάσεις της εφοδιαστικής αλυσίδας λογισμικού και η σύγκλιση της πληροφορικής με την επιχειρησιακή τεχνολογία (OT) εισήγαγαν νέα διοικητικά πεδία και νέες εξαρτήσεις.

Κάθε εποχή στηρίχθηκε σε μια παραδοχή εμπιστοσύνης που οι επιτιθέμενοι τελικά διέψευσαν. Ο σημερινός οργανισμός είναι υβριδικός, έντονα διασυνδεδεμένος και εξαρτημένος από εξωτερικούς παρόχους ταυτότητας, λογισμικού και υποδομών. Το πρακτικό συμπέρασμα είναι ότι η θέση μέσα στο δίκτυο δεν αρκεί πλέον για να θεμελιώσει εμπιστοσύνη· η ταυτότητα, η κατάσταση της συσκευής και η εξουσιοδότηση πρέπει να επαληθεύονται συνεχώς. Αυτή είναι η βασική ιδέα της αρχιτεκτονικής μηδενικής εμπιστοσύνης (Zero Trust), την οποία εξετάζει αναλυτικά το Κεφάλαιο 13.

\begin{table}[htbp]
  \centering\small
  \caption{Παραδοχές εμπιστοσύνης και οι αστοχίες που τις διέψευσαν}
  \begin{tabularx}{\linewidth}{>{\raggedright\arraybackslash}X >{\raggedright\arraybackslash}X >{\raggedright\arraybackslash}X >{\raggedright\arraybackslash}X}
    \toprule
    \textbf{Περίοδος} & \textbf{Παραδοχή εμπιστοσύνης} & \textbf{Χαρακτηριστική αστοχία} & \textbf{Καινοτομία ελέγχου που προέκυψε} \\
    \midrule
    1970–1980 & Αξιόπιστοι υπολογιστές και χρήστες & Creeper, Morris Worm & Έλεγχος πρόσβασης, CERT \\
    1990–2000 & Ασφαλές εσωτερικό, εχθρικό εξωτερικό & Code Red, SQL Slammer & Τείχη προστασίας, διαχείριση ενημερώσεων, antivirus \\
    2010–2020 & Αξιόπιστοι προμηθευτές, περίμετρος VPN & Παραβίαση Target, WannaCry & Κατάτμηση δικτύου, EDR, MFA \\
    2020 κ.ε. & Ποτέ μην εμπιστεύεσαι, πάντα επαλήθευε & SolarWinds, λυτρισμικό σε OT & Zero Trust, SBOM, XDR \\
    \bottomrule
  \end{tabularx}
\end{table}
\section{Απειλή, ευπάθεια και κίνδυνος: το βασικό λεξιλόγιο}\label{sec:1.5}
Τρεις όροι χρησιμοποιούνται συχνά εναλλακτικά στην καθημερινή γλώσσα, έχουν όμως διακριτό τεχνικό περιεχόμενο. \textbf{Απειλή} είναι κάθε περίσταση, γεγονός ή δράστης που μπορεί να προκαλέσει βλάβη —μια εγκληματική ομάδα, ένας δυσαρεστημένος υπάλληλος, μια πλημμύρα. \textbf{Ευπάθεια} είναι μια αδυναμία που μπορεί να εκμεταλλευτεί ή να ενεργοποιήσει μια απειλή —μια υπηρεσία χωρίς ενημερώσεις, ένας αδύναμος κωδικός πρόσβασης, η απουσία αντιγράφων ασφαλείας. \textbf{Κίνδυνος} είναι η δυνατότητα δυσμενούς επίπτωσης υπό συνθήκες αβεβαιότητας· υπάρχει μόνο όταν μια σχετική απειλή συναντά μια εκμεταλλεύσιμη ευπάθεια σε ένα περιουσιακό στοιχείο που έχει αξία.

Η γνωστή έκφραση \emph{κίνδυνος ≈ πιθανότητα × επίπτωση} είναι χρήσιμο μοντέλο ιεράρχησης και όχι φυσικός νόμος. Κάθε οργανισμός οφείλει να ορίζει τις δικές του κλίμακες, να δηλώνει τις παραδοχές του και να αναγνωρίζει την αβεβαιότητα. \textbf{Μέτρο ελέγχου} είναι κάθε ενέργεια που μεταβάλλει την πιθανότητα, την επίπτωση ή και τα δύο: ο έλεγχος ταυτότητας πολλαπλών παραγόντων μειώνει την πιθανότητα κατάληψης λογαριασμού, ενώ τα αμετάβλητα αντίγραφα ασφαλείας μειώνουν την επίπτωση ενός λυτρισμικού.

Μετά την εφαρμογή των μέτρων απομένει ο \textbf{υπολειπόμενος κίνδυνος}, ο οποίος πρέπει να αντιμετωπίζεται ρητά. Υπάρχουν τέσσερις επιλογές: η περαιτέρω \emph{μείωση}, η \emph{αποδοχή} (από στέλεχος που έχει τη σχετική αρμοδιότητα), η \emph{μεταφορά} όπου αυτό έχει νόημα (για παράδειγμα, μέσω ασφάλισης ή συμβάσεων) και η \emph{αποφυγή}, δηλαδή η εγκατάλειψη της επικίνδυνης δραστηριότητας. Τα ώριμα προγράμματα ασφάλειας αξιολογούν το σύνολο των μέτρων τους με βάση την κάλυψη, την ανεξαρτησία και την πραγματική λειτουργική αποτελεσματικότητα —όχι με βάση τον αριθμό των προϊόντων που έχουν εγκαταστήσει.

\section{Η εξάδα του Parker και ο κύβος του McCumber}\label{sec:1.6}
Η τριάδα CIA είναι αναγκαία, αλλά δεν είναι πάντοτε επαρκής. Ο Donn Parker πρότεινε την \textbf{εξάδα του Parker} (Parkerian Hexad), η οποία προσθέτει στις τρεις κλασικές ιδιότητες την \emph{αυθεντικότητα}, την \emph{κατοχή (έλεγχο)} και τη \emph{χρησιμότητα}. Οι προσθήκες αυτές αποτυπώνουν καταστάσεις που η τριάδα περιγράφει ελλιπώς. Όταν κλαπεί ένας κρυπτογραφημένος φορητός υπολογιστής, η εμπιστευτικότητα μπορεί να παραμένει άθικτη, όμως η κατοχή έχει σαφώς απολεσθεί. Όταν τα κρυπτογραφημένα δεδομένα διασώζονται αλλά τα κλειδιά έχουν χαθεί οριστικά, τα δεδομένα παραμένουν εμπιστευτικά, έχουν όμως χάσει κάθε χρησιμότητα.

Ο \textbf{κύβος του McCumber} προσφέρει μια συμπληρωματική, τρισδιάστατη οπτική. Εξετάζει την ασφάλεια στις τρεις \emph{καταστάσεις} της πληροφορίας (αποθήκευση, επεξεργασία, μετάδοση), στις τρεις \emph{κατηγορίες μέτρων προστασίας} (τεχνολογία, πολιτικές και πρακτικές, άνθρωποι) και στους ίδιους τους στόχους ασφάλειας. Η συνδυαστική χρήση των δύο μοντέλων αποκαλύπτει ελλιπείς ισχυρισμούς όπως «τα δεδομένα είναι κρυπτογραφημένα, άρα είναι ασφαλή». Ο αξιολογητής καλείται να ρωτήσει: ποιες ιδιότητες προστατεύονται, σε ποια κατάσταση της πληροφορίας και με ποιον συνδυασμό τεχνολογίας, διαδικασιών και ανθρώπινης συμπεριφοράς;

\section{Τυπικά μοντέλα ασφάλειας: Bell–LaPadula, Biba και Clark–Wilson}\label{sec:1.7}
Οι αρχές μάς λένε \emph{τι} επιδιώκουμε· τα τυπικά μοντέλα ορίζουν \emph{με ακρίβεια} ποιες συμπεριφορές ενός συστήματος επιτρέπονται, ώστε ένας σχεδιασμός να μπορεί να αναλυθεί και, ιδανικά, να αποδειχθεί ορθός. Το μοντέλο \textbf{Bell–LaPadula (BLP)}, που αναπτύχθηκε για στρατιωτικά συστήματα τη δεκαετία του 1970, προστατεύει την εμπιστευτικότητα με δύο κανόνες: \emph{όχι ανάγνωση προς τα πάνω} (ένα υποκείμενο δεν μπορεί να διαβάσει αντικείμενα υψηλότερης διαβάθμισης) και \emph{όχι εγγραφή προς τα κάτω} (ένα υποκείμενο δεν μπορεί να γράψει πληροφορία σε χαμηλότερη διαβάθμιση, κάτι που θα διέρρεε απόρρητα).

Το μοντέλο \textbf{Biba} αποτελεί το αντίστοιχο του BLP για την ακεραιότητα και αντιστρέφει τους κανόνες του: \emph{όχι ανάγνωση προς τα κάτω} (μια αξιόπιστη διεργασία δεν πρέπει να καταναλώνει λιγότερο αξιόπιστα δεδομένα) και \emph{όχι εγγραφή προς τα πάνω} (ένα μη αξιόπιστο υποκείμενο δεν μπορεί να τροποποιεί πιο αξιόπιστα αντικείμενα). Στα εμπορικά περιβάλλοντα, ωστόσο, ενδιαφέρουν λιγότερο τα επίπεδα διαβάθμισης και περισσότερο η ορθότητα των συναλλαγών. Γι' αυτό το μοντέλο \textbf{Clark–Wilson} επιβάλλει την ακεραιότητα μέσω \emph{καλά σχηματισμένων συναλλαγών}, οι οποίες εκτελούνται μόνο από εξουσιοδοτημένους χρήστες μέσω πιστοποιημένων προγραμμάτων, σε συνδυασμό με τον \emph{διαχωρισμό καθηκόντων} και την πλήρη καταγραφή ελέγχου.

\begin{notebox}{Γιατί έχει πρακτική σημασία}
Τα υποχρεωτικά επίπεδα ακεραιότητας των Windows (Κεφάλαιο 2) προέρχονται άμεσα από το μοντέλο Biba, ενώ η έγκριση τραπεζικών μεταφορών από δύο πρόσωπα —δημιουργό και ελεγκτή (Κεφάλαιο 5)— αποτελεί κλασικό μέτρο τύπου Clark–Wilson.
\end{notebox}

\section{Κριτήρια αξιολόγησης: από το TCSEC στα Common Criteria}\label{sec:1.8}
Πώς μπορεί ένας πελάτης να γνωρίζει αν ένα προϊόν παρέχει πράγματι την ασφάλεια που ισχυρίζεται ο κατασκευαστής του; Τα σχήματα αξιολόγησης απαντούν σε αυτό το ερώτημα μέσω ανεξάρτητης και επαναλήψιμης αξιολόγησης. Τα αμερικανικά \textbf{Trusted Computer System Evaluation Criteria (TCSEC)}, γνωστά ως \emph{Πορτοκαλί Βιβλίο} (Orange Book, 1983), κατέτασσαν τα λειτουργικά συστήματα σε κατηγορίες από D (ελάχιστη προστασία) έως A1 (τυπικά επαληθευμένος σχεδιασμός). Η Ευρώπη ακολούθησε με το \textbf{ITSEC}, το οποίο διαχώρισε τη λειτουργικότητα από τη διασφάλιση.

Και τα δύο αντικαταστάθηκαν από τα διεθνή \textbf{Common Criteria (ISO/IEC 15408)}. Στο πλαίσιο αυτό, ένα \emph{Προφίλ Προστασίας} (Protection Profile) περιγράφει τις απαιτήσεις ασφάλειας για μια κατηγορία προϊόντων (για παράδειγμα, τείχη προστασίας), ενώ ο \emph{Στόχος Ασφάλειας} (Security Target) του κατασκευαστή δηλώνει πώς ένα συγκεκριμένο προϊόν τις ικανοποιεί. Διαπιστευμένα εργαστήρια αξιολογούν στη συνέχεια το προϊόν σε ένα \textbf{Επίπεδο Διασφάλισης Αξιολόγησης (EAL1–EAL7)}. Είναι κρίσιμο να κατανοηθεί τι σημαίνει το EAL: μετρά την \emph{αυστηρότητα} της αξιολόγησης και όχι την \emph{ισχύ} του προϊόντος. Ένα προϊόν EAL4 έχει εξεταστεί διεξοδικότερα από ένα προϊόν EAL2, αλλά μόνο ως προς τις απαιτήσεις που δήλωσε ο δικός του Στόχος Ασφάλειας.

\section*{Βασικοί όροι}
\begin{description}[style=nextline,leftmargin=1.2em]
  \item[Τριάδα CIA] Εμπιστευτικότητα, ακεραιότητα και διαθεσιμότητα: οι τρεις βασικοί στόχοι της ασφάλειας πληροφοριών.
  \item[Μη αποποίηση] Αποδεικτικά στοιχεία που εμποδίζουν ένα μέρος να αρνηθεί αξιόπιστα μια ενέργεια που εκτέλεσε.
  \item[Ελάχιστο προνόμιο] Παραχώρηση μόνο της ελάχιστης αναγκαίας εξουσίας για μια εργασία και για τον ελάχιστο αναγκαίο χρόνο.
  \item[Άμυνα σε βάθος] Διαστρωμάτωση ανεξάρτητων και ποικίλων μέτρων, ώστε καμία μεμονωμένη αστοχία να μην είναι καθοριστική.
  \item[Υπολειπόμενος κίνδυνος] Ο κίνδυνος που απομένει μετά την εφαρμογή των μέτρων και πρέπει να αντιμετωπιστεί ρητά.
  \item[EAL] Επίπεδο Διασφάλισης Αξιολόγησης (Common Criteria): η αυστηρότητα μιας ανεξάρτητης αξιολόγησης προϊόντος.
\end{description}

\section*{Σύνοψη κεφαλαίου}
\begin{itemize}
  \item Η κυβερνοασφάλεια, η διασφάλιση πληροφοριών και η κυβερνοανθεκτικότητα απαντούν σε διαφορετικά ερωτήματα: την αποτροπή της παραβίασης, την εμπιστοσύνη στην πληροφορία και τη συνέχιση της λειτουργίας παρά την αστοχία.
  \item Η τριάδα CIA παραμένει το θεμέλιο, όμως η αυθεντικότητα, η λογοδοσία, η μη αποποίηση, η κατοχή και η χρησιμότητα είναι απαραίτητες για την περιγραφή πολλών πραγματικών καταστάσεων.
  \item Τα μέτρα ασφάλειας επιλέγονται πρώτα με βάση τις αρχές —ελάχιστο προνόμιο, άμυνα σε βάθος, ασφαλείς προεπιλογές, ασφάλεια εκ σχεδιασμού— και μόνο έπειτα με βάση τα προϊόντα.
  \item Ο κίνδυνος προκύπτει όταν μια απειλή συναντά μια ευπάθεια σε ένα πολύτιμο περιουσιακό στοιχείο· ο υπολειπόμενος κίνδυνος πρέπει πάντοτε να μειώνεται, να γίνεται αποδεκτός, να μεταφέρεται ή να αποφεύγεται ρητά.
  \item Τα τυπικά μοντέλα και τα σχήματα αξιολόγησης καθιστούν τους ισχυρισμούς ασφάλειας ακριβείς και ανεξάρτητα επαληθεύσιμους.
\end{itemize}

\section*{Ερωτήσεις ανασκόπησης}
\begin{enumerate}
  \item Εξηγήστε, με δικό σας παράδειγμα, γιατί η κυβερνοανθεκτικότητα παραμένει αναγκαία ακόμη και σε έναν οργανισμό με εξαιρετικά προληπτικά μέτρα.
  \item Περιγράψτε ένα σενάριο στο οποίο η βελτίωση της διαθεσιμότητας αποδυναμώνει την εμπιστευτικότητα. Πώς θα συμφιλιώνατε τους δύο στόχους;
  \item Ποια αρχή των Saltzer και Schroeder παραβιάζεται όταν όλοι οι διαχειριστές μοιράζονται έναν κοινό κωδικό διαχειριστή τομέα; Τεκμηριώστε την απάντησή σας.
  \item Γιατί η διαβάθμιση EAL των Common Criteria δεν αρκεί από μόνη της για να κρίνετε αν ένα προϊόν είναι κατάλληλο για το δικό σας περιβάλλον;
\end{enumerate}

% ============================================================
% Chapter 2: Operating System Security: POSIX and Windows Architecture
% Language: Greek — edit freely; regenerate from src/content if preferred.
% ============================================================
\chapter{Ασφάλεια Λειτουργικών Συστημάτων: Αρχιτεκτονική POSIX και Windows}\label{ch:2}
\chaptermeta{Διακριτικός έλεγχος πρόσβασης · Ειδικά δικαιώματα · ACL και capabilities · Υποχρεωτικοί έλεγχοι · Token, UAC και επίπεδα ακεραιότητας στα Windows}{Επίπεδο: Εισαγωγικό–Μεσαίο επίπεδο · Συστήματα}{Χρόνος μελέτης: 7–9 ώρες}

Το λειτουργικό σύστημα είναι το πρώτο απτό όριο ασφάλειας που συναντά ο φοιτητής. Κάθε αρχείο, κάθε διεργασία και κάθε δικτυακή υποδοχή (socket) ελέγχεται τελικά από τον πυρήνα, ο οποίος αποφασίζει ποιος μπορεί να κάνει τι. Αν ο έλεγχος πρόσβασης του λειτουργικού συστήματος είναι λανθασμένα διαμορφωμένος, οι άμυνες υψηλότερου επιπέδου, όπως τα τείχη προστασίας ή το λογισμικό προστασίας από ιούς, συχνά παρακάμπτονται με ελάχιστη προσπάθεια.

Το κεφάλαιο συγκρίνει τις δύο κυρίαρχες οικογένειες ελέγχου πρόσβασης. Αρχικά εξηγεί το μοντέλο POSIX που χρησιμοποιούν το Linux, το macOS και τα υπόλοιπα συστήματα τύπου Unix, συμπεριλαμβανομένων των ειδικών δικαιωμάτων που αποτελούν συχνή αιτία κλιμάκωσης προνομίων. Στη συνέχεια προχωρά πέρα από τα απλά δικαιώματα, στις λίστες ελέγχου πρόσβασης (ACL), στις δυνατότητες (capabilities) του Linux και στα πλαίσια υποχρεωτικού ελέγχου πρόσβασης. Το δεύτερο μισό αφιερώνεται στα Windows: αναγνωριστικά ασφάλειας, διακριτικά πρόσβασης (access tokens), Έλεγχος Λογαριασμού Χρήστη (UAC), το μητρώο (registry) και τα υποχρεωτικά επίπεδα ακεραιότητας.

\begin{outcomesbox}{Μαθησιακά αποτελέσματα}
Μετά τη μελέτη του κεφαλαίου θα πρέπει να μπορείτε:
\begin{itemize}
  \item Να διαβάζετε, να ερμηνεύετε και να ορίζετε δικαιώματα POSIX σε συμβολική και οκταδική μορφή.
  \item Να εξηγείτε τον σκοπό και τους κινδύνους των ειδικών δικαιωμάτων SetUID, SetGID και sticky bit.
  \item Να περιγράφετε πώς οι ACL, οι δυνατότητες (capabilities) και ο υποχρεωτικός έλεγχος πρόσβασης (SELinux, AppArmor) εξειδικεύουν το βασικό μοντέλο δικαιωμάτων.
  \item Να εξηγείτε τον ρόλο των SID, των διακριτικών πρόσβασης, του UAC και των επιπέδων ακεραιότητας στα Windows.
  \item Να εντοπίζετε συνήθεις εσφαλμένες ρυθμίσεις λειτουργικών συστημάτων που επιτρέπουν κλιμάκωση προνομίων ή πλευρική κίνηση.
\end{itemize}
\end{outcomesbox}

\section{Διακριτικός έλεγχος πρόσβασης POSIX: μοντέλο και λειτουργία}\label{sec:2.1}
Τα συστήματα POSIX εφαρμόζουν \textbf{διακριτικό έλεγχο πρόσβασης} (Discretionary Access Control, DAC): ο κάτοχος ενός αρχείου αποφασίζει ποιοι άλλοι μπορούν να έχουν πρόσβαση σε αυτό. Κάθε αρχείο έχει έναν κάτοχο (αναγνωριστικό χρήστη), μια ομάδα (αναγνωριστικό ομάδας) και τρία σύνολα δικαιωμάτων —για τον \emph{κάτοχο}, την \emph{ομάδα} και τους \emph{λοιπούς}. Κάθε σύνολο περιλαμβάνει τρία δικαιώματα: \textbf{ανάγνωση (r)}, \textbf{εγγραφή (w)} και \textbf{εκτέλεση (x)}. Στους καταλόγους η σημασία διαφοροποιείται ελαφρώς: η ανάγνωση επιτρέπει την εμφάνιση των ονομάτων, η εγγραφή τη δημιουργία ή διαγραφή εγγραφών και η εκτέλεση την είσοδο στον κατάλογο και την πρόσβαση στα αρχεία που περιέχει.

Τα δικαιώματα γράφονται συνήθως σε οκταδική μορφή, όπου η ανάγνωση αντιστοιχεί στο 4, η εγγραφή στο 2 και η εκτέλεση στο 1. Επομένως, η τιμή \texttt{640} σημαίνει \emph{rw-} για τον κάτοχο, \emph{r--} για την ομάδα και \emph{---} για όλους τους υπόλοιπους —μια τυπική ρύθμιση για αρχείο ρυθμίσεων που περιέχει μυστικά. Όταν μια διεργασία ζητά πρόσβαση, ο πυρήνας ελέγχει τις κατηγορίες με σειρά (κάτοχος, ομάδα, λοιποί) και εφαρμόζει μόνο την \emph{πρώτη} κατηγορία που ταιριάζει. Κατά συνέπεια, ένας κάτοχος που έχει αφαιρέσει από τον εαυτό του το δικαίωμα ανάγνωσης δεν αποκτά πρόσβαση, ακόμη κι αν οι «λοιποί» μπορούν να διαβάσουν το αρχείο.

\begin{examplebox}{Παράδειγμα εφαρμογής}
Η εντολή \texttt{chmod 750 /srv/reports} παρέχει στον κάτοχο πλήρη έλεγχο (7 = 4+2+1), επιτρέπει στα μέλη της ομάδας να βλέπουν και να εισέρχονται στον κατάλογο (5 = 4+1) και απαγορεύει κάθε πρόσβαση στους υπόλοιπους χρήστες (0). Ένας διακομιστής ιστού που εκτελείται με άσχετο λογαριασμό δεν θα μπορεί, επομένως, να διαβάσει τις αναφορές, ακόμη κι αν μια ευπάθεια του επιτρέπει να μαντέψει τη διαδρομή τους.
\end{examplebox}

\section{Ειδικά δικαιώματα: SetUID, SetGID και sticky bit}\label{sec:2.2}
Τρία πρόσθετα bit τροποποιούν τη συμπεριφορά των βασικών δικαιωμάτων. Το \textbf{SetUID} (οκταδικό 4000) προκαλεί την εκτέλεση ενός προγράμματος με τα προνόμια του \emph{κατόχου} του και όχι του χρήστη που το εκκίνησε. Με αυτόν τον τρόπο ένας απλός χρήστης μπορεί να αλλάξει τον κωδικό του με την εντολή \texttt{passwd}, ένα πρόγραμμα που ανήκει στον root και πρέπει να γράψει στο προστατευμένο αρχείο \texttt{/etc/shadow}. Το \textbf{SetGID} (2000) λειτουργεί ανάλογα για την ομάδα· όταν εφαρμόζεται σε κατάλογο, κάνει επιπλέον τα νέα αρχεία να κληρονομούν την ομάδα του καταλόγου, κάτι χρήσιμο για κοινόχρηστους φακέλους έργων. Το \textbf{sticky bit} (1000) σε κατάλογο με δικαίωμα εγγραφής για όλους, όπως ο \texttt{/tmp}, επιτρέπει σε κάθε χρήστη να διαγράφει μόνο τα δικά του αρχεία.

Τα προγράμματα SetUID που ανήκουν στον root είναι από τους πιο ελκυστικούς στόχους για τους επιτιθέμενους, επειδή οποιοδήποτε σφάλμα τους —υπερχείλιση ενδιάμεσης μνήμης, μη ασφαλής κλήση άλλου προγράμματος, διαδρομή βιβλιοθήκης με δικαίωμα εγγραφής— οδηγεί αμέσως σε πλήρη διαχειριστικό έλεγχο. Γι' αυτό οι διαχειριστές συστημάτων οφείλουν να τα καταγράφουν περιοδικά (για παράδειγμα με την εντολή \texttt{find / -perm -4000 -type f}), να αφαιρούν το bit όπου δεν είναι απολύτως αναγκαίο και να προτιμούν λεπτομερέστερους μηχανισμούς, όπως οι δυνατότητες (capabilities).

\section{Πέρα από τα βασικά δικαιώματα: ACL, capabilities και υποχρεωτικός έλεγχος πρόσβασης}\label{sec:2.3}
Το μοντέλο κάτοχος–ομάδα–λοιποί είναι απλό, αλλά δεν μπορεί να εκφράσει κανόνες όπως \emph{η Αλίκη μπορεί να διαβάσει αυτό το αρχείο και ο Βασίλης να το τροποποιήσει, ενώ κανείς άλλος δεν μπορεί να κάνει τίποτα από τα δύο}. Οι \textbf{λίστες ελέγχου πρόσβασης POSIX (ACL)} επιλύουν το πρόβλημα, προσθέτοντας εγγραφές για συγκεκριμένους χρήστες και ομάδες, οι οποίες διαχειρίζονται με την εντολή \texttt{setfacl} και εμφανίζονται με την \texttt{getfacl}. Μια εγγραφή \emph{μάσκας} περιορίζει τα μέγιστα δικαιώματα που μπορεί να λάβει οποιαδήποτε ονομαστική εγγραφή, επιτρέποντας στους διαχειριστές να περιορίζουν γρήγορα την πρόσβαση.

Οι \textbf{δυνατότητες (capabilities) του Linux} διασπούν το παντοδύναμο προνόμιο του root σε περίπου σαράντα ανεξάρτητες μονάδες. Ένας διακομιστής ιστού που χρειάζεται μόνο να δεσμεύσει τη θύρα 80 μπορεί να λάβει τη δυνατότητα \texttt{CAP\_NET\_BIND\_SERVICE} αντί να εκτελείται ως root, ώστε η παραβίασή του να μην παρέχει έλεγχο ολόκληρου του συστήματος. Τέλος, τα πλαίσια \textbf{υποχρεωτικού ελέγχου πρόσβασης} (Mandatory Access Control, MAC), όπως τα \textbf{SELinux} και \textbf{AppArmor}, επιβάλλουν μια πολιτική σε επίπεδο συστήματος την οποία δεν μπορεί να παρακάμψει ούτε ο κάτοχος του αρχείου. Υπό καθεστώς MAC, μια παραβιασμένη διεργασία παραμένει περιορισμένη στους πόρους που της επιτρέπει ρητά η πολιτική —εφαρμογή της αρχής των ασφαλών προεπιλογών που εισήχθη στο Κεφάλαιο 1.

\section{Αρχιτεκτονική ασφάλειας των Windows: ταυτότητες, token, UAC και μητρώο}\label{sec:2.4}
Τα Windows αναγνωρίζουν κάθε χρήστη, ομάδα και υπολογιστή μέσω ενός \textbf{αναγνωριστικού ασφάλειας} (Security Identifier, SID), μιας μοναδικής τιμής που παραμένει σταθερή ακόμη κι αν ο λογαριασμός μετονομαστεί. Όταν ένας χρήστης συνδέεται, η Τοπική Αρχή Ασφάλειας (LSA) δημιουργεί ένα \textbf{διακριτικό πρόσβασης} (access token), το οποίο περιέχει το SID του χρήστη, τα SID των ομάδων του και τα προνόμιά του. Κάθε διεργασία που εκκινεί ο χρήστης κληρονομεί αντίγραφο αυτού του διακριτικού. Αντικείμενα όπως αρχεία, κλειδιά μητρώου και υπηρεσίες φέρουν έναν \textbf{περιγραφέα ασφάλειας} (security descriptor), του οποίου η \textbf{λίστα διακριτικού ελέγχου πρόσβασης} (DACL) περιέχει εγγραφές που επιτρέπουν ή απαγορεύουν την πρόσβαση. Όταν μια διεργασία ζητά πρόσβαση, η Οθόνη Αναφοράς Ασφάλειας (Security Reference Monitor) συγκρίνει το διακριτικό με τη DACL· οι ρητές απαγορεύσεις αξιολογούνται πριν από τις άδειες.

Ο \textbf{Έλεγχος Λογαριασμού Χρήστη} (User Account Control, UAC) αντιμετωπίζει ένα ιστορικό πρόβλημα: για χρόνια, οι περισσότεροι χρήστες των Windows εργάζονταν μόνιμα ως διαχειριστές. Με το UAC, ένας διαχειριστής λαμβάνει κατά τη σύνδεση δύο διακριτικά —ένα \emph{φιλτραρισμένο} διακριτικό απλού χρήστη για την καθημερινή εργασία και ένα \emph{πλήρες} διακριτικό που ενεργοποιείται μόνο μετά από ρητή συγκατάθεση. Το \textbf{μητρώο} (registry), η ιεραρχική βάση δεδομένων που αποθηκεύει τις ρυθμίσεις του συστήματος και των εφαρμογών, προστατεύεται από τους ίδιους περιγραφείς ασφάλειας. Κλειδιά μητρώου όπως το \texttt{HKLM\textbackslash{}...\textbackslash{}Run} αποτελούν αγαπημένο σημείο εγκατάστασης μόνιμης παρουσίας (persistence) για το κακόβουλο λογισμικό· γι' αυτό τα δικαιώματα και οι μεταβολές τους πρέπει να παρακολουθούνται.

\section{Τα Windows σε βάθος: επίπεδα ακεραιότητας, προνόμια και επιφάνεια πλευρικής κίνησης}\label{sec:2.5}
Από τα Windows Vista και έπειτα, κάθε διεργασία και κάθε αντικείμενο φέρει επιπλέον ένα \textbf{υποχρεωτικό επίπεδο ακεραιότητας}: \emph{Χαμηλό}, \emph{Μεσαίο}, \emph{Υψηλό} ή \emph{Συστήματος}. Μια διεργασία δεν μπορεί να γράψει σε αντικείμενο υψηλότερου επιπέδου ακεραιότητας, ανεξάρτητα από το τι ορίζει η DACL. Πρόκειται για άμεση εφαρμογή του μοντέλου Biba (Κεφάλαιο 1). Οι φυλλομετρητές ιστού αξιοποιούν αυτόν τον μηχανισμό, εκτελώντας τις διεργασίες απόδοσης περιεχομένου σε Χαμηλό επίπεδο ακεραιότητας, ώστε κακόβουλο περιεχόμενο ιστού να μην μπορεί να τροποποιήσει τα αρχεία του χρήστη ακόμη κι αν επιτύχει εκτέλεση κώδικα.

Ορισμένα \textbf{προνόμια} ενός διακριτικού είναι τόσο ισχυρά, ώστε ισοδυναμούν ουσιαστικά με πλήρη έλεγχο. Το \texttt{SeDebugPrivilege} επιτρέπει την ανάγνωση της μνήμης οποιασδήποτε διεργασίας —συμπεριλαμβανομένης της LSASS, όπου διατηρούνται διαπιστευτήρια— ενώ τα \texttt{SeImpersonatePrivilege} και \texttt{SeBackupPrivilege} έχουν επανειλημμένα χρησιμοποιηθεί για κλιμάκωση προνομίων. Σε ένα περιβάλλον τομέα (domain), αυτές οι τοπικές αδυναμίες μετατρέπονται σε δικτυακό πρόβλημα: τα αποθηκευμένα διαπιστευτήρια, οι κοινοί κωδικοί τοπικού διαχειριστή και τα πρωτόκολλα απομακρυσμένης διαχείρισης (SMB, WMI, WinRM, RDP) επιτρέπουν σε έναν επιτιθέμενο να κινηθεί πλευρικά από έναν παραβιασμένο σταθμό εργασίας σε πολλούς άλλους. Τα αντίμετρα περιλαμβάνουν μοναδικούς κωδικούς τοπικού διαχειριστή (Windows LAPS), το Credential Guard, την κλιμακωτή (tiered) διαχείριση και τον περιορισμό των λογαριασμών που επιτρέπεται να συνδέονται σε κάθε σύστημα.

\begin{table}[htbp]
  \centering\small
  \caption{Σύγκριση ελέγχου πρόσβασης σε POSIX και Windows}
  \begin{tabularx}{\linewidth}{>{\raggedright\arraybackslash}X >{\raggedright\arraybackslash}X >{\raggedright\arraybackslash}X}
    \toprule
    \textbf{Χαρακτηριστικό} & \textbf{POSIX (Linux)} & \textbf{Windows} \\
    \midrule
    Ταυτότητα & UID / GID & SID \\
    Πολιτική ανά αντικείμενο & Βασικά δικαιώματα + προαιρετική ACL & Περιγραφέας ασφάλειας με DACL \\
    Ανύψωση προνομίων & SetUID, sudo, capabilities & Συγκατάθεση UAC, προνόμια στο διακριτικό \\
    Υποχρεωτικός έλεγχος & SELinux, AppArmor & Επίπεδα ακεραιότητας, AppContainer \\
    Συνήθης οδός κλιμάκωσης & Ευάλωτο εκτελέσιμο SetUID & Κατάχρηση προνομίου διακριτικού, αδύναμη ACL υπηρεσίας \\
    \bottomrule
  \end{tabularx}
\end{table}
\section*{Βασικοί όροι}
\begin{description}[style=nextline,leftmargin=1.2em]
  \item[DAC] Διακριτικός έλεγχος πρόσβασης: ο κάτοχος ενός πόρου αποφασίζει ποιος έχει πρόσβαση σε αυτόν.
  \item[SetUID] Ειδικό δικαίωμα που κάνει ένα πρόγραμμα να εκτελείται με τα προνόμια του κατόχου του.
  \item[Δυνατότητα (capability)] Αυτοτελής μονάδα του προνομίου root που μπορεί να παραχωρηθεί χωριστά.
  \item[MAC] Υποχρεωτικός έλεγχος πρόσβασης: πολιτική σε επίπεδο συστήματος που δεν μπορεί να παρακάμψει ο κάτοχος (SELinux, AppArmor).
  \item[Διακριτικό πρόσβασης] Δομή των Windows που περιέχει την ταυτότητα, τις ομάδες και τα προνόμια μιας διεργασίας.
  \item[Επίπεδο ακεραιότητας] Ετικέτα των Windows (Χαμηλό–Συστήματος) που εμποδίζει διεργασίες χαμηλότερου επιπέδου να γράφουν σε αντικείμενα υψηλότερου.
\end{description}

\section*{Σύνοψη κεφαλαίου}
\begin{itemize}
  \item Τα συστήματα POSIX ελέγχουν την πρόσβαση μέσω δικαιωμάτων για κάτοχο, ομάδα και λοιπούς· αποφασίζει η πρώτη κατηγορία που ταιριάζει.
  \item Τα SetUID και SetGID είναι αναγκαία σε λίγα σημεία, αλλά αποτελούν βασικό στόχο κλιμάκωσης προνομίων και πρέπει να καταγράφονται.
  \item Οι ACL προσθέτουν λεπτομερείς διακριτικούς κανόνες, οι δυνατότητες διασπούν το προνόμιο του root και ο MAC περιορίζει ακόμη και παραβιασμένες διεργασίες.
  \item Τα Windows συνδυάζουν SID, διακριτικά πρόσβασης και DACL· το UAC διαχωρίζει την καθημερινή εργασία από τις διαχειριστικές ενέργειες.
  \item Τα επίπεδα ακεραιότητας και η προσεκτική διαχείριση προνομίων περιορίζουν τόσο την τοπική κλιμάκωση όσο και την πλευρική κίνηση σε έναν τομέα.
\end{itemize}

\section*{Ερωτήσεις ανασκόπησης}
\begin{enumerate}
  \item Ένα αρχείο έχει δικαιώματα \texttt{604} και κάτοχο την alice. Μπορεί η alice να το διαβάσει; Μπορεί ο bob; Εξηγήστε τη σειρά αξιολόγησης που οδηγεί στην απάντησή σας.
  \item Γιατί ένα ευάλωτο εκτελέσιμο SetUID που ανήκει στον root είναι πιο επικίνδυνο από την ίδια ευπάθεια σε ένα συνηθισμένο πρόγραμμα;
  \item Συγκρίνετε τις δυνατότητες του Linux με τα προνόμια διακριτικού των Windows. Ποια αρχή ασφάλειας υλοποιούν και οι δύο μηχανισμοί;
  \item Εξηγήστε πώς οι κοινοί κωδικοί τοπικού διαχειριστή διευκολύνουν την πλευρική κίνηση και περιγράψτε δύο αντίμετρα.
\end{enumerate}

\part{Αντίπαλοι και Επιθέσεις}
% ============================================================
% Chapter 3: Threat Actors, Motivations and Attack Lifecycle Models
% Language: Greek — edit freely; regenerate from src/content if preferred.
% ============================================================
\chapter{Δράστες Απειλών, Κίνητρα και Μοντέλα Κύκλου Ζωής Επιθέσεων}\label{ch:3}
\chaptermeta{Ταξινόμηση δραστών · Το φάσμα των χάκερ · Cyber Kill Chain · MITRE ATT\&CK · Μοντέλο Διαμαντιού και D3FEND · Πληροφορίες απειλών και ανταλλαγή τους}{Επίπεδο: Εισαγωγικό–Μεσαίο επίπεδο · Σκέψη του αντιπάλου}{Χρόνος μελέτης: 6–8 ώρες}

Οι άμυνες αποκτούν νόημα μόνο σε σχέση με τους αντιπάλους που καλούνται να σταματήσουν. Μια μικρή επιχείρηση και ένας εθνικός πάροχος ενέργειας αντιμετωπίζουν πολύ διαφορετικούς αντιπάλους, με διαφορετικούς στόχους, πόρους και υπομονή. Για τον λόγο αυτό, το κεφάλαιο μετατοπίζει την οπτική από το σύστημα στον επιτιθέμενο.

Αρχικά ταξινομούμε τους δράστες απειλών με βάση το κίνητρο και τις ικανότητές τους και αποσαφηνίζουμε τις ηθικές και νομικές διακρίσεις πίσω από τη γνωστή ορολογία των «καπέλων». Στη συνέχεια μελετάμε τρία συμπληρωματικά μοντέλα που περιγράφουν την εξέλιξη μιας επίθεσης: το Cyber Kill Chain της Lockheed Martin, τη βάση γνώσης MITRE ATT\&CK και το Μοντέλο Διαμαντιού για την ανάλυση εισβολών. Το κεφάλαιο ολοκληρώνεται με τις πληροφορίες κυβερνοαπειλών (Cyber Threat Intelligence) —πώς παράγεται, αξιολογείται και ανταλλάσσεται υπεύθυνα η γνώση για τους αντιπάλους.

\begin{outcomesbox}{Μαθησιακά αποτελέσματα}
Μετά τη μελέτη του κεφαλαίου θα πρέπει να μπορείτε:
\begin{itemize}
  \item Να ταξινομείτε τους δράστες απειλών με βάση το κίνητρο, τις ικανότητες και την επιμονή τους και να τους συσχετίζετε με κατάλληλες άμυνες.
  \item Να διακρίνετε τη δραστηριότητα white-, grey- και black-hat με κριτήριο την εξουσιοδότηση και την πρόθεση.
  \item Να αντιστοιχίζετε μια επίθεση στις φάσεις του Cyber Kill Chain και στις τακτικές και τεχνικές του ATT\&CK.
  \item Να χρησιμοποιείτε το Μοντέλο Διαμαντιού για τη δομημένη ανάλυση μιας εισβολής.
  \item Να περιγράφετε τον κύκλο ζωής των πληροφοριών απειλών και να εφαρμόζετε το Πρωτόκολλο Φωτεινού Σηματοδότη (TLP) κατά την ανταλλαγή τους.
\end{itemize}
\end{outcomesbox}

\section{Ταξινόμηση των δραστών απειλών}\label{sec:3.1}
Οι δράστες απειλών ομαδοποιούνται χρήσιμα με βάση το τι επιδιώκουν και τι μπορούν να κάνουν. Οι \textbf{κρατικά υποστηριζόμενες ομάδες}, γνωστές και ως \emph{προηγμένες επίμονες απειλές} (Advanced Persistent Threats, APT), επιδιώκουν κατασκοπεία, δολιοφθορά ή επηρεασμό· διαθέτουν τους πόρους για εκμεταλλεύσεις μηδενικής ημέρας (zero-day) και για επιχειρήσεις μέσω της εφοδιαστικής αλυσίδας και μπορούν να παραμένουν σε ένα δίκτυο για μήνες ή χρόνια. Οι \textbf{κυβερνοεγκληματίες} έχουν οικονομικά κίνητρα· η οικονομία του λυτρισμικού ως υπηρεσίας (Ransomware-as-a-Service) έχει δώσει ακόμη και σε συνεργάτες μέτριας δεξιότητας πρόσβαση σε επαγγελματικά εργαλεία. Οι \textbf{χακτιβιστές} επιδιώκουν δημοσιότητα για έναν πολιτικό ή κοινωνικό σκοπό, συνήθως μέσω επιθέσεων άρνησης υπηρεσίας, αλλοίωσης ιστοτόπων ή διαρροών.

Οι \textbf{εσωτερικοί δράστες} χρειάζονται ιδιαίτερη προσοχή, διότι διαθέτουν ήδη νόμιμη πρόσβαση και, επομένως, ακυρώνουν τις παραδοχές που βασίζονται στην περίμετρο. Ένας \emph{κακόβουλος} εσωτερικός δράστης μπορεί να υποκλέψει δεδομένα ή να δολιοφθείρει συστήματα από απληστία ή μνησικακία, ενώ ένας \emph{αμελής} προκαλεί ζημιά ακούσια, για παράδειγμα μέσω λανθασμένης ρύθμισης ή επειδή ενέδωσε σε ένα μήνυμα ηλεκτρονικού ψαρέματος. Η κατανόηση των δραστών που αφορούν πραγματικά έναν οργανισμό τού επιτρέπει να επιλέγει μέτρα ανάλογα με τον πραγματικό κίνδυνο και όχι με τους πιο εντυπωσιακούς τίτλους ειδήσεων.

\begin{table}[htbp]
  \centering\small
  \caption{Προφίλ δραστών απειλών}
  \begin{tabularx}{\linewidth}{>{\raggedright\arraybackslash}X >{\raggedright\arraybackslash}X >{\raggedright\arraybackslash}X >{\raggedright\arraybackslash}X >{\raggedright\arraybackslash}X}
    \toprule
    \textbf{Δράστης} & \textbf{Κίνητρο} & \textbf{Ικανότητα} & \textbf{Τυπική διάρκεια παραμονής} & \textbf{Χαρακτηριστικές τεχνικές} \\
    \midrule
    Κρατική ομάδα APT & Κατασκοπεία, δολιοφθορά, επηρεασμός & Πολύ υψηλή & Μήνες έως χρόνια & Στοχευμένο ψάρεμα, χρήση νόμιμων εργαλείων του συστήματος, κρυφή διοίκηση και έλεγχος \\
    Κυβερνοεγκληματίας & Οικονομικό όφελος & Μέτρια έως υψηλή & Ημέρες έως εβδομάδες & Ψάρεμα, λυτρισμικό, κλοπή δεδομένων για εκβιασμό \\
    Χακτιβιστής & Ιδεολογία, δημοσιότητα & Χαμηλή έως μέτρια & Ώρες έως ημέρες & DDoS, αλλοίωση ιστοτόπων, δημοσιοποίηση προσωπικών στοιχείων \\
    Κακόβουλος εσωτερικός & Όφελος, μνησικακία & Μέτρια (εξουσιοδοτημένη πρόσβαση) & Μεταβλητή & Κατάχρηση προνομίων, διαρροή δεδομένων \\
    Αμελής εσωτερικός & Κανένα (σφάλμα) & Χαμηλή (ακούσια) & — & Λανθασμένες ρυθμίσεις, επαναχρησιμοποίηση κωδικών \\
    \bottomrule
  \end{tabularx}
\end{table}
\section{Το φάσμα των χάκερ: ηθική και νομιμότητα}\label{sec:3.2}
Η δημοφιλής ορολογία των «καπέλων» ταξινομεί την \textbf{εξουσιοδότηση και την πρόθεση}, όχι την τεχνική δεξιότητα. Ένας \textbf{white-hat} (ηθικός) χάκερ ελέγχει συστήματα μόνο με ρητή, γραπτή άδεια και εντός συμφωνημένου πεδίου, αναφέροντας τα ευρήματα ώστε να διορθωθούν. Ένας \textbf{black-hat} χάκερ ενεργεί χωρίς εξουσιοδότηση και με κακόβουλη ή ιδιοτελή πρόθεση. Ανάμεσά τους βρίσκεται ο \textbf{grey-hat}, ο οποίος μπορεί να ενεργεί χωρίς άδεια αλλά χωρίς σαφή κακοβουλία —για παράδειγμα, σαρώνει τον ιστότοπο μιας εταιρείας και στη συνέχεια την ενημερώνει για μια ευπάθεια.

Οι φοιτητές πρέπει να κατανοήσουν ότι οι καλές προθέσεις δεν δημιουργούν νομική εξουσιοδότηση. Στις περισσότερες έννομες τάξεις, συμπεριλαμβανομένων εκείνων που ενσωματώνουν την Οδηγία της ΕΕ για τις επιθέσεις κατά συστημάτων πληροφοριών, η πρόσβαση σε σύστημα χωρίς άδεια αποτελεί ποινικό αδίκημα ακόμη κι αν δεν προκληθεί ζημιά. Γι' αυτό η επαγγελματική δοκιμή διείσδυσης (Κεφάλαιο 10) ξεκινά πάντοτε με υπογεγραμμένο πεδίο εφαρμογής και κανόνες εμπλοκής, και γι' αυτό υπάρχουν τα προγράμματα συντονισμένης γνωστοποίησης ευπαθειών και τα προγράμματα επιβράβευσης (bug bounties): μετατρέπουν την περιέργεια του grey-hat σε εξουσιοδοτημένη και νομικά προστατευμένη έρευνα.

\section{Το Cyber Kill Chain της Lockheed Martin}\label{sec:3.3}
Το \textbf{Cyber Kill Chain}, που δημοσίευσε η Lockheed Martin το 2011, περιγράφει μια στοχευμένη εισβολή ως επτά διαδοχικές φάσεις: \textbf{αναγνώριση} (συλλογή πληροφοριών για τον στόχο), \textbf{οπλοποίηση} (σύζευξη μιας εκμετάλλευσης με ένα ωφέλιμο φορτίο), \textbf{παράδοση} (για παράδειγμα, ένα συνημμένο ηλεκτρονικού ταχυδρομείου), \textbf{εκμετάλλευση} (ενεργοποίηση της ευπάθειας), \textbf{εγκατάσταση} (εδραίωση μόνιμης παρουσίας), \textbf{διοίκηση και έλεγχος} (δημιουργία απομακρυσμένου καναλιού) και \textbf{ενέργειες επί των στόχων} (κλοπή δεδομένων, κρυπτογράφηση, δολιοφθορά).

Η βασική ιδέα του μοντέλου είναι αμυντική: επειδή ο επιτιθέμενος πρέπει να επιτύχει σε \emph{κάθε} φάση, ο αμυνόμενος αρκεί να σπάσει την αλυσίδα \emph{μία} φορά. Η έγκαιρη διακοπή είναι φθηνότερη —ο αποκλεισμός ενός κακόβουλου συνημμένου στη φάση της παράδοσης αποτρέπει μια δαπανηρή απόκριση σε περιστατικό αργότερα. Το μοντέλο έχει όμως περιορισμούς. Σχεδιάστηκε με βάση εισβολές με κακόβουλο λογισμικό που προέρχονται από το εξωτερικό και περιγράφει λιγότερο ικανοποιητικά τις εσωτερικές απειλές, την κατάχρηση διαπιστευτηρίων και τις επιθέσεις στο νέφος. Για τον λόγο αυτό συνδυάζεται συνήθως με το λεπτομερέστερο πλαίσιο ATT\&CK.

\section{MITRE ATT\&CK: τακτικές, τεχνικές και διαδικασίες}\label{sec:3.4}
Το \textbf{MITRE ATT\&CK} είναι μια δημόσια διαθέσιμη βάση γνώσης για τη συμπεριφορά των αντιπάλων, η οποία έχει δημιουργηθεί από παρατηρήσεις πραγματικών εισβολών. Οργανώνεται ως πίνακας. Οι στήλες είναι οι \textbf{τακτικές} —οι βραχυπρόθεσμοι στόχοι του επιτιθέμενου, όπως η \emph{Αρχική Πρόσβαση}, η \emph{Εκτέλεση}, η \emph{Μόνιμη Παρουσία}, η \emph{Κλιμάκωση Προνομίων}, η \emph{Αποφυγή Άμυνας}, η \emph{Πρόσβαση σε Διαπιστευτήρια}, η \emph{Πλευρική Κίνηση} και η \emph{Διαρροή Δεδομένων}. Κάθε στήλη περιέχει \textbf{τεχνικές} που περιγράφουν \emph{πώς} επιτυγχάνεται ο στόχος· για παράδειγμα, η T1059.001 δηλώνει εκτέλεση εντολών μέσω PowerShell. Οι \textbf{διαδικασίες} είναι οι συγκεκριμένοι τρόποι με τους οποίους μια ομάδα έχει χρησιμοποιήσει μια τεχνική.

Το ATT\&CK παρέχει στους αμυνόμενους μια κοινή γλώσσα. Μια ομάδα ασφάλειας μπορεί να αντιστοιχίσει τους κανόνες ανίχνευσής της σε τεχνικές και να διαπιστώσει αμέσως πού υπάρχουν κενά κάλυψης· μια αναφορά πληροφοριών απειλών μπορεί να δηλώνει ποιες τεχνικές χρησιμοποιεί μια ομάδα· και μια κόκκινη ομάδα (red team) μπορεί να προσομοιώσει αυτές τις τεχνικές για να ελέγξει αν οι άμυνες λειτουργούν πράγματι. Η έμφαση στη \emph{συμπεριφορά} και όχι σε εύκολα μεταβαλλόμενες ενδείξεις, όπως οι τιμές κατακερματισμού αρχείων, είναι αυτό που καθιστά το πλαίσιο ανθεκτικό στον χρόνο.

\section{Το Μοντέλο Διαμαντιού και το MITRE D3FEND}\label{sec:3.5}
Το \textbf{Μοντέλο Διαμαντιού για την Ανάλυση Εισβολών} (Diamond Model) αναπαριστά κάθε κακόβουλο γεγονός μέσω τεσσάρων συνδεδεμένων στοιχείων: του \textbf{αντιπάλου}, της \textbf{δυνατότητας} (εργαλεία και κακόβουλο λογισμικό), της \textbf{υποδομής} (ονόματα τομέα, διακομιστές, διευθύνσεις IP) και του \textbf{θύματος}. Η δύναμη του μοντέλου βρίσκεται στη \emph{μετάβαση} (pivoting) από στοιχείο σε στοιχείο: μόλις ο αναλυτής γνωρίζει μία κορυφή, οι συνδέσεις υποδεικνύουν πού να αναζητήσει στη συνέχεια. Ένα δείγμα κακόβουλου λογισμικού (δυνατότητα) μπορεί να αποκαλύψει έναν τομέα διοίκησης και ελέγχου (υποδομή), του οποίου τα στοιχεία καταχώρισης μπορεί να οδηγούν σε άλλα θύματα και, τελικά, σε ένα προφίλ αντιπάλου.

Ενώ το ATT\&CK καταγράφει την επιθετική συμπεριφορά, το \textbf{MITRE D3FEND} καταγράφει αμυντικά αντίμετρα —όπως η \emph{ανάλυση δημιουργίας διεργασιών} ή η \emph{θωράκιση διαπιστευτηρίων}— και τα συνδέει με τις επιθετικές τεχνικές που αντιμετωπίζουν. Σε συνδυασμό, το Kill Chain απαντά στο \emph{πότε} να παρέμβουμε, το ATT\&CK στο \emph{τι} κάνει ο αντίπαλος, το Μοντέλο Διαμαντιού στο \emph{ποιος και με ποια μέσα} και το D3FEND στο \emph{ποιο μέτρο} αντιμετωπίζει κάθε συμπεριφορά.

\section{Πληροφορίες κυβερνοαπειλών: κύκλος ζωής, είδη και υπεύθυνη ανταλλαγή}\label{sec:3.6}
Οι \textbf{πληροφορίες κυβερνοαπειλών} (Cyber Threat Intelligence, CTI) είναι πληροφορίες για τους αντιπάλους που έχουν συλλεχθεί, αναλυθεί και τοποθετηθεί σε πλαίσιο, ώστε να υποστηρίζουν τη λήψη αποφάσεων. Παράγονται μέσω ενός \textbf{κύκλου ζωής πληροφοριών}: \emph{καθοδήγηση} (προσδιορισμός του τι χρειάζεται να γνωρίζουν όσοι λαμβάνουν αποφάσεις), \emph{συλλογή}, \emph{επεξεργασία}, \emph{ανάλυση}, \emph{διάχυση} και \emph{ανατροφοδότηση}. Οι πληροφορίες διακρίνονται συχνά σε τρία επίπεδα. Οι \textbf{στρατηγικές} ενημερώνουν τη διοίκηση για τάσεις και κινδύνους· οι \textbf{επιχειρησιακές} περιγράφουν εκστρατείες και δράστες· οι \textbf{τακτικές} παρέχουν τεχνικές λεπτομέρειες, όπως τεχνικές επίθεσης και ενδείξεις παραβίασης.

Οι πληροφορίες αποκτούν αξία όταν ανταλλάσσονται, όμως η ανταλλαγή πρέπει να σέβεται τις επιθυμίες της πηγής. Το \textbf{Πρωτόκολλο Φωτεινού Σηματοδότη} (Traffic Light Protocol, TLP 2.0) χαρακτηρίζει τις πληροφορίες ως \emph{TLP:RED} (μόνο ονομαστικοί παραλήπτες), \emph{TLP:AMBER} (ο οργανισμός του παραλήπτη, βάσει ανάγκης γνώσης), \emph{TLP:GREEN} (η ευρύτερη κοινότητα) ή \emph{TLP:CLEAR} (δημόσιες). Η μηχανικά αναγνώσιμη ανταλλαγή βασίζεται στο \textbf{STIX}, μια τυποποιημένη γλώσσα περιγραφής πληροφοριών απειλών, και στο \textbf{TAXII}, ένα πρωτόκολλο μεταφοράς τους. Κλαδικές κοινότητες, όπως τα Κέντρα Ανταλλαγής και Ανάλυσης Πληροφοριών (ISAC), επιτρέπουν σε οργανισμούς που αντιμετωπίζουν παρόμοιες απειλές να προειδοποιούν γρήγορα ο ένας τον άλλον.

\section*{Βασικοί όροι}
\begin{description}[style=nextline,leftmargin=1.2em]
  \item[APT] Προηγμένη επίμονη απειλή: δράστης με σημαντικούς πόρους που διατηρεί μακροχρόνια κρυφή πρόσβαση.
  \item[Kill Chain] Μοντέλο επτά φάσεων μιας στοχευμένης εισβολής, από την αναγνώριση έως τις ενέργειες επί των στόχων.
  \item[TTP] Τακτικές, τεχνικές και διαδικασίες: η συμπεριφορική περιγραφή του τρόπου δράσης ενός αντιπάλου.
  \item[Μοντέλο Διαμαντιού] Αναλυτικό μοντέλο που συνδέει αντίπαλο, δυνατότητα, υποδομή και θύμα.
  \item[TLP] Πρωτόκολλο Φωτεινού Σηματοδότη: χαρακτηρισμοί που ορίζουν πόσο ευρέως μπορεί να αναδιανεμηθεί μια πληροφορία.
  \item[STIX / TAXII] Τυποποιημένη γλώσσα (STIX) και πρωτόκολλο μεταφοράς (TAXII) για την ανταλλαγή πληροφοριών απειλών.
\end{description}

\section*{Σύνοψη κεφαλαίου}
\begin{itemize}
  \item Οι δράστες απειλών διαφέρουν ως προς το κίνητρο, τις ικανότητες και την επιμονή· οι άμυνες πρέπει να είναι ανάλογες με τους δράστες που πράγματι αφορούν τον οργανισμό.
  \item Η ορολογία των «καπέλων» αφορά την εξουσιοδότηση και την πρόθεση· η μη εξουσιοδοτημένη πρόσβαση είναι παράνομη ανεξάρτητα από το κίνητρο.
  \item Το Kill Chain δείχνει ότι η διακοπή οποιασδήποτε φάσης σταματά την εισβολή και ότι η έγκαιρη διακοπή κοστίζει λιγότερο.
  \item Το ATT\&CK περιγράφει λεπτομερώς τη συμπεριφορά των αντιπάλων και παρέχει κοινή γλώσσα για την ανίχνευση, τις πληροφορίες απειλών και τις δοκιμές.
  \item Το Μοντέλο Διαμαντιού υποστηρίζει τη μετάβαση μεταξύ στοιχείων κατά την ανάλυση, ενώ οι πληροφορίες πρέπει να ανταλλάσσονται με σαφείς κανόνες, όπως το TLP.
\end{itemize}

\section*{Ερωτήσεις ανασκόπησης}
\begin{enumerate}
  \item Γιατί οι εσωτερικές απειλές αντιμετωπίζονται δύσκολα με μέτρα που βασίζονται στην περίμετρο; Προτείνετε δύο μέτρα αποτελεσματικά απέναντί τους.
  \item Αντιστοιχίστε μια τυπική επίθεση που ξεκινά με ηλεκτρονικό ψάρεμα και καταλήγει σε λυτρισμικό στις επτά φάσεις του Cyber Kill Chain.
  \item Εξηγήστε γιατί οι ανιχνεύσεις που βασίζονται σε τεχνικές του ATT\&CK είναι πιο ανθεκτικές από εκείνες που βασίζονται σε τιμές κατακερματισμού αρχείων.
  \item Ένας συνεργάτης σάς αποστέλλει αναφορά με χαρακτηρισμό TLP:AMBER. Μπορείτε να την προωθήσετε σε έναν προμηθευτή; Τεκμηριώστε την απάντησή σας.
\end{enumerate}

% ============================================================
% Chapter 4: Taxonomy of Malware and Malicious Code
% Language: Greek — edit freely; regenerate from src/content if preferred.
% ============================================================
\chapter{Ταξινόμηση του Κακόβουλου Λογισμικού και του Κακόβουλου Κώδικα}\label{ch:4}
\chaptermeta{Ιοί · Σκουλήκια (worms) · Δούρειοι ίπποι (trojans) · Λογισμικό κατασκοπείας · Λυτρισμικό · Επιθέσεις χωρίς αρχεία · Rootkits · Ενδείξεις παραβίασης · Stuxnet, WannaCry, Emotet}{Επίπεδο: Μεσαίο επίπεδο · Κακόβουλος κώδικας}{Χρόνος μελέτης: 8–10 ώρες}

Το κακόβουλο λογισμικό είναι το πιο ορατό εργαλείο των κυβερνοεπιθέσεων, ωστόσο τα μέσα ενημέρωσης χρησιμοποιούν συχνά την ορολογία του χαλαρά. Ο χαρακτηρισμός κάθε μόλυνσης ως «ιού» δεν είναι μόνο ανακριβής· μπορεί να οδηγήσει σε λανθασμένη απόκριση. Ένα σκουλήκι περιορίζεται με την απομόνωση δικτύων, ενώ ένας δούρειος ίππος σταματά με την παρεμπόδιση της εκτέλεσής του. Η ακριβής ταξινόμηση αποτελεί επομένως επιχειρησιακή αναγκαιότητα και όχι ακαδημαϊκή άσκηση.

Το κεφάλαιο οργανώνει το κακόβουλο λογισμικό με βάση τον \emph{μηχανισμό} του —πώς διαδίδεται, πώς κρύβεται και τι κάνει— και όχι με βάση τα ονόματα μεμονωμένων οικογενειών. Στη συνέχεια εξετάζει τις προηγμένες τεχνικές που επιτρέπουν στα σύγχρονα εμφυτεύματα να αποφεύγουν την ανίχνευση, εξηγεί πώς αναγνωρίζεται μια παραβίαση σε υπολογιστές και δίκτυα, αναλύει τη σύγχρονη επιχείρηση λυτρισμικού και αντλεί διδάγματα από τρεις ιστορικές περιπτώσεις.

\begin{outcomesbox}{Μαθησιακά αποτελέσματα}
Μετά τη μελέτη του κεφαλαίου θα πρέπει να μπορείτε:
\begin{itemize}
  \item Να διακρίνετε τις κατηγορίες κακόβουλου λογισμικού με βάση την αναπαραγωγή, τη διάδοση και το ωφέλιμο φορτίο τους και να επιλέγετε την κατάλληλη στρατηγική περιορισμού.
  \item Να εξηγείτε την εκτέλεση χωρίς αρχεία, τον πολυμορφισμό και τα rootkits, καθώς και τον λόγο για τον οποίο παρακάμπτουν την ανίχνευση με υπογραφές.
  \item Να αναγνωρίζετε ενδείξεις παραβίασης σε υπολογιστές και δίκτυα και να εξηγείτε γιατί απαιτείται συσχέτισή τους.
  \item Να περιγράφετε τα στάδια μιας σύγχρονης επίθεσης λυτρισμικού και τη σωστή σειρά απόκρισης.
  \item Να αντλείτε αμυντικά διδάγματα από τις περιπτώσεις Stuxnet, WannaCry και Emotet.
\end{itemize}
\end{outcomesbox}

\section{Βασικές κατηγορίες: τι αναπαράγεται, τι εξαπατά, τι αποφέρει κέρδος}\label{sec:4.1}
Ένας \textbf{ιός} προσκολλάται σε ένα αρχείο-ξενιστή ή έγγραφο και αναπαράγεται μόνο όταν αυτό εκτελεστεί· επομένως, εξαρτάται συνήθως από κάποια ενέργεια του χρήστη. Αντίθετα, ένα \textbf{σκουλήκι} (worm) είναι αυτοτελές και διαδίδεται από μόνο του στα δίκτυα, εκμεταλλευόμενο ευάλωτες υπηρεσίες· το Morris Worm και η διάδοση του WannaCry μέσω του πρωτοκόλλου SMB είναι κλασικά παραδείγματα. Ένας \textbf{δούρειος ίππος} (trojan) δεν αναπαράγεται καθόλου. Παρουσιάζεται ως νόμιμο λογισμικό, ώστε να πείσει το θύμα να τον εκτελέσει, και η επικινδυνότητά του έγκειται στην κοινωνική του αληθοφάνεια και στα ωφέλιμα φορτία που μεταφέρει.

Άλλες κατηγορίες ορίζονται από το τι κάνουν και όχι από το πώς διαδίδονται. Το \textbf{λογισμικό κατασκοπείας} (spyware) συλλέγει κρυφά τη δραστηριότητα ή τα δεδομένα του χρήστη· το \textbf{διαφημιστικό λογισμικό} (adware) εκμεταλλεύεται οικονομικά την προσοχή του θύματος και συχνά λειτουργεί ως πύλη για λογισμικό κατασκοπείας· οι \textbf{καταγραφείς πληκτρολογήσεων} (keyloggers) καταγράφουν ό,τι πληκτρολογείται για να υποκλέψουν κωδικούς. Το \textbf{λυτρισμικό} (ransomware) στερεί την πρόσβαση στα δεδομένα, είτε κρυπτογραφώντας αρχεία (\emph{κρυπτογραφικό λυτρισμικό}) είτε κλειδώνοντας το σύστημα (\emph{λυτρισμικό κλειδώματος}), και σήμερα συνήθως συνδυάζει την κρυπτογράφηση με κλοπή δεδομένων —τον λεγόμενο \emph{διπλό εκβιασμό}.

\begin{table}[htbp]
  \centering\small
  \caption{Κατηγορίες κακόβουλου λογισμικού και κύριες άμυνες}
  \begin{tabularx}{\linewidth}{>{\raggedright\arraybackslash}X >{\raggedright\arraybackslash}X >{\raggedright\arraybackslash}X >{\raggedright\arraybackslash}X}
    \toprule
    \textbf{Κατηγορία} & \textbf{Αυτοαναπαράγεται;} & \textbf{Απαιτεί ενέργεια χρήστη;} & \textbf{Κύρια άμυνα} \\
    \midrule
    Ιός & Ναι, μέσω αρχείων-ξενιστών & Συνήθως & Παρεμπόδιση εκτέλεσης, antivirus, ενημερώσεις \\
    Σκουλήκι & Ναι, μέσω δικτύου & Όχι & Κατάτμηση δικτύου, άμεση ενημέρωση υπηρεσιών, IDS \\
    Δούρειος ίππος & Όχι & Ναι (εξαπάτηση) & Λίστες επιτρεπόμενων εφαρμογών, υπογραφή κώδικα, ευαισθητοποίηση \\
    Λυτρισμικό & Ενίοτε (όπως σκουλήκι) & Συχνά (ψάρεμα, εκτεθειμένο RDP) & Αμετάβλητα αντίγραφα ασφαλείας, MFA, EDR, ενημερώσεις \\
    Λογισμικό κατασκοπείας / keylogger & Όχι & Ναι (ενσωματωμένο σε λογισμικό) & Ελάχιστο προνόμιο, παρακολούθηση εξερχόμενης κίνησης \\
    \bottomrule
  \end{tabularx}
\end{table}
\section{Προηγμένοι μηχανισμοί: εκτέλεση χωρίς αρχεία, πολυμορφισμός και rootkits}\label{sec:4.2}
Τα παραδοσιακά προϊόντα προστασίας από ιούς σαρώνουν αρχεία στον δίσκο και τα συγκρίνουν με γνωστές υπογραφές. Τα σύγχρονα εμφυτεύματα σχεδιάζονται ακριβώς για να παρακάμπτουν αυτή την προσέγγιση. Το \textbf{κακόβουλο λογισμικό χωρίς αρχεία} (fileless malware) βρίσκεται στη μνήμη, στο μητρώο, σε συνδρομές WMI ή σε προγραμματισμένες εργασίες και κάνει κατάχρηση νόμιμων εργαλείων του συστήματος, όπως το PowerShell, το \texttt{mshta} ή το \texttt{wmic}. Η πρακτική αυτή ονομάζεται \emph{living off the land}, επειδή ο επιτιθέμενος χρησιμοποιεί ό,τι είναι ήδη εγκατεστημένο. Τεχνικές όπως η \emph{ανακλαστική έγχυση DLL} (reflective DLL injection) και η \emph{εκκένωση διεργασίας} (process hollowing) φορτώνουν κακόβουλο κώδικα στη μνήμη μιας αξιόπιστης διεργασίας, ώστε κανένα ύποπτο εκτελέσιμο να μη γράφεται ποτέ στον δίσκο.

Ο \textbf{πολυμορφικός} και ο \textbf{μεταμορφικός} κώδικας μεταβάλλει τη δυαδική του αναπαράσταση σε κάθε μόλυνση —μέσω κρυπτογραφικών περιτυλιγμάτων ή αντικατάστασης εντολών— ενώ η συμπεριφορά του παραμένει ίδια. Μια υπογραφή που γράφτηκε για ένα αντίγραφο αποτυγχάνει επομένως να ταιριάξει με το επόμενο· γι' αυτό οι αμυνόμενοι στρέφονται στην ανίχνευση συμπεριφοράς και σε κανόνες προτύπων όπως οι κανόνες YARA (Κεφάλαιο 11). Τα \textbf{rootkits} προχωρούν ένα βήμα παραπέρα: λειτουργώντας σε επίπεδο χρήστη, στον πυρήνα, στη διαδικασία εκκίνησης ή ακόμη και στο υλικολογισμικό, παρεμβάλλονται στα ερωτήματα προς το σύστημα για να αποκρύψουν αρχεία, διεργασίες και δικτυακές συνδέσεις. Επειδή ένα rootkit ελέγχει τι αναφέρει το παραβιασμένο σύστημα για τον εαυτό του, πρέπει να ανιχνεύεται \emph{εκτός} αυτού του συστήματος —μέσω ασφαλούς εκκίνησης (secure boot), εγκληματολογικής ανάλυσης μνήμης ή ανάλυσης του δίσκου εκτός λειτουργίας.

\section{Ενδείξεις παραβίασης σε υπολογιστές και δίκτυα}\label{sec:4.3}
Μια παραβίαση συνήθως αποκαλύπτεται ως \textbf{απόκλιση από τη φυσιολογική συμπεριφορά αναφοράς} (baseline). Σε έναν υπολογιστή, τυπικές ενδείξεις είναι η παρατεταμένη δραστηριότητα επεξεργαστή, δίσκου ή δικτύου χωρίς αντίστοιχο φόρτο εργασίας του χρήστη· νέοι μηχανισμοί μόνιμης παρουσίας, όπως κλειδιά \emph{Run} στο μητρώο, υπηρεσίες, εργασίες cron ή συνδρομές WMI· ασυνήθιστα δέντρα διεργασιών, για παράδειγμα ένας επεξεργαστής κειμένου που εκκινεί γραμμή εντολών· και στοιχεία παραποίησης των αμυνών, όπως απενεργοποιημένος πράκτορας προστασίας ή διαγραμμένο αρχείο καταγραφής ασφάλειας (συμβάν 1102 στα Windows). Στο δίκτυο, ύποπτες ενδείξεις αποτελούν οι τακτικές συνδέσεις «σηματοδότησης» (beaconing) σε σταθερά διαστήματα, τα πολύ μεγάλα ή φαινομενικά τυχαία ονόματα τομέα και τα αποτυπώματα πελάτη TLS που δεν αντιστοιχούν σε κανένα νόμιμο λογισμικό του οργανισμού.

Μια σημαντική αρχή της ανάλυσης είναι ότι \textbf{μια μεμονωμένη ένδειξη απλώς υποδηλώνει· πολλές συσχετισμένες ενδείξεις επιβεβαιώνουν}. Οι διαχειριστές εκτελούν περιστασιακά PowerShell και το νόμιμο λογισμικό επικοινωνεί μερικές φορές με ασυνήθιστους τομείς. Η βεβαιότητα αυξάνεται όταν ανεξάρτητες παρατηρήσεις οδηγούν στο ίδιο συμπέρασμα —για παράδειγμα, ένα έγγραφο Office εκκινεί PowerShell, το οποίο δημιουργεί ένα κλειδί Run και αρχίζει να επικοινωνεί κάθε εξήντα δευτερόλεπτα με έναν πρόσφατα καταχωρισμένο τομέα. Το Κεφάλαιο 11 μετατρέπει αυτή την αρχή σε κανόνες ανίχνευσης και το Κεφάλαιο 12 την εφαρμόζει στην εγκληματολογική ανάλυση.

\section{Droppers, loaders, trojans απομακρυσμένης πρόσβασης και wipers}\label{sec:4.4}
Οι σύγχρονες εισβολές σπάνια πραγματοποιούνται από ένα μόνο πρόγραμμα. Αντίθετα, μια αλυσίδα εξειδικευμένων στοιχείων μεταβιβάζει τον έλεγχο από το ένα στάδιο στο επόμενο. Ένας \textbf{dropper} μεταφέρει το ωφέλιμο φορτίο μέσα του και το εγγράφει στο σύστημα. Ένας \textbf{downloader} ή \textbf{loader} —γνωστά παραδείγματα είναι τα Emotet, TrickBot και QakBot— ανακτά πρόσθετες μονάδες από το διαδίκτυο μετά την αρχική μόλυνση, γεγονός που επέτρεψε την ανάπτυξη μιας οικονομίας \emph{εγκληματικού λογισμικού ως υπηρεσίας}, στην οποία μια ομάδα πωλεί πρόσβαση σε άλλη. Ένας \textbf{δούρειος ίππος απομακρυσμένης πρόσβασης} (Remote Access Trojan, RAT) παρέχει στον επιτιθέμενο διαδραστικό έλεγχο: καταγραφή πληκτρολογήσεων, λήψη στιγμιοτύπων οθόνης, μεταφορά αρχείων και χρήση του θύματος ως ενδιάμεσου κόμβου.

Ένας \textbf{wiper}, όπως τα Shamoon, WhisperGate ή HermeticWiper, σχεδιάζεται αποκλειστικά για καταστροφή, αν και μπορεί να μεταμφιέζεται σε λυτρισμικό. Η σωστή αναγνώριση είναι εδώ κρίσιμη: επειδή δεν υπάρχει κλειδί αποκρυπτογράφησης, ένα σχέδιο απόκρισης που εξετάζει την καταβολή λύτρων δεν είναι απλώς μάταιο αλλά και επιζήμιο, διότι καθυστερεί την ανάκαμψη από τα αντίγραφα ασφαλείας. Κατά την ανάλυση ενός περιστατικού, οι αναλυτές πρέπει επομένως να κατονομάζουν το \emph{στάδιο} που παρατήρησαν και όχι μόνο την οικογένεια —η διατύπωση «loader του QakBot, πριν από την εγκατάσταση RAT» είναι πολύ πιο χρήσιμη για την ομάδα απόκρισης από το σκέτο «QakBot».

\section{Ανατομία του σύγχρονου λυτρισμικού}\label{sec:4.5}
Το σύγχρονο κρυπτογραφικό λυτρισμικό βασίζεται στην \textbf{υβριδική κρυπτογράφηση}. Κάθε αρχείο κρυπτογραφείται με έναν ταχύ συμμετρικό αλγόριθμο (AES ή ChaCha20) και ένα μοναδικό κλειδί· τα κλειδιά αυτά κρυπτογραφούνται με τη σειρά τους με το δημόσιο κλειδί του επιτιθέμενου (RSA ή ελλειπτικών καμπυλών), ώστε μόνο εκείνος να μπορεί να τα ανακτήσει. Πριν αρχίσει η κρυπτογράφηση, το κακόβουλο λογισμικό διαγράφει συνήθως τα σκιώδη αντίγραφα (shadow copies) και τους καταλόγους αντιγράφων ασφαλείας των Windows (για παράδειγμα με την εντολή \texttt{vssadmin delete shadows}), ώστε να αποτρέψει την εύκολη ανάκτηση, και επιχειρεί να απενεργοποιήσει τα εργαλεία ασφάλειας, ενίοτε επανεκκινώντας το σύστημα σε ασφαλή λειτουργία, όπου οι πράκτορες προστασίας δεν εκτελούνται.

Η πίεση προς το θύμα ασκείται μέσω \textbf{πολλαπλού εκβιασμού}: τα δεδομένα κρυπτογραφούνται, τα κλεμμένα δεδομένα απειλούνται με δημοσιοποίηση σε ιστότοπο διαρροών και ορισμένες ομάδες προσθέτουν επιθέσεις άρνησης υπηρεσίας ή επικοινωνούν απευθείας με τους πελάτες του θύματος. Η σωστή απόκριση ακολουθεί αυστηρή σειρά: πρώτα \textbf{απομονώνονται} οι επηρεασμένοι υπολογιστές και οι παραβιασμένες ταυτότητες· έπειτα \textbf{διαφυλάσσονται τα αποδεικτικά στοιχεία}, συμπεριλαμβανομένων των ειδώλων μνήμης και των σημειωμάτων λύτρων, που περιέχουν αναγνωριστικά του θύματος· στη συνέχεια \textbf{προσδιορίζεται το στέλεχος} και ελέγχεται αν υπάρχει δωρεάν εργαλείο αποκρυπτογράφησης (για παράδειγμα στην πύλη No More Ransom)· και μόνο τότε γίνεται \textbf{αποκατάσταση} από αντίγραφα ασφαλείας που είναι γνωστό ότι δεν έχουν μολυνθεί. Η καταβολή λύτρων δεν εγγυάται ούτε την αποκρυπτογράφηση ούτε τη διαγραφή των κλεμμένων δεδομένων και μπορεί επιπλέον να δημιουργήσει νομικά ζητήματα.

\begin{keybox}{Βασική αρχή}
Το πιο αποτελεσματικό τεχνικό μέτρο απέναντι στο λυτρισμικό είναι ένα αντίγραφο ασφαλείας που ο επιτιθέμενος δεν μπορεί να προσεγγίσει ή να τροποποιήσει: εκτός σύνδεσης ή αμετάβλητο, τακτικά δοκιμασμένο και προστατευμένο με ξεχωριστά διαπιστευτήρια (βλ. Κεφάλαιο 13).
\end{keybox}

\section{Ιστορικές μελέτες περίπτωσης: Stuxnet, WannaCry και Emotet}\label{sec:4.6}
Το \textbf{Stuxnet} (ανακαλύφθηκε το 2010) στόχευε προγραμματιζόμενους λογικούς ελεγκτές της Siemens που λειτουργούσαν φυγοκεντρητές εμπλουτισμού ουρανίου. Χρησιμοποίησε αρκετές ευπάθειες μηδενικής ημέρας, διέσχισε ένα φυσικά απομονωμένο δίκτυο (air gap) μέσω μέσων USB, έφερε οδηγούς υπογεγραμμένους με κλεμμένα πιστοποιητικά και μετέβαλε τη φυσική συμπεριφορά των μηχανημάτων, ενώ εμφάνιζε στους χειριστές κανονικές ενδείξεις. Το δίδαγμά του είναι ότι τα περιβάλλοντα επιχειρησιακής τεχνολογίας δεν μπορούν να βασίζονται μόνο στην απομόνωση (Κεφάλαιο 13).

Το \textbf{WannaCry} (2017) συνδύασε λυτρισμικό με σκουλήκι που εκμεταλλευόταν την ευπάθεια MS17-010 του SMBv1 (την εκμετάλλευση \emph{EternalBlue}). Μέσα σε λίγες ώρες διατάραξε οργανισμούς σε όλο τον κόσμο, μεταξύ των οποίων νοσοκομεία του Εθνικού Συστήματος Υγείας του Ηνωμένου Βασιλείου, παρότι η σχετική ενημέρωση ήταν διαθέσιμη εδώ και δύο μήνες. Το δίδαγμα είναι ότι η καθυστερημένη εφαρμογή ενημερώσεων και τα επίπεδα, μη κατατμημένα δίκτυα αλληλοενισχύονται: μόλις ένα σκουλήκι εισέλθει, τίποτα δεν επιβραδύνει τη διάδοσή του. Το \textbf{Emotet} (2014–2021) εξελίχθηκε από τραπεζικός δούρειος ίππος σε υπηρεσία φόρτωσης (loader-as-a-service) που διανεμόταν μέσω κακόβουλων μηνυμάτων ηλεκτρονικού ταχυδρομείου και εγκαθιστούσε το TrickBot και, τελικά, το λυτρισμικό Ryuk, έως ότου μια διεθνής επιχείρηση των διωκτικών αρχών εξάρθρωσε την υποδομή του. Το δίδαγμά του είναι ότι τα εγκληματικά οικοσυστήματα λειτουργούν όπως οι εφοδιαστικές αλυσίδες και ότι η άμυνα στο αρχικό στάδιο του loader αποτρέπει όλα όσα ακολουθούν.

\section*{Βασικοί όροι}
\begin{description}[style=nextline,leftmargin=1.2em]
  \item[Σκουλήκι (worm)] Αυτοτελές κακόβουλο λογισμικό που διαδίδεται στα δίκτυα χωρίς ενέργεια του χρήστη.
  \item[Δούρειος ίππος] Κακόβουλο λογισμικό μεταμφιεσμένο σε νόμιμο, το οποίο δεν αυτοαναπαράγεται.
  \item[Κακόβουλο λογισμικό χωρίς αρχεία] Κακόβουλος κώδικας που εκτελείται από τη μνήμη ή μέσω εργαλείων του συστήματος χωρίς να γράφει εκτελέσιμα στον δίσκο.
  \item[Rootkit] Κακόβουλο λογισμικό που κρύβει την παρουσία του παρεμβαίνοντας στις αναφορές του ίδιου του λειτουργικού συστήματος.
  \item[Loader] Στάδιο κακόβουλου λογισμικού που κατεβάζει και εκτελεί πρόσθετες μονάδες μετά την αρχική μόλυνση.
  \item[Διπλός εκβιασμός] Τακτική λυτρισμικού που συνδυάζει την κρυπτογράφηση με την απειλή δημοσιοποίησης κλεμμένων δεδομένων.
\end{description}

\section*{Σύνοψη κεφαλαίου}
\begin{itemize}
  \item Το κακόβουλο λογισμικό ταξινομείται καλύτερα με βάση τον μηχανισμό —αναπαραγωγή, διάδοση και ωφέλιμο φορτίο— επειδή κάθε μηχανισμός απαιτεί διαφορετική στρατηγική περιορισμού.
  \item Οι τεχνικές χωρίς αρχεία, ο πολυμορφισμός και τα rootkits παρακάμπτουν τις υπογραφές αρχείων και απαιτούν ανίχνευση συμπεριφοράς και ανίχνευση εκτός του παραβιασμένου συστήματος.
  \item Οι ενδείξεις παραβίασης αποκτούν νόημα μέσω της συσχέτισης· μια μεμονωμένη ανωμαλία είναι απλώς ένα στοιχείο προς διερεύνηση.
  \item Οι σύγχρονες εισβολές είναι σταδιακές αλυσίδες από droppers, loaders, RAT και τελικά φορτία, όπως λυτρισμικό ή wipers.
  \item Τα Stuxnet, WannaCry και Emotet αναδεικνύουν τα όρια της απομόνωσης, το κόστος της καθυστερημένης ενημέρωσης και τον χαρακτήρα εφοδιαστικής αλυσίδας του κυβερνοεγκλήματος.
\end{itemize}

\section*{Ερωτήσεις ανασκόπησης}
\begin{enumerate}
  \item Γιατί η στρατηγική περιορισμού ενός σκουληκιού διαφέρει από εκείνη ενός δούρειου ίππου; Δώστε ένα συγκεκριμένο μέτρο για καθένα.
  \item Εξηγήστε γιατί η εκκένωση διεργασίας (process hollowing) παρακάμπτει τις λίστες επιτρεπόμενων εφαρμογών που βασίζονται σε τιμές κατακερματισμού.
  \item Αναφέρετε τρεις ενδείξεις σε υπολογιστή και δύο δικτυακές ενδείξεις που, συνδυαστικά, θα σας έπειθαν ότι ένας σταθμός εργασίας έχει παραβιαστεί.
  \item Γιατί είναι επικίνδυνο να αντιμετωπιστεί ένας wiper σαν να ήταν λυτρισμικό;
\end{enumerate}

% ============================================================
% Chapter 5: Human Factors, Social Engineering and Digital Safety
% Language: Greek — edit freely; regenerate from src/content if preferred.
% ============================================================
\chapter{Ανθρώπινος Παράγοντας, Κοινωνική Μηχανική και Ψηφιακή Ασφάλεια}\label{ch:5}
\chaptermeta{Ψυχολογία της πειθούς · Ηλεκτρονικό ψάρεμα και οι παραλλαγές του · Φυσικά κανάλια · Stalkerware και μεταδεδομένα · Παραβίαση επιχειρηματικού email · Κουλτούρα ασφάλειας}{Επίπεδο: Εισαγωγικό επίπεδο · Ανθρωποκεντρική προσέγγιση}{Χρόνος μελέτης: 5–7 ώρες}

Τα τεχνικά μέτρα προστατεύουν τα συστήματα, όμως τα συστήματα τα χειρίζονται άνθρωποι, και οι άνθρωποι μπορούν να πειστούν. Η κοινωνική μηχανική εκμεταλλεύεται την εμπιστοσύνη, την προθυμία για βοήθεια, τον φόβο και τη συνήθεια, ώστε ένας νόμιμος χρήστης να εκτελέσει μια ενέργεια που ωφελεί τον επιτιθέμενο —να ανοίξει ένα συνημμένο, να αποκαλύψει έναν κωδικό, να εγκρίνει μια πληρωμή. Οι κλαδικές αναφορές παραβιάσεων διαπιστώνουν σταθερά ότι ο ανθρώπινος παράγοντας εμπλέκεται στην πλειονότητα των επιτυχημένων εισβολών.

Το κεφάλαιο εξηγεί γιατί λειτουργεί η κοινωνική μηχανική, παρουσιάζει τις κύριες μεθόδους επίθεσης και εξετάζει δύο απειλές που έχουν αυξηθεί σημαντικά τα τελευταία χρόνια: την παραβίαση επιχειρηματικού ηλεκτρονικού ταχυδρομείου και την «κόπωση» του ελέγχου ταυτότητας πολλαπλών παραγόντων. Εξετάζει επίσης την προσωπική ψηφιακή ασφάλεια, συμπεριλαμβανομένου του λογισμικού παρακολούθησης (stalkerware) και της διαρροής μεταδεδομένων. Το κεφάλαιο κλείνει με ένα εποικοδομητικό μήνυμα: οι χρήστες δεν είναι ο «πιο αδύναμος κρίκος» που πρέπει να κατηγορείται, αλλά ένα επίπεδο ανίχνευσης που μπορεί να ενισχυθεί μέσω καλού σχεδιασμού και μιας θετικής κουλτούρας αναφοράς.

\begin{outcomesbox}{Μαθησιακά αποτελέσματα}
Μετά τη μελέτη του κεφαλαίου θα πρέπει να μπορείτε:
\begin{itemize}
  \item Να εξηγείτε τις ψυχολογικές αρχές που εκμεταλλεύονται οι κοινωνικοί μηχανικοί.
  \item Να διακρίνετε το ηλεκτρονικό ψάρεμα, το στοχευμένο ψάρεμα, το whaling, το vishing, το smishing, την κατασκευή προσχήματος και το δόλωμα.
  \item Να περιγράφετε τα φυσικά κανάλια επίθεσης και τα κανάλια παρατήρησης, καθώς και τα αντίμετρά τους.
  \item Να αναλύετε μια παραβίαση επιχειρηματικού email και μια επίθεση κόπωσης MFA και να προτείνετε πολυεπίπεδες άμυνες.
  \item Να σχεδιάζετε μέτρα ευαισθητοποίησης και δείκτες που μειώνουν την ευαλωτότητα χωρίς να επιρρίπτουν ευθύνες στους χρήστες.
\end{itemize}
\end{outcomesbox}

\section{Η ψυχολογία της εκμετάλλευσης}\label{sec:5.1}
Η κοινωνική μηχανική επιτυγχάνει επειδή στοχεύει νοητικές συντομεύσεις που υπό κανονικές συνθήκες μάς εξυπηρετούν. Ο ψυχολόγος Robert Cialdini περιέγραψε ορισμένες \textbf{αρχές της πειθούς} που οι επιτιθέμενοι εκμεταλλεύονται συστηματικά. \emph{Αυθεντία}: οι άνθρωποι τείνουν να υπακούουν σε οδηγίες που φαίνεται να προέρχονται από έναν προϊστάμενο, το τμήμα πληροφορικής ή την αστυνομία. \emph{Επείγον και σπανιότητα}: μια προθεσμία («ο λογαριασμός σας θα κλείσει σε μία ώρα») μειώνει τον διαθέσιμο χρόνο για προσεκτική σκέψη. \emph{Κοινωνική απόδειξη}: υποθέτουμε ότι μια ενέργεια είναι ασφαλής αν φαίνεται ότι την κάνουν και άλλοι. \emph{Συμπάθεια} και \emph{αμοιβαιότητα}: είμαστε πιο πρόθυμοι να βοηθήσουμε κάποιον που είναι φιλικός ή που μας έχει κάνει μια μικρή χάρη.

Οι αρχές αυτές είναι ισχυρές επειδή παρακάμπτουν τη συνειδητή σκέψη και ενεργοποιούν γρήγορες, αυτόματες αντιδράσεις. Το άγχος, η κούραση και ο υπερβολικός όγκος πληροφοριών καθιστούν όλους πιο ευάλωτους, ανεξάρτητα από την ευφυΐα ή τις τεχνικές τους γνώσεις. Οι αποτελεσματικές άμυνες δεν βασίζονται, επομένως, στη διαρκή εγρήγορση των ανθρώπων· εισάγουν \emph{τριβή την κατάλληλη στιγμή} —για παράδειγμα, υποχρεωτική επιβεβαιωτική κλήση πριν από οποιαδήποτε αλλαγή τραπεζικών στοιχείων— ώστε η αυτόματη αντίδραση να διακόπτεται ακριβώς όταν έχει σημασία.

\section{Μέθοδοι επίθεσης: από το μαζικό ψάρεμα έως την κατασκευή προσχήματος}\label{sec:5.2}
Το \textbf{ηλεκτρονικό ψάρεμα} (phishing) είναι η μαζική αποστολή δόλιων μηνυμάτων που μιμούνται αξιόπιστους οργανισμούς, με σκοπό την υποκλοπή διαπιστευτηρίων ή την εγκατάσταση κακόβουλου λογισμικού. Το \textbf{στοχευμένο ψάρεμα} (spear-phishing) απευθύνεται σε συγκεκριμένο πρόσωπο ή ομάδα και χρησιμοποιεί στοιχεία που έχουν συλλεχθεί με έρευνα —ονόματα συναδέλφων, τρέχοντα έργα— ώστε να φαίνεται αξιόπιστο. Το \textbf{whaling} είναι στοχευμένο ψάρεμα που απευθύνεται σε ανώτατα στελέχη. Οι ίδιες τεχνικές εφαρμόζονται και σε άλλα κανάλια: το \textbf{vishing} χρησιμοποιεί τηλεφωνικές κλήσεις, συχνά με πλαστογραφημένη αναγνώριση καλούντος και, όλο και περισσότερο, με φωνές που παράγονται από τεχνητή νοημοσύνη, ενώ το \textbf{smishing} χρησιμοποιεί SMS ή εφαρμογές ανταλλαγής μηνυμάτων.

Η \textbf{κατασκευή προσχήματος} (pretexting) συνίσταται στην επινόηση ενός αληθοφανούς σεναρίου —ένας ελεγκτής που χρειάζεται μια αναφορά, ένας τεχνικός που πρέπει να «επαληθεύσει» έναν κωδικό— για την απόκτηση πληροφοριών ή πρόσβασης. Το \textbf{δόλωμα} (baiting) αφήνει ένα ελκυστικό αντικείμενο, όπως ένα USB με την ετικέτα «Μισθοδοσία 2026», σε σημείο όπου το θύμα θα το βρει και θα το συνδέσει στον υπολογιστή του. Το κοινό στοιχείο είναι ότι ο επιτιθέμενος δημιουργεί μια κατάσταση στην οποία η επιβλαβής ενέργεια μοιάζει φυσιολογική, εξυπηρετική ή επείγουσα.

\section{Φυσικά κανάλια και κανάλια παρατήρησης}\label{sec:5.3}
Δεν πραγματοποιούνται όλες οι επιθέσεις κοινωνικής μηχανικής στο διαδίκτυο. Το \textbf{tailgating} (ή \emph{piggybacking}) σημαίνει να ακολουθεί κανείς ένα εξουσιοδοτημένο πρόσωπο μέσα από μια ασφαλισμένη πόρτα, συχνά κρατώντας κούτες, ώστε το να του κρατήσει κάποιος την πόρτα να φαίνεται απλή ευγένεια. Η \textbf{παρακολούθηση πάνω από τον ώμο} (shoulder surfing) είναι η παρατήρηση μιας οθόνης ή ενός πληκτρολογίου για την υποκλοπή κωδικών ή PIN, ενώ η \textbf{έρευνα στα απορρίμματα} (dumpster diving) είναι η αναζήτηση χρήσιμων πληροφοριών σε πεταμένα έγγραφα και εξοπλισμό. Η φυσική πρόσβαση είναι ιδιαίτερα επικίνδυνη, επειδή μπορεί να παρακάμψει εντελώς πολλά δικτυακά μέτρα —ένας επιτιθέμενος που φτάνει σε έναν ξεκλείδωτο σταθμό εργασίας ή σε μια δικτυακή θύρα βρίσκεται ήδη «μέσα».

Τα αντίμετρα συνδυάζουν τεχνολογία και συμπεριφορά: θαλάμους διπλής πόρτας (mantraps) και περιστροφικές πύλες που επιτρέπουν τη διέλευση ενός ατόμου κάθε φορά, εμφανείς κάρτες επισκεπτών, πολιτική καθαρού γραφείου και κλειδωμένης οθόνης, φίλτρα απορρήτου στις οθόνες, καταστροφή εγγράφων σε τεμαχιστές εγκάρσιας κοπής και πιστοποιημένη καταστροφή αποθηκευτικών μέσων. Το προσωπικό πρέπει να αισθάνεται ότι δικαιούται να ζητήσει ευγενικά στοιχεία από ένα άγνωστο πρόσωπο· αυτό γίνεται ευκολότερο όταν η διοίκηση το υποστηρίζει ρητά.

\section{Προσωπική ψηφιακή ασφάλεια: stalkerware και διαρροή μεταδεδομένων}\label{sec:5.4}
Ορισμένες από τις πιο επιβλαβείς επιθέσεις δεν πραγματοποιούνται από απομακρυσμένους εγκληματίες, αλλά από πρόσωπα του στενού περιβάλλοντος του θύματος. Το \textbf{stalkerware} είναι εμπορικά διαθέσιμο λογισμικό το οποίο, μόλις εγκατασταθεί σε ένα κινητό τηλέφωνο, προωθεί κρυφά μηνύματα, τοποθεσία, φωτογραφίες και ιστορικό κλήσεων σε άλλο πρόσωπο. Συνδέεται συχνά με την κακοποίηση από σύντροφο. Προειδοποιητικά σημάδια αποτελούν η ανεξήγητη εξάντληση της μπαταρίας, άγνωστες εφαρμογές με δικαιώματα διαχειριστή συσκευής και ένας σύντροφος που γνωρίζει πράγματα που δεν θα έπρεπε. Επειδή η αφαίρεση του stalkerware μπορεί να ειδοποιήσει τον θύτη και να αυξήσει τον κίνδυνο, τα θύματα πρέπει να αναζητούν υποστήριξη από εξειδικευμένους φορείς πριν προβούν σε οποιαδήποτε ενέργεια.

Τα \textbf{μεταδεδομένα} —δεδομένα για τα δεδομένα— μπορεί να αποκαλύπτουν πολύ περισσότερα από όσα σκοπεύαμε. Τα μεταδεδομένα EXIF μιας φωτογραφίας μπορεί να περιέχουν τις ακριβείς συντεταγμένες GPS του σημείου λήψης, το μοντέλο της συσκευής και την ώρα· τα έγγραφα γραφείου μπορεί να περιλαμβάνουν το όνομα του συντάκτη και το ιστορικό αναθεωρήσεων. Πριν από τη δημόσια κοινοποίηση αρχείων, οι χρήστες πρέπει να αφαιρούν αυτά τα μεταδεδομένα, για παράδειγμα με την εντολή \texttt{exiftool -all=} ή με τις ενσωματωμένες λειτουργίες «κατάργησης ιδιοτήτων» των λειτουργικών συστημάτων και των σουιτών γραφείου.

\section{Παραβίαση επιχειρηματικού email και κόπωση MFA}\label{sec:5.5}
Η \textbf{παραβίαση επιχειρηματικού ηλεκτρονικού ταχυδρομείου} (Business Email Compromise, BEC) είναι μια απάτη κατά την οποία ο επιτιθέμενος υποδύεται ένα ανώτατο στέλεχος, έναν προμηθευτή ή έναν δικηγόρο —είτε παραβιάζοντας ένα πραγματικό γραμματοκιβώτιο είτε καταχωρίζοντας ένα παρόμοιο όνομα τομέα— και ζητά μια επείγουσα πληρωμή ή αλλαγή τραπεζικών στοιχείων. Η BEC συχνά δεν περιλαμβάνει καθόλου κακόβουλο λογισμικό και, επομένως, διαφεύγει από τα τεχνικά φίλτρα· βασίζεται αποκλειστικά στην αυθεντία και στο επείγον. Οι αποτελεσματικές άμυνες είναι τόσο διαδικαστικές όσο και τεχνικές: οι αλλαγές στοιχείων πληρωμής επαληθεύονται μέσω γνωστού αριθμού τηλεφώνου, οι μεταφορές υψηλής αξίας απαιτούν έγκριση από δύο πρόσωπα (\emph{δημιουργός–ελεγκτής}) και ο έλεγχος αυθεντικότητας email (SPF, DKIM και πολιτική DMARC \texttt{p=reject}) δυσχεραίνει σημαντικά την πλαστογράφηση τομέα.

Η \textbf{κόπωση MFA} (MFA fatigue ή \emph{push bombing}) στοχεύει χρήστες που επαληθεύουν την ταυτότητά τους μέσω ειδοποιήσεων push. Έχοντας υποκλέψει έναν κωδικό, ο επιτιθέμενος προκαλεί δεκάδες αιτήματα έγκρισης σύνδεσης, συχνά αργά τη νύχτα, έως ότου ο εξαντλημένος χρήστης πατήσει «Έγκριση» —μερικές φορές μετά από τηλεφώνημα ενός ψεύτικου γραφείου υποστήριξης. Τα αντίμετρα περιλαμβάνουν την \emph{αντιστοίχιση αριθμού} (ο χρήστης πρέπει να πληκτρολογήσει έναν κωδικό που εμφανίζεται στην οθόνη σύνδεσης), όρια στον αριθμό των αιτημάτων και μεθόδους ανθεκτικές στο ψάρεμα, όπως τα κλειδιά ασφαλείας FIDO2 (Κεφάλαιο 8).

\section{Οικοδόμηση κουλτούρας ασφάλειας που ενθαρρύνει την αναφορά}\label{sec:5.6}
Τα προγράμματα ευαισθητοποίησης συχνά μετρούν το λάθος πράγμα. Το \textbf{ποσοστό κλικ} στις προσομοιωμένες εκστρατείες ηλεκτρονικού ψαρέματος αντανακλά κυρίως τη δυσκολία της προσομοίωσης· μπορεί να αυξηθεί ή να μειωθεί κατά βούληση. Πολύ πιο ουσιαστικός δείκτης είναι το \textbf{ποσοστό αναφοράς} —το ποσοστό των χρηστών που αναφέρουν ένα ύποπτο μήνυμα— σε συνδυασμό με τον \textbf{χρόνο έως την πρώτη αναφορά}. Μία μόνο έγκαιρη αναφορά επιτρέπει στην ομάδα ασφάλειας να αφαιρέσει ένα κακόβουλο μήνυμα από όλα τα γραμματοκιβώτια πριν το ανοίξουν οι περισσότεροι παραλήπτες.

Η κουλτούρα καθορίζει αν οι άνθρωποι θα προβούν σε αναφορά. Αν οι εργαζόμενοι που πατούν σε έναν σύνδεσμο ψαρέματος εκτίθενται δημόσια ή τιμωρούνται, θα αποκρύπτουν τα λάθη τους και ο οργανισμός θα χάνει την πιο έγκαιρη προειδοποίησή του. Μια κουλτούρα που ενθαρρύνει την αναφορά την καθιστά εύκολη (ένα κουμπί στο πρόγραμμα ηλεκτρονικού ταχυδρομείου), ευχαριστεί τους ανθρώπους για κάθε αναφορά —ακόμη και για τους ψευδείς συναγερμούς— και αντιμετωπίζει τα λάθη ως ευκαιρίες μάθησης. Με αυτόν τον τρόπο οι χρήστες γίνονται ενεργό επίπεδο ανίχνευσης και όχι αδυναμία του οργανισμού.

\section*{Βασικοί όροι}
\begin{description}[style=nextline,leftmargin=1.2em]
  \item[Κοινωνική μηχανική] Χειραγώγηση ανθρώπων ώστε να εκτελέσουν ενέργειες ή να αποκαλύψουν πληροφορίες προς όφελος του επιτιθέμενου.
  \item[Στοχευμένο ψάρεμα] Ηλεκτρονικό ψάρεμα προσαρμοσμένο σε συγκεκριμένο πρόσωπο ή ομάδα με στοιχεία που συλλέχθηκαν με έρευνα.
  \item[Κατασκευή προσχήματος] Επινόηση αληθοφανούς σεναρίου για την απόκτηση πληροφοριών ή πρόσβασης.
  \item[BEC] Παραβίαση επιχειρηματικού email: απάτη με υπόδυση αξιόπιστων προσώπων για εκτροπή πληρωμών.
  \item[Κόπωση MFA] Κατάκλυση ενός χρήστη με αιτήματα επαλήθευσης έως ότου εγκρίνει κάποιο.
  \item[DMARC] Πολιτική email που ορίζει στους παραλήπτες πώς να χειρίζονται μηνύματα που αποτυγχάνουν στους ελέγχους SPF/DKIM.
\end{description}

\section*{Σύνοψη κεφαλαίου}
\begin{itemize}
  \item Η κοινωνική μηχανική εκμεταλλεύεται αυτόματες αντιδράσεις απέναντι στην αυθεντία, το επείγον, την κοινωνική απόδειξη, τη συμπάθεια και την αμοιβαιότητα.
  \item Το ηλεκτρονικό ψάρεμα εμφανίζεται σε πολλές μορφές και κανάλια· η κατασκευή προσχήματος και το δόλωμα δημιουργούν καταστάσεις όπου οι επιβλαβείς ενέργειες μοιάζουν φυσιολογικές.
  \item Η φυσική πρόσβαση μπορεί να παρακάμψει τα δικτυακά μέτρα και πρέπει να προστατεύεται τόσο με τεχνολογία όσο και με τη συμπεριφορά του προσωπικού.
  \item Η BEC και η κόπωση MFA αντιμετωπίζονται με διαδικασίες επαλήθευσης, διπλή έγκριση, έλεγχο αυθεντικότητας email και MFA ανθεκτικό στο ψάρεμα.
  \item Ο κρίσιμος δείκτης είναι το ποσοστό αναφοράς και όχι το ποσοστό κλικ· μια κουλτούρα χωρίς απόδοση ευθυνών μετατρέπει τους χρήστες σε επίπεδο ανίχνευσης.
\end{itemize}

\section*{Ερωτήσεις ανασκόπησης}
\begin{enumerate}
  \item Εντοπίστε τις αρχές πειθούς που χρησιμοποιούνται στο μήνυμα: «Είμαι ο Διευθύνων Σύμβουλος. Χρειάζομαι την πληρωμή του προμηθευτή μέσα στα επόμενα 30 λεπτά —μην τηλεφωνήσετε, είμαι σε σύσκεψη».
  \item Γιατί η BEC δύσκολα σταματά με λογισμικό προστασίας από ιούς ή με περιβάλλοντα απομονωμένης εκτέλεσης; Ποια μέτρα είναι αποτελεσματικά;
  \item Εξηγήστε πώς η αντιστοίχιση αριθμού αποτρέπει μια βασική επίθεση κόπωσης MFA.
  \item Τεκμηριώστε γιατί το ποσοστό αναφοράς είναι καλύτερος δείκτης ευαισθητοποίησης από το ποσοστό κλικ.
\end{enumerate}

% ============================================================
% Chapter 6: Network Security: Attacks, Defences and Botnet Topologies
% Language: Greek — edit freely; regenerate from src/content if preferred.
% ============================================================
\chapter{Ασφάλεια Δικτύων: Επιθέσεις, Άμυνες και Τοπολογίες Botnet}\label{ch:6}
\chaptermeta{DoS και DDoS · Ενίσχυση · Botnet και διοίκηση-έλεγχος · Επιθέσεις ενδιάμεσου · Τείχη προστασίας · VPN · Κατάτμηση, NAC και ασφάλεια ασύρματων δικτύων}{Επίπεδο: Μεσαίο επίπεδο · Δίκτυα}{Χρόνος μελέτης: 8–10 ώρες}

Τα δίκτυα συνδέουν τα πάντα, και ό,τι συνδέεται μπορεί επίσης να δεχθεί επίθεση. Το κεφάλαιο εξετάζει πώς οι αντίπαλοι κάνουν κατάχρηση των δικτυακών πρωτοκόλλων για να διαταράξουν τη διαθεσιμότητα, να υποκλέψουν ή να αλλοιώσουν την κίνηση και να συντονίσουν χιλιάδες παραβιασμένες συσκευές. Προϋποθέτει βασική εξοικείωση με το μοντέλο TCP/IP· όπου μια λεπτομέρεια πρωτοκόλλου είναι απαραίτητη, εξηγείται στο κείμενο.

Το πρώτο μισό του κεφαλαίου εστιάζει στις επιθέσεις: την άρνηση υπηρεσίας και τις κατανεμημένες και ενισχυμένες μορφές της, την αρχιτεκτονική και τον κύκλο ζωής των botnet, καθώς και τις τεχνικές ενδιάμεσου (man-in-the-middle). Το δεύτερο μισό στρέφεται στην άμυνα. Δείχνει πώς λειτουργεί στην πράξη η πολυεπίπεδη αντιμετώπιση επιθέσεων DDoS και παρουσιάζει τα δομικά στοιχεία της δικτυακής άμυνας —τείχη προστασίας, εικονικά ιδιωτικά δίκτυα, κατάτμηση, έλεγχο πρόσβασης στο δίκτυο και ασφάλεια ασύρματων δικτύων— με σαφή διατύπωση του τι μπορεί και τι δεν μπορεί να επιτύχει το καθένα.

\begin{outcomesbox}{Μαθησιακά αποτελέσματα}
Μετά τη μελέτη του κεφαλαίου θα πρέπει να μπορείτε:
\begin{itemize}
  \item Να διακρίνετε τις επιθέσεις άρνησης υπηρεσίας ογκομετρικού τύπου, επιπέδου πρωτοκόλλου και επιπέδου εφαρμογής και να εξηγείτε την ενίσχυση (amplification).
  \item Να συγκρίνετε τις κεντρικές και τις ομότιμες αρχιτεκτονικές botnet και να περιγράφετε τον κύκλο ζωής ενός botnet.
  \item Να εξηγείτε την πλαστογράφηση ARP, την πλαστογράφηση DNS και την απογύμνωση TLS, καθώς και τα μέτρα που τις αποτρέπουν.
  \item Να σχεδιάζετε μια πολυεπίπεδη άμυνα κατά DDoS, από την άκρη του δικτύου έως την εφαρμογή.
  \item Να περιγράφετε τον ρόλο και τα όρια των τειχών προστασίας, των VPN, της κατάτμησης, του 802.1X και του WPA3.
\end{itemize}
\end{outcomesbox}

\section{Άρνηση υπηρεσίας: DoS, DDoS και ενίσχυση}\label{sec:6.1}
Μια επίθεση \textbf{άρνησης υπηρεσίας} (Denial of Service, DoS) στοχεύει να καταστήσει μια υπηρεσία μη διαθέσιμη στους νόμιμους χρήστες της. Όταν η κίνηση της επίθεσης προέρχεται ταυτόχρονα από πολλές πηγές, ονομάζεται \textbf{κατανεμημένη άρνηση υπηρεσίας} (Distributed Denial of Service, DDoS) και γίνεται πολύ δυσκολότερο να αποκλειστεί, επειδή δεν υπάρχει μία πηγή που να μπορεί απλώς να φιλτραριστεί. Οι επιθέσεις DDoS χωρίζονται σε τρεις γενικές οικογένειες. Οι \textbf{ογκομετρικές} κατακλύζουν το εύρος ζώνης του θύματος. Οι επιθέσεις \textbf{επιπέδου πρωτοκόλλου} εξαντλούν τους πόρους κατάστασης διακομιστών ή δικτυακών συσκευών· η κλασική \emph{πλημμύρα SYN} στέλνει πολλά αιτήματα σύνδεσης TCP χωρίς να ολοκληρώνει τη χειραψία, γεμίζοντας τον πίνακα ημιανοιχτών συνδέσεων του διακομιστή. Οι επιθέσεις \textbf{επιπέδου εφαρμογής} στέλνουν φαινομενικά νόμιμα αλλά απαιτητικά αιτήματα —όπως ερωτήματα αναζήτησης— που εξαντλούν την ίδια την εφαρμογή.

Η \textbf{ενίσχυση} (amplification) επιτρέπει σε έναν επιτιθέμενο με περιορισμένο εύρος ζώνης να παράγει τεράστιες πλημμύρες κίνησης. Ο επιτιθέμενος στέλνει μικρά αιτήματα με πλαστογραφημένη διεύθυνση προέλευσης (τη διεύθυνση του θύματος) σε δημόσιους διακομιστές, οι οποίοι επιστρέφουν πολύ μεγαλύτερες απαντήσεις. Πρωτόκολλα όπως τα DNS, NTP και memcached έχουν χρησιμοποιηθεί με αυτόν τον τρόπο, με συντελεστές ενίσχυσης από μερικές δεκάδες έως δεκάδες χιλιάδες. Δύο μέτρα αντιμετωπίζουν τη ρίζα του προβλήματος: τα δίκτυα πρέπει να απορρίπτουν εξερχόμενα πακέτα με πλαστογραφημένες διευθύνσεις προέλευσης (\emph{φιλτράρισμα εισόδου}, BCP 38) και οι διαχειριστές δεν πρέπει να εκθέτουν στο διαδίκτυο ανοιχτούς αναλυτές DNS ή άλλους ενισχυτές.

\section{Botnet: αρχιτεκτονικές, κύκλος ζωής και στρατολόγηση συσκευών IoT}\label{sec:6.2}
Ένα \textbf{botnet} είναι ένα δίκτυο παραβιασμένων συσκευών («bots») που ελέγχονται απομακρυσμένα από έναν διαχειριστή. Σε μια \textbf{κεντρική} αρχιτεκτονική, τα bots λαμβάνουν εντολές από έναν ή λίγους διακομιστές διοίκησης και ελέγχου (Command and Control, C2), ιστορικά μέσω IRC και σήμερα κυρίως μέσω HTTPS. Ο σχεδιασμός αυτός είναι απλός αλλά εύθραυστος: η κατάσχεση των διακομιστών C2 απενεργοποιεί το botnet. Τα \textbf{ομότιμα} (peer-to-peer) botnet κατανέμουν τις εντολές στα ίδια τα bots, κάτι που καθιστά την εξάρθρωσή τους πολύ δυσκολότερη. Οι διαχειριστές χρησιμοποιούν επίσης \emph{αλγορίθμους παραγωγής ονομάτων τομέα} (DGA), οι οποίοι δημιουργούν χιλιάδες υποψήφια ονόματα κάθε ημέρα, ώστε οι αμυνόμενοι να μην μπορούν να τα αποκλείσουν όλα εκ των προτέρων.

Ο κύκλος ζωής ενός botnet περιλαμβάνει τη \textbf{στρατολόγηση} (μόλυνση νέων συσκευών), τη \textbf{διοίκηση και τον έλεγχο}, την \textbf{αξιοποίηση} (DDoS επί πληρωμή, ανεπιθύμητη αλληλογραφία, μαζική δοκιμή διαπιστευτηρίων, εξόρυξη κρυπτονομισμάτων) και τη \textbf{συντήρηση} (ενημέρωση και προστασία των bots). Το botnet \textbf{Mirai} του 2016 έδειξε πόσο αδύναμο ήταν το Διαδίκτυο των Πραγμάτων: απλώς δοκίμαζε μια λίστα περίπου εξήντα εργοστασιακών ονομάτων χρήστη και κωδικών σε κάμερες και δρομολογητές, στρατολόγησε εκατοντάδες χιλιάδες συσκευές και εξαπέλυσε επιθέσεις που διατάραξαν σημαντικές υπηρεσίες του διαδικτύου. Το δίδαγμα για κατασκευαστές και διαχειριστές είναι σαφές —μοναδικά διαπιστευτήρια, αυτόματες ενημερώσεις και καμία περιττή απομακρυσμένη υπηρεσία.

\section{Υποκλοπή και αλλοίωση: επιθέσεις ενδιάμεσου}\label{sec:6.3}
Σε μια \textbf{επίθεση ενδιάμεσου} (Man-in-the-Middle, MITM), ο αντίπαλος παρεμβάλλεται ανάμεσα σε δύο μέρη που επικοινωνούν, ώστε η κίνηση να διέρχεται από ένα σύστημα που ελέγχει. Σε ένα τοπικό δίκτυο, η \textbf{πλαστογράφηση ARP} το επιτυγχάνει αποστέλλοντας ψευδή μηνύματα που συσχετίζουν τη φυσική διεύθυνση του επιτιθέμενου με τη διεύθυνση IP της προεπιλεγμένης πύλης· τα θύματα στέλνουν στη συνέχεια την κίνησή τους στον επιτιθέμενο. Η \textbf{πλαστογράφηση DNS} επιστρέφει ψευδείς απαντήσεις σε ερωτήματα ονομάτων, ανακατευθύνοντας τους χρήστες σε κακόβουλους διακομιστές. Σε δημόσια δίκτυα Wi-Fi, ένα σημείο πρόσβασης \emph{«κακός δίδυμος»} (evil twin) με γνώριμο όνομα δικτύου μπορεί να προσελκύσει ανυποψίαστους χρήστες.

Μόλις βρεθεί στη μέση, ο επιτιθέμενος μπορεί είτε να \textbf{υποκλέπτει} παθητικά τη μη κρυπτογραφημένη κίνηση είτε να την τροποποιεί ενεργά —για παράδειγμα μέσω της \emph{απογύμνωσης TLS} (TLS stripping), η οποία υποβαθμίζει μια σύνδεση από HTTPS σε HTTP. Η κρυπτογραφία είναι το θεμελιώδες αντίμετρο: το σωστά επικυρωμένο TLS (Κεφάλαιο 7) καθιστά την υποκλαπείσα κίνηση μη αναγνώσιμη και την αλλοίωση ανιχνεύσιμη, ενώ ο μηχανισμός \textbf{HTTP Strict Transport Security (HSTS)} αποτρέπει την υποβάθμιση. Σε επίπεδο δικτύου, η \emph{δυναμική επιθεώρηση ARP} και η \emph{παρακολούθηση DHCP} (DHCP snooping) στους μεταγωγείς, η επικύρωση DNSSEC και ο έλεγχος ταυτότητας 802.1X περιορίζουν τις ευκαιρίες υποκλοπής.

\section{Πολυεπίπεδη άμυνα κατά DDoS στην πράξη}\label{sec:6.4}
Καμία μεμονωμένη συσκευή δεν μπορεί να απορροφήσει κάθε επίθεση DDoS· γι' αυτό η άμυνα οργανώνεται σε επίπεδα, καθένα από τα οποία χειρίζεται αυτό που κάνει καλύτερα. Στο \textbf{ανάντη} επίπεδο, οι πάροχοι διαδικτύου και οι υπηρεσίες «καθαρισμού» κίνησης στο νέφος απορροφούν τις ογκομετρικές πλημμύρες, αξιοποιώντας χωρητικότητα πολύ μεγαλύτερη από τη σύνδεση οποιουδήποτε μεμονωμένου οργανισμού· η δρομολόγηση anycast κατανέμει το φορτίο σε πολλά κέντρα δεδομένων. Στην \textbf{άκρη του δικτύου}, οι δρομολογητές εφαρμόζουν λίστες πρόσβασης και όρια ρυθμού, ενώ η \emph{απομακρυσμένα ενεργοποιούμενη «μαύρη τρύπα»} (remote-triggered black-holing) μπορεί να θυσιάσει μία διεύθυνση που δέχεται επίθεση για να προστατεύσει το υπόλοιπο δίκτυο.

Πιο κοντά στην υπηρεσία, \textbf{άμυνες επιπέδου πρωτοκόλλου}, όπως τα \emph{SYN cookies}, επιτρέπουν στους διακομιστές να χειρίζονται πλημμύρες ημιανοιχτών συνδέσεων χωρίς να αποθηκεύουν κατάσταση. Στο επίπεδο της \textbf{εφαρμογής}, τα δίκτυα διανομής περιεχομένου (CDN), τα τείχη προστασίας εφαρμογών ιστού (WAF), η προσωρινή αποθήκευση και ο περιορισμός ρυθμού ανά πελάτη φιλτράρουν τα απαιτητικά αιτήματα, ενώ δοκιμασίες όπως τα CAPTCHA διαχωρίζουν τους ανθρώπους από τα bots. Τέλος, ένα τεκμηριωμένο \textbf{εγχειρίδιο ενεργειών} (runbook) —ποιον καλούμε, πώς ενεργοποιούμε τον καθαρισμό κίνησης, πώς επικοινωνούμε με τους πελάτες— εξασφαλίζει ότι ο οργανισμός αντιδρά σε λεπτά και όχι σε ώρες.

\section{Τείχη προστασίας και εικονικά ιδιωτικά δίκτυα}\label{sec:6.5}
Ένα \textbf{τείχος προστασίας} (firewall) επιβάλλει μια πολιτική σχετικά με το ποια κίνηση επιτρέπεται να διέρχεται μεταξύ ζωνών του δικτύου. Τα \emph{φίλτρα πακέτων} εξετάζουν μεμονωμένα πακέτα με βάση τη διεύθυνση, τη θύρα και το πρωτόκολλο. Τα τείχη προστασίας \emph{με διατήρηση κατάστασης} (stateful) παρακολουθούν τις συνδέσεις, ώστε μια απάντηση να επιτρέπεται μόνο αν αντιστοιχεί σε αίτημα που έχει σταλεί. Τα \emph{τείχη προστασίας νέας γενιάς} αναγνωρίζουν επιπλέον εφαρμογές και χρήστες και μπορεί να ενσωματώνουν πρόληψη εισβολών. Ανεξάρτητα από τη γενιά, ένας κανόνας σχεδίασης είναι θεμελιώδης: η πολιτική πρέπει να \textbf{απορρίπτει εξ ορισμού} και να επιτρέπει μόνο τις ρητά αναγκαίες ροές, σύμφωνα με την αρχή των ασφαλών προεπιλογών.

Ένα \textbf{εικονικό ιδιωτικό δίκτυο} (Virtual Private Network, VPN) δημιουργεί ένα κρυπτογραφημένο τούνελ μέσα από ένα μη αξιόπιστο δίκτυο. Το \textbf{IPsec} λειτουργεί στο επίπεδο δικτύου και χρησιμοποιείται συχνά για συνδέσεις μεταξύ εγκαταστάσεων· τα \textbf{VPN που βασίζονται σε TLS} είναι βολικά για απομακρυσμένους χρήστες· το \textbf{WireGuard} είναι ένα σύγχρονο πρωτόκολλο με σκόπιμα μικρό όγκο κώδικα και σταθερή, σύγχρονη κρυπτογραφία. Είναι σημαντικό να κατανοηθεί τι \emph{δεν} κάνει ένα VPN: προστατεύει την κίνηση κατά τη μεταφορά, αλλά δεν καθιστά αξιόπιστη τη συσκευή που συνδέεται. Ένας παραβιασμένος φορητός υπολογιστής που συνδέεται μέσω VPN φέρνει τον επιτιθέμενο μέσα στο δίκτυο —ένας από τους λόγους για τους οποίους οι οργανισμοί μεταβαίνουν σε πρόσβαση μηδενικής εμπιστοσύνης (Κεφάλαιο 13).

\section{Κατάτμηση, έλεγχος πρόσβασης στο δίκτυο και ασφάλεια ασύρματων δικτύων}\label{sec:6.6}
Η \textbf{κατάτμηση δικτύου} (segmentation) διαιρεί ένα δίκτυο σε ζώνες —για παράδειγμα σταθμούς εργασίας χρηστών, διακομιστές, διεπαφές διαχείρισης και Wi-Fi επισκεπτών— και επιτρέπει μόνο την κίνηση που πραγματικά χρειάζεται κάθε ζώνη. Η κατάτμηση δεν αποτρέπει την πρώτη παραβίαση, περιορίζει όμως το πόσο μακριά μπορεί να κινηθεί ο επιτιθέμενος στη συνέχεια· η περίπτωση του WannaCry στο Κεφάλαιο 4 έδειξε το κόστος των επίπεδων δικτύων. Η \emph{μικροκατάτμηση} εφαρμόζει την ίδια ιδέα έως το επίπεδο των μεμονωμένων φόρτων εργασίας.

Ο \textbf{έλεγχος πρόσβασης στο δίκτυο} (Network Access Control, NAC) αποφασίζει αν μια συσκευή επιτρέπεται καν να συνδεθεί στο δίκτυο. Το πρότυπο \textbf{IEEE 802.1X} απαιτεί από κάθε συσκευή να επαληθεύσει την ταυτότητά της —συνήθως με πιστοποιητικό— πριν ενεργοποιηθεί η θύρα του μεταγωγέα ή η ασύρματη σύνδεση, και μπορεί να τοποθετεί μη συμμορφούμενες συσκευές σε δίκτυο καραντίνας. Η ασφάλεια των ασύρματων δικτύων έχει εξελιχθεί σημαντικά: το WEP είναι πλήρως παραβιασμένο, το WPA2 είναι ευάλωτο σε μαντεψιά κωδικών εκτός σύνδεσης όταν χρησιμοποιούνται αδύναμες φράσεις πρόσβασης, ενώ το \textbf{WPA3} εισάγει τη χειραψία SAE, η οποία αντιστέκεται σε τέτοιες μαντεψιές και παρέχει εμπρόσθια μυστικότητα. Στις επιχειρήσεις, η συνιστώμενη ρύθμιση παραμένει το WPA2/WPA3-Enterprise με 802.1X.

\section*{Βασικοί όροι}
\begin{description}[style=nextline,leftmargin=1.2em]
  \item[DDoS] Κατανεμημένη άρνηση υπηρεσίας: επίθεση κατά της διαθεσιμότητας από πολλές πηγές ταυτόχρονα.
  \item[Ενίσχυση] Αξιοποίηση τρίτων διακομιστών που επιστρέφουν μεγάλες απαντήσεις σε μικρά, πλαστογραφημένα αιτήματα.
  \item[Botnet / C2] Δίκτυο παραβιασμένων συσκευών και το κανάλι διοίκησης και ελέγχου που τις κατευθύνει.
  \item[Πλαστογράφηση ARP] Αποστολή ψευδών μηνυμάτων ARP για την ανακατεύθυνση της τοπικής κίνησης μέσω του επιτιθέμενου.
  \item[Τείχος προστασίας με κατάσταση] Τείχος προστασίας που παρακολουθεί τις συνδέσεις και επιτρέπει απαντήσεις μόνο σε νόμιμα αιτήματα.
  \item[802.1X] Έλεγχος πρόσβασης ανά θύρα που επαληθεύει την ταυτότητα των συσκευών πριν τους παραχωρήσει συνδεσιμότητα.
\end{description}

\section*{Σύνοψη κεφαλαίου}
\begin{itemize}
  \item Οι επιθέσεις DDoS στοχεύουν το εύρος ζώνης, την κατάσταση των πρωτοκόλλων ή τους πόρους της εφαρμογής· η ενίσχυση πολλαπλασιάζει τη δύναμη του επιτιθέμενου μέσω πλαστογράφησης.
  \item Τα botnet ελέγχονται μέσω κεντρικής ή ομότιμης διοίκησης-ελέγχου και αξιοποιούνται με πολλούς τρόπους· το Mirai ανέδειξε την αδυναμία των εργοστασιακών διαπιστευτηρίων στις συσκευές IoT.
  \item Οι επιθέσεις ενδιάμεσου εκμεταλλεύονται την εμπιστοσύνη στα τοπικά πρωτόκολλα· τα κύρια αντίμετρα είναι το επικυρωμένο TLS, το HSTS και οι προστασίες στους μεταγωγείς.
  \item Η άμυνα κατά DDoS είναι πολυεπίπεδη, από τον καθαρισμό κίνησης στο ανάντη δίκτυο έως τον περιορισμό ρυθμού στην εφαρμογή, και εξαρτάται από ένα προετοιμασμένο εγχειρίδιο ενεργειών.
  \item Τα τείχη προστασίας, τα VPN, η κατάτμηση και ο NAC επιλύουν συγκεκριμένα προβλήματα —κανένα δεν καθιστά αξιόπιστη μια παραβιασμένη συσκευή.
\end{itemize}

\section*{Ερωτήσεις ανασκόπησης}
\begin{enumerate}
  \item Εξηγήστε γιατί ένας επιτιθέμενος που χρησιμοποιεί ενίσχυση DNS πρέπει να πλαστογραφεί τη διεύθυνση IP προέλευσης των αιτημάτων του.
  \item Γιατί τα ομότιμα botnet εξαρθρώνονται δυσκολότερα από τα κεντρικά;
  \item Περιγράψτε πώς το HSTS αποτρέπει μια επίθεση απογύμνωσης TLS σε δημόσιο δίκτυο Wi-Fi.
  \item Μια εταιρεία δηλώνει: «Όλοι οι απομακρυσμένοι εργαζόμενοι χρησιμοποιούν το VPN μας, άρα οι φορητοί τους υπολογιστές είναι ασφαλείς». Αξιολογήστε κριτικά αυτή τη δήλωση.
\end{enumerate}

\part{Μηχανισμοί Προστασίας και Ασφαλής Μηχανική}
% ============================================================
% Chapter 7: Cryptographic Mechanisms and Public Key Infrastructure
% Language: Greek — edit freely; regenerate from src/content if preferred.
% ============================================================
\chapter{Κρυπτογραφικοί Μηχανισμοί και Υποδομή Δημόσιου Κλειδιού}\label{ch:7}
\chaptermeta{Συμμετρική και ασύμμετρη κρυπτογράφηση · Συναρτήσεις κατακερματισμού · Ψηφιακές υπογραφές · ΥΔΚ και X.509 · TLS 1.3 · Διαχείριση κλειδιών}{Επίπεδο: Μεσαίο επίπεδο · Εφαρμοσμένη κρυπτογραφία}{Χρόνος μελέτης: 9–11 ώρες}

Η κρυπτογραφία αποτελεί το μαθηματικό θεμέλιο της εμπιστευτικότητας, της ακεραιότητας και της αυθεντικότητας στα ψηφιακά συστήματα. Το Κεφάλαιο 6 ολοκληρώθηκε με επιθέσεις που υποκλέπτουν και αλλοιώνουν την κίνηση· το παρόν κεφάλαιο παρέχει τα εργαλεία που καθιστούν άχρηστη μια τέτοια υποκλοπή. Στόχος δεν είναι η παραγωγή των υποκείμενων μαθηματικών, αλλά η ακριβής κατανόηση του τι εγγυάται κάθε κρυπτογραφικό δομικό στοιχείο, πώς συνδυάζονται τα στοιχεία αυτά και πού αποτυγχάνουν τα πραγματικά συστήματα.

Το κεφάλαιο αντιπαραβάλλει αρχικά τη συμμετρική με την ασύμμετρη κρυπτογράφηση και εξηγεί τις συναρτήσεις κατακερματισμού και τις ψηφιακές υπογραφές. Στη συνέχεια δείχνει πώς η υποδομή δημόσιου κλειδιού (ΥΔΚ ή PKI) συνδέει κλειδιά με ταυτότητες μέσω πιστοποιητικών και αναλύει τη χειραψία του TLS 1.3, που προστατεύει σήμερα το μεγαλύτερο μέρος της κίνησης στο διαδίκτυο. Κλείνει με τη διαχείριση κλειδιών —το πεδίο όπου, στην πράξη, σημειώνονται οι περισσότερες κρυπτογραφικές αποτυχίες.

\begin{outcomesbox}{Μαθησιακά αποτελέσματα}
Μετά τη μελέτη του κεφαλαίου θα πρέπει να μπορείτε:
\begin{itemize}
  \item Να εξηγείτε τη διαφορά μεταξύ συμμετρικής και ασύμμετρης κρυπτογράφησης και γιατί τα πρακτικά συστήματα τις συνδυάζουν.
  \item Να διατυπώνετε τις ιδιότητες ασφάλειας των κρυπτογραφικών συναρτήσεων κατακερματισμού και να διακρίνετε τον κατακερματισμό από την κρυπτογράφηση.
  \item Να περιγράφετε πώς οι ψηφιακές υπογραφές παρέχουν ακεραιότητα, αυθεντικότητα και (υπό προϋποθέσεις) μη αποποίηση.
  \item Να εξηγείτε τον ρόλο των αρχών πιστοποίησης, των αρχών καταχώρισης, των αλυσίδων πιστοποιητικών και της ανάκλησης σε μια ΥΔΚ.
  \item Να σκιαγραφείτε τη χειραψία του TLS 1.3 και να εντοπίζετε συνήθεις αστοχίες στη διαχείριση κλειδιών.
\end{itemize}
\end{outcomesbox}

\section{Συμμετρική και ασύμμετρη κρυπτογράφηση}\label{sec:7.1}
Στη \textbf{συμμετρική κρυπτογράφηση} χρησιμοποιείται το ίδιο μυστικό κλειδί για την κρυπτογράφηση και την αποκρυπτογράφηση. Σύγχρονοι συμμετρικοί αλγόριθμοι, όπως οι \textbf{AES} και \textbf{ChaCha20}, είναι εξαιρετικά γρήγοροι και, με κλειδιά 128 ή 256 bit, θεωρούνται ασφαλείς απέναντι σε όλες τις γνωστές πρακτικές επιθέσεις. Η αδυναμία τους είναι οργανωτική: τα δύο μέρη πρέπει να μοιράζονται ήδη το κλειδί, ενώ ένα δίκτυο \emph{n} συμμετεχόντων θα χρειαζόταν περίπου n²/2 διαφορετικά κλειδιά. Οι συμμετρικοί αλγόριθμοι πρέπει επίσης να χρησιμοποιούνται με κατάλληλο \emph{τρόπο λειτουργίας}· οι τρόποι με ενσωματωμένη αυθεντικοποίηση, όπως ο AES-GCM, προστατεύουν τόσο την εμπιστευτικότητα όσο και την ακεραιότητα.

Η \textbf{ασύμμετρη κρυπτογράφηση} (κρυπτογράφηση δημόσιου κλειδιού) χρησιμοποιεί ένα ζεύγος μαθηματικά συσχετισμένων κλειδιών. Το \textbf{δημόσιο κλειδί} μπορεί να διανέμεται ελεύθερα, ενώ το \textbf{ιδιωτικό κλειδί} παραμένει μυστικό. Δεδομένα που κρυπτογραφούνται με το δημόσιο κλειδί αποκρυπτογραφούνται μόνο με το ιδιωτικό. Αλγόριθμοι όπως ο \textbf{RSA} και η \textbf{κρυπτογραφία ελλειπτικών καμπυλών (ECC)} επιλύουν το πρόβλημα της διανομής κλειδιών, είναι όμως χιλιάδες φορές πιο αργοί από τους συμμετρικούς. Γι' αυτό τα πρακτικά συστήματα χρησιμοποιούν \textbf{υβριδική κρυπτογράφηση}: οι ασύμμετρες τεχνικές δημιουργούν ή προστατεύουν ένα τυχαίο συμμετρικό \emph{κλειδί συνεδρίας} και αυτό το κλειδί κρυπτογραφεί στη συνέχεια τον κύριο όγκο των δεδομένων. Κάθε σύνδεση TLS —και, όπως έδειξε το Κεφάλαιο 4, κάθε σύγχρονη επίθεση λυτρισμικού— ακολουθεί αυτό το πρότυπο.

\section{Κρυπτογραφικές συναρτήσεις κατακερματισμού}\label{sec:7.2}
Μια \textbf{κρυπτογραφική συνάρτηση κατακερματισμού} αντιστοιχίζει είσοδο οποιουδήποτε μήκους σε έξοδο σταθερού μήκους, τη \emph{σύνοψη} (digest). Η SHA-256, για παράδειγμα, παράγει πάντοτε 256 bit. Τρεις ιδιότητες καθιστούν μια συνάρτηση κατακερματισμού κατάλληλη για εφαρμογές ασφάλειας. \textbf{Αντίσταση προεικόνας}: δεδομένης μιας σύνοψης, είναι υπολογιστικά ανέφικτο να βρεθεί είσοδος που την παράγει. \textbf{Αντίσταση δεύτερης προεικόνας}: δεδομένης μιας εισόδου, είναι ανέφικτο να βρεθεί διαφορετική είσοδος με την ίδια σύνοψη. \textbf{Αντίσταση συγκρούσεων}: είναι ανέφικτο να βρεθούν \emph{οποιεσδήποτε} δύο είσοδοι με την ίδια σύνοψη. Οι MD5 και SHA-1 έχουν απολέσει την αντίσταση συγκρούσεων και δεν πρέπει πλέον να χρησιμοποιούνται σε υπογραφές· οι SHA-256, SHA-3 και BLAKE2 παραμένουν ασφαλείς.

Ο κατακερματισμός δεν είναι κρυπτογράφηση: δεν υπάρχει κλειδί ούτε τρόπος «αποκρυπτογράφησης» μιας σύνοψης. Οι συνόψεις χρησιμοποιούνται για την επαλήθευση της ακεραιότητας αρχείων, για την ταυτοποίηση αποδεικτικών στοιχείων στην εγκληματολογική ανάλυση (Κεφάλαιο 12) και ως δομικά στοιχεία των υπογραφών. Σε συνδυασμό με ένα μυστικό κλειδί σχηματίζουν έναν κώδικα \textbf{HMAC}, ο οποίος αποδεικνύει ότι ένα μήνυμα προέρχεται από κάποιον που γνωρίζει το κλειδί. Για την αποθήκευση κωδικών πρόσβασης οι συνήθεις γρήγορες συναρτήσεις είναι ακατάλληλες, επειδή οι επιτιθέμενοι μπορούν να δοκιμάζουν δισεκατομμύρια μαντεψιές ανά δευτερόλεπτο. Ειδικές \textbf{συναρτήσεις κατακερματισμού κωδικών}, όπως οι Argon2, scrypt ή bcrypt, είναι σκόπιμα αργές και χρησιμοποιούν ένα μοναδικό τυχαίο \emph{αλάτι} (salt) για κάθε κωδικό.

\section{Ψηφιακές υπογραφές και μη αποποίηση}\label{sec:7.3}
Μια \textbf{ψηφιακή υπογραφή} αντιστρέφει τους ρόλους του ζεύγους κλειδιών. Ο υπογράφων υπολογίζει τη σύνοψη του μηνύματος και τη μετασχηματίζει με το \emph{ιδιωτικό} του κλειδί· οποιοσδήποτε μπορεί να επαληθεύσει το αποτέλεσμα με το αντίστοιχο \emph{δημόσιο} κλειδί. Μια έγκυρη υπογραφή αποδεικνύει τρία πράγματα: ότι το μήνυμα δεν έχει αλλοιωθεί από τη στιγμή της υπογραφής (\textbf{ακεραιότητα}), ότι υπογράφηκε από τον κάτοχο του ιδιωτικού κλειδιού (\textbf{αυθεντικότητα}) και ότι ο υπογράφων δεν μπορεί εύκολα να ισχυριστεί αργότερα το αντίθετο (\textbf{μη αποποίηση}). Συνήθεις αλγόριθμοι είναι οι RSA-PSS, ECDSA και Ed25519.

Η μη αποποίηση, ωστόσο, είναι τόσο ισχυρή όσο η διαδικασία που την περιβάλλει. Αν το ιδιωτικό κλειδί ήταν αποθηκευμένο χωρίς προστασία σε κοινόχρηστο υπολογιστή, ο υπογράφων μπορεί εύλογα να υποστηρίξει ότι το χρησιμοποίησε κάποιος άλλος. Γι' αυτό νομικά πλαίσια όπως ο Κανονισμός \textbf{eIDAS} της ΕΕ διακρίνουν τις απλές ηλεκτρονικές υπογραφές από τις \emph{εγκεκριμένες} ηλεκτρονικές υπογραφές, οι οποίες απαιτούν εγκεκριμένο πιστοποιητικό και πιστοποιημένη διάταξη δημιουργίας υπογραφής και έχουν έννομο αποτέλεσμα ισοδύναμο με εκείνο της ιδιόχειρης υπογραφής.

\section{Υποδομή δημόσιου κλειδιού και ο κύκλος ζωής των πιστοποιητικών X.509}\label{sec:7.4}
Η κρυπτογραφία δημόσιου κλειδιού θέτει ένα νέο ερώτημα: πώς γνωρίζουμε ότι ένα δημόσιο κλειδί ανήκει πράγματι σε αυτόν που νομίζουμε; Ένας ενδιάμεσος επιτιθέμενος θα μπορούσε να το αντικαταστήσει με το δικό του. Μια \textbf{υποδομή δημόσιου κλειδιού} (ΥΔΚ ή PKI) απαντά σε αυτό το ερώτημα: μια αξιόπιστη \textbf{αρχή πιστοποίησης} (Certificate Authority, CA) υπογράφει ένα \textbf{πιστοποιητικό} που συνδέει ένα δημόσιο κλειδί με μια ταυτότητα, όπως ένα όνομα τομέα. Μια \textbf{αρχή καταχώρισης} (Registration Authority, RA) επαληθεύει την ταυτότητα του αιτούντος πριν η CA εκδώσει το πιστοποιητικό. Τα πιστοποιητικά ακολουθούν το πρότυπο \textbf{X.509} και περιέχουν, μεταξύ άλλων, το υποκείμενο, τον εκδότη, την περίοδο ισχύος, το δημόσιο κλειδί και επεκτάσεις όπως το \emph{Εναλλακτικό Όνομα Υποκειμένου} (Subject Alternative Name, SAN), που απαριθμεί τα ονόματα τομέα που καλύπτει το πιστοποιητικό.

Η εμπιστοσύνη οργανώνεται ως \textbf{αλυσίδα}: μια \emph{ριζική CA} εκτός σύνδεσης, της οποίας το πιστοποιητικό είναι προεγκατεστημένο στα λειτουργικά συστήματα και στους φυλλομετρητές, υπογράφει \emph{ενδιάμεσες CA}, οι οποίες με τη σειρά τους υπογράφουν τα πιστοποιητικά τελικών οντοτήτων. Η επαλήθευση ακολουθεί την αλυσίδα προς τα πάνω έως μια αξιόπιστη ρίζα. Τα πιστοποιητικά έχουν \textbf{κύκλο ζωής} —αίτηση (μέσω αιτήματος υπογραφής πιστοποιητικού, CSR), έκδοση, εγκατάσταση, ανανέωση και, σε περίπτωση διαρροής του κλειδιού, \textbf{ανάκληση}, η οποία γνωστοποιείται μέσω λιστών ανάκλησης (CRL) ή του πρωτοκόλλου άμεσης κατάστασης (OCSP). Επειδή ο έλεγχος ανάκλησης είναι στην πράξη αναξιόπιστος, ο κλάδος κινείται προς πιστοποιητικά μικρής διάρκειας που ανανεώνονται αυτόματα με το πρωτόκολλο ACME.

\section{Το TLS 1.3 και οι εφαρμογές της ΥΔΚ}\label{sec:7.5}
Το \textbf{Transport Layer Security (TLS)} προστατεύει την περιήγηση στον ιστό (HTTPS), τη μεταφορά ηλεκτρονικού ταχυδρομείου, τις διεπαφές προγραμματισμού (API) και πολλά άλλα πρωτόκολλα. Η έκδοση 1.3 (2018) απλοποίησε το πρωτόκολλο και αφαίρεσε παρωχημένες και ευάλωτες επιλογές. Η χειραψία απαιτεί μόνο έναν κύκλο επικοινωνίας. Ο πελάτης στέλνει ένα μήνυμα \emph{ClientHello} με τις σουίτες κρυπτογράφησης που υποστηρίζει και ένα εφήμερο μερίδιο κλειδιού Diffie–Hellman· ο διακομιστής απαντά με το δικό του μερίδιο, το πιστοποιητικό του και μια υπογραφή που αποδεικνύει ότι κατέχει το ιδιωτικό κλειδί. Στη συνέχεια οι δύο πλευρές παράγουν τα ίδια κλειδιά συνεδρίας. Επειδή τα κλειδιά Diffie–Hellman είναι εφήμερα, το TLS 1.3 παρέχει πάντοτε \textbf{εμπρόσθια μυστικότητα} (forward secrecy): η μεταγενέστερη υποκλοπή του ιδιωτικού κλειδιού του διακομιστή δεν επιτρέπει την αποκρυπτογράφηση κίνησης που είχε καταγραφεί στο παρελθόν.

Οι ίδιες αρχές της ΥΔΚ υποστηρίζουν και άλλες εφαρμογές. Το \textbf{S/MIME} υπογράφει και κρυπτογραφεί μηνύματα ηλεκτρονικού ταχυδρομείου· η \textbf{υπογραφή κώδικα} επιτρέπει στα λειτουργικά συστήματα να επαληθεύουν ότι ένα λογισμικό προέρχεται από γνωστό εκδότη και δεν έχει τροποποιηθεί —γι' αυτό τα κλεμμένα πιστοποιητικά που χρησιμοποίησε το Stuxnet ήταν τόσο επιζήμια· και το \textbf{αμοιβαίο TLS (mTLS)} απαιτεί πιστοποιητικό τόσο από τον διακομιστή όσο και από τον πελάτη, πρότυπο που χρησιμοποιείται ευρέως για τον έλεγχο ταυτότητας μεταξύ υπηρεσιών στις σύγχρονες αρχιτεκτονικές.

\section{Διαχείριση κλειδιών: εκεί όπου η κρυπτογραφία πραγματικά αποτυγχάνει}\label{sec:7.6}
Οι σύγχρονοι αλγόριθμοι σπάνια παραβιάζονται μαθηματικά· τα συστήματα αποτυγχάνουν επειδή τα κλειδιά παράγονται πρόχειρα, αποθηκεύονται απρόσεκτα ή δεν ανανεώνονται ποτέ. Τυπικές αστοχίες είναι τα ιδιωτικά κλειδιά που δημοσιεύονται κατά λάθος σε δημόσια αποθετήρια κώδικα, τα ενσωματωμένα διαπιστευτήρια στις εφαρμογές, οι προβλέψιμες γεννήτριες τυχαίων αριθμών και τα πιστοποιητικά που λήγουν χωρίς να το αντιληφθεί κανείς, προκαλώντας διακοπές λειτουργίας. Η \textbf{διαχείριση κλειδιών} καλύπτει ολόκληρο τον κύκλο ζωής: ασφαλή παραγωγή από κρυπτογραφικά ασφαλή πηγή τυχαιότητας, προστατευμένη αποθήκευση, ελεγχόμενη διανομή, τακτική εναλλαγή, ανάκληση και, τέλος, καταστροφή.

Τα κλειδιά υψηλής αξίας πρέπει να φυλάσσονται σε ειδικό υλικό. Μια \textbf{μονάδα ασφάλειας υλικού} (Hardware Security Module, HSM) ή μια υπηρεσία διαχείρισης κλειδιών στο νέφος εκτελεί τις κρυπτογραφικές λειτουργίες εσωτερικά, ώστε το ιδιωτικό κλειδί να μην εγκαταλείπει ποτέ τη συσκευή. Η πρόσβαση στα κλειδιά πρέπει να ακολουθεί τις αρχές του ελάχιστου προνομίου και του διαχωρισμού καθηκόντων, και κάθε χρήση πρέπει να καταγράφεται. Με ματιά στο μέλλον, οι οργανισμοί οφείλουν επίσης να τηρούν απογραφή του πού χρησιμοποιείται κάθε αλγόριθμος (\emph{κρυπτογραφική ευελιξία}), ώστε οι ευάλωτοι αλγόριθμοι να μπορούν να αντικατασταθούν —για παράδειγμα, καθώς υιοθετούνται οι μετακβαντικοί αλγόριθμοι που τυποποίησε το NIST.

\section*{Βασικοί όροι}
\begin{description}[style=nextline,leftmargin=1.2em]
  \item[Υβριδική κρυπτογράφηση] Χρήση ασύμμετρης κρυπτογραφίας για την προστασία ενός συμμετρικού κλειδιού συνεδρίας που κρυπτογραφεί τα δεδομένα.
  \item[Αντίσταση συγκρούσεων] Η αδυναμία εύρεσης δύο διαφορετικών εισόδων με την ίδια σύνοψη.
  \item[Ψηφιακή υπογραφή] Τιμή που υπολογίζεται με ιδιωτικό κλειδί και αποδεικνύει την ακεραιότητα και την προέλευση ενός μηνύματος.
  \item[Αρχή πιστοποίησης] Αξιόπιστη οντότητα που υπογράφει πιστοποιητικά τα οποία συνδέουν δημόσια κλειδιά με ταυτότητες.
  \item[Εμπρόσθια μυστικότητα] Η ιδιότητα κατά την οποία η διαρροή μακροπρόθεσμων κλειδιών δεν εκθέτει την κίνηση προηγούμενων συνεδριών.
  \item[HSM] Μονάδα ασφάλειας υλικού: ανθεκτική σε παραβίαση συσκευή που φυλάσσει κλειδιά και εκτελεί κρυπτογραφικές λειτουργίες.
\end{description}

\section*{Σύνοψη κεφαλαίου}
\begin{itemize}
  \item Οι συμμετρικοί αλγόριθμοι είναι γρήγοροι αλλά απαιτούν κοινό κλειδί· οι ασύμμετροι επιλύουν τη διανομή κλειδιών· τα υβριδικά σχήματα συνδυάζουν και τα δύο.
  \item Οι συναρτήσεις κατακερματισμού παρέχουν έλεγχο ακεραιότητας και όχι εμπιστευτικότητα· οι κωδικοί απαιτούν αργές συναρτήσεις κατακερματισμού με αλάτι.
  \item Οι ψηφιακές υπογραφές παρέχουν ακεραιότητα και αυθεντικότητα· η μη αποποίηση εξαρτάται επιπλέον από τη φύλαξη των κλειδιών και το νομικό πλαίσιο.
  \item Η ΥΔΚ συνδέει κλειδιά με ταυτότητες μέσω αλυσίδων πιστοποιητικών· η διαχείριση του κύκλου ζωής και η ανάκληση είναι απαραίτητες.
  \item Το TLS 1.3 προσφέρει γρήγορη χειραψία με εμπρόσθια μυστικότητα, όμως οι περισσότερες πραγματικές αποτυχίες οφείλονται σε κακή διαχείριση κλειδιών.
\end{itemize}

\section*{Ερωτήσεις ανασκόπησης}
\begin{enumerate}
  \item Γιατί τα πρακτικά συστήματα δεν κρυπτογραφούν μεγάλα αρχεία απευθείας με RSA;
  \item Ένας προγραμματιστής αποθηκεύει τους κωδικούς ως συνόψεις SHA-256 χωρίς αλάτι. Εξηγήστε τον κίνδυνο και προτείνετε καλύτερο σχεδιασμό.
  \item Περιγράψτε βήμα προς βήμα πώς ένας φυλλομετρητής επαληθεύει την αλυσίδα πιστοποιητικών ενός ιστοτόπου.
  \item Εξηγήστε την εμπρόσθια μυστικότητα και γιατί το TLS 1.3 κατάργησε τη στατική ανταλλαγή κλειδιών RSA.
\end{enumerate}

% ============================================================
% Chapter 8: Identity, Directory Services and Federated Single Sign-On
% Language: Greek — edit freely; regenerate from src/content if preferred.
% ============================================================
\chapter{Ταυτότητα, Υπηρεσίες Καταλόγου και Ομοσπονδιακή Ενιαία Σύνδεση}\label{ch:8}
\chaptermeta{IAM · Παράγοντες ταυτοποίησης και MFA · FIDO2 · RBAC και ABAC · LDAP και Active Directory · Kerberos · SAML, OAuth 2.0 και OpenID Connect}{Επίπεδο: Μεσαίο επίπεδο · Ταυτότητα}{Χρόνος μελέτης: 9–11 ώρες}

Καθώς οι οργανισμοί μεταφέρονται σε υπηρεσίες νέφους και στην απομακρυσμένη εργασία, η περίμετρος του δικτύου χάνει το νόημά της και \textbf{η ταυτότητα γίνεται η νέα περίμετρος}. Οι περισσότερες σύγχρονες εισβολές δεν σπάνε την κρυπτογράφηση ούτε εκμεταλλεύονται εξωτικές ευπάθειες· συνδέονται με κλεμμένα ή καταχρηστικά χρησιμοποιούμενα διαπιστευτήρια. Η κατανόηση του τρόπου με τον οποίο δημιουργούνται, επαληθεύονται, εξουσιοδοτούνται και ομοσπονδοποιούνται οι ταυτότητες είναι επομένως κεντρική για τη σύγχρονη άμυνα.

Το κεφάλαιο στηρίζεται στα κρυπτογραφικά θεμέλια του Κεφαλαίου 7. Ξεκινά με την αρχιτεκτονική της διαχείρισης ταυτότητας και πρόσβασης και με τους παράγοντες που χρησιμοποιούνται για την ταυτοποίηση, συμπεριλαμβανομένου του ανθεκτικού στο ψάρεμα FIDO2. Στη συνέχεια συγκρίνει τα μοντέλα ελέγχου πρόσβασης, εξηγεί τις υπηρεσίες καταλόγου και το πρωτόκολλο Kerberos στο οποίο στηρίζονται οι τομείς Windows —μαζί με τις επιθέσεις εναντίον του— και ολοκληρώνεται με τα πρωτόκολλα ομοσπονδίας SAML, OAuth 2.0 και OpenID Connect, που επιτρέπουν την ενιαία σύνδεση πέρα από τα όρια ενός οργανισμού.

\begin{outcomesbox}{Μαθησιακά αποτελέσματα}
Μετά τη μελέτη του κεφαλαίου θα πρέπει να μπορείτε:
\begin{itemize}
  \item Να διακρίνετε την αναγνώριση, την ταυτοποίηση (αυθεντικοποίηση), την εξουσιοδότηση και την καταγραφή.
  \item Να συγκρίνετε τους παράγοντες ταυτοποίησης και να εξηγείτε γιατί το FIDO2 αντιστέκεται στο ηλεκτρονικό ψάρεμα.
  \item Να επιλέγετε μεταξύ RBAC, ABAC και ελέγχου πρόσβασης βάσει κανόνων για ένα δεδομένο σενάριο.
  \item Να εξηγείτε τη ροή εισιτηρίων του Kerberos και τις αρχές πίσω από τις επιθέσεις Kerberoasting, pass-the-ticket και golden ticket.
  \item Να περιγράφετε τον ρόλο των SAML, OAuth 2.0 και OpenID Connect, καθώς και συνήθεις αδυναμίες της ομοσπονδίας ταυτοτήτων.
\end{itemize}
\end{outcomesbox}

\section{Αρχιτεκτονική IAM: αναγνώριση, ταυτοποίηση, εξουσιοδότηση και καταγραφή}\label{sec:8.1}
Η \textbf{διαχείριση ταυτότητας και πρόσβασης} (Identity and Access Management, IAM) οργανώνεται γύρω από τέσσερα διακριτά βήματα. Η \textbf{αναγνώριση} είναι η δήλωση μιας ταυτότητας, όπως η πληκτρολόγηση ενός ονόματος χρήστη. Η \textbf{ταυτοποίηση} ή \textbf{αυθεντικοποίηση} είναι η επαλήθευση αυτής της δήλωσης, για παράδειγμα με έλεγχο ενός κωδικού ή ενός κλειδιού ασφαλείας. Η \textbf{εξουσιοδότηση} καθορίζει τι επιτρέπεται να κάνει το υποκείμενο που ταυτοποιήθηκε. Η \textbf{καταγραφή} (ή λογιστική παρακολούθηση) καταγράφει τι έκανε πράγματι το υποκείμενο, υποστηρίζοντας τη λογοδοσία και τις έρευνες. Η σύγχυση αυτών των βημάτων οδηγεί σε σχεδιαστικά σφάλματα —για παράδειγμα, στην αντιμετώπιση ενός επιτυχώς ταυτοποιημένου χρήστη ως αυτομάτως εξουσιοδοτημένου για κάθε πόρο.

Η IAM καλύπτει επίσης τον \textbf{κύκλο ζωής της ταυτότητας}: οι λογαριασμοί πρέπει να δημιουργούνται όταν κάποιος εντάσσεται στον οργανισμό, να προσαρμόζονται όταν αλλάζει ρόλο και να απενεργοποιούνται άμεσα όταν αποχωρεί. Οι ορφανοί λογαριασμοί πρώην εργαζομένων και η σταδιακή συσσώρευση δικαιωμάτων κατά τη διάρκεια μιας επαγγελματικής πορείας —η λεγόμενη \emph{διολίσθηση προνομίων} (privilege creep)— συγκαταλέγονται στα πιο συνηθισμένα ευρήματα των ελέγχων ασφάλειας. Το καθιερωμένο αντίμετρο είναι οι περιοδικές \textbf{επισκοπήσεις πρόσβασης}, κατά τις οποίες οι προϊστάμενοι επιβεβαιώνουν ότι κάθε δικαίωμα εξακολουθεί να είναι αναγκαίο.

\section{Παράγοντες ταυτοποίησης, MFA και FIDO2 ανθεκτικό στο ψάρεμα}\label{sec:8.2}
Οι παράγοντες ταυτοποίησης ομαδοποιούνται παραδοσιακά σε τρεις κατηγορίες: \textbf{κάτι που γνωρίζεις} (κωδικός ή PIN), \textbf{κάτι που έχεις} (κινητό τηλέφωνο, έξυπνη κάρτα ή κλειδί ασφαλείας) και \textbf{κάτι που είσαι} (βιομετρικό χαρακτηριστικό, όπως το δακτυλικό αποτύπωμα). Ο \textbf{έλεγχος ταυτότητας πολλαπλών παραγόντων} (Multi-Factor Authentication, MFA) συνδυάζει παράγοντες από \emph{διαφορετικές} κατηγορίες, ώστε η υποκλοπή ενός παράγοντα να μην αρκεί. Ωστόσο, δεν είναι όλες οι μορφές MFA εξίσου ισχυρές. Οι κωδικοί μίας χρήσης μέσω SMS μπορούν να υποκλαπούν με αντικατάσταση κάρτας SIM (SIM swapping)· οι χρονικά μεταβαλλόμενοι κωδικοί εφαρμογών ταυτοποίησης (TOTP, σύμφωνα με τα πρότυπα OATH) και οι ειδοποιήσεις push μπορούν ακόμη να αναμεταδοθούν από έναν διακομιστή ψαρέματος σε πραγματικό χρόνο.

Το \textbf{FIDO2/WebAuthn} αλλάζει ριζικά αυτή την εικόνα. Κατά την εγγραφή, ο αυθεντικοποιητής —ένα κλειδί ασφαλείας ή το ασφαλές υλικό ενός κινητού ή φορητού υπολογιστή— δημιουργεί ένα ζεύγος κλειδιών ειδικά για τον συγκεκριμένο ιστότοπο. Κατά τη σύνδεση, υπογράφει μια πρόκληση που περιλαμβάνει την προέλευση (origin) του ιστοτόπου. Ένας ψεύτικος ιστότοπος με διαφορετικό όνομα τομέα λαμβάνει επομένως μια υπογραφή άχρηστη για τον πραγματικό ιστότοπο, και δεν υπάρχει κοινό μυστικό προς υποκλοπή. Γι' αυτό το FIDO2 χαρακτηρίζεται \textbf{ανθεκτικό στο ηλεκτρονικό ψάρεμα} και γι' αυτό τα \emph{passkeys} που βασίζονται σε αυτό αντικαθιστούν τους κωδικούς σε πολλές υπηρεσίες.

\section{Μοντέλα ελέγχου πρόσβασης: RBAC, ABAC και έλεγχος βάσει κανόνων}\label{sec:8.3}
Αφού ταυτοποιηθεί ένα υποκείμενο, το σύστημα πρέπει να αποφασίσει τι επιτρέπεται να κάνει. Στον \textbf{έλεγχο πρόσβασης βάσει ρόλων} (Role-Based Access Control, RBAC), τα δικαιώματα ανατίθενται σε ρόλους (για παράδειγμα \emph{νοσηλευτής}, \emph{λογιστής}, \emph{διαχειριστής βάσης δεδομένων}) και οι χρήστες λαμβάνουν ρόλους. Το RBAC είναι εύκολο στην κατανόηση και στον έλεγχο, όμως οι μεγάλοι οργανισμοί μπορεί να αντιμετωπίσουν \emph{έκρηξη ρόλων} όταν κάθε εξαίρεση απαιτεί νέο ρόλο. Ο \textbf{έλεγχος πρόσβασης βάσει χαρακτηριστικών} (Attribute-Based Access Control, ABAC) αξιολογεί πολιτικές επί χαρακτηριστικών του υποκειμένου, του πόρου, της ενέργειας και του περιβάλλοντος —για παράδειγμα, «οι ιατροί μπορούν να διαβάζουν τους φακέλους ασθενών της δικής τους κλινικής, κατά τη διάρκεια της βάρδιάς τους, από διαχειριζόμενη συσκευή». Το ABAC είναι πιο εκφραστικό, αλλά δυσκολότερο στην ανάλυση και στον έλεγχο.

Ο \textbf{έλεγχος πρόσβασης βάσει κανόνων} εφαρμόζει γενικούς κανόνες που δεν εξαρτώνται από τον μεμονωμένο χρήστη, όπως οι κανόνες ενός τείχους προστασίας ή ο κανόνας «καμία σύνδεση μεταξύ μεσονυκτίου και 5 π.μ.». Στην πράξη οι οργανισμοί συνδυάζουν τα μοντέλα: οι ρόλοι παρέχουν μια σταθερή βάση, τα χαρακτηριστικά εξειδικεύουν τις αποφάσεις ανάλογα με το πλαίσιο και οι γενικοί κανόνες επιβάλλουν περιορισμούς για ολόκληρο τον οργανισμό.

\section{Υπηρεσίες καταλόγου: X.500, LDAP και Active Directory}\label{sec:8.4}
Μια \textbf{υπηρεσία καταλόγου} είναι μια εξειδικευμένη, ιεραρχική βάση δεδομένων που αποθηκεύει πληροφορίες για χρήστες, ομάδες, υπολογιστές και άλλους πόρους και είναι βελτιστοποιημένη για συχνές αναγνώσεις. Το εννοιολογικό της μοντέλο προέρχεται από τα πρότυπα \textbf{X.500}, ενώ το \textbf{Lightweight Directory Access Protocol (LDAP)} είναι ο συνήθης τρόπος υποβολής ερωτημάτων και τροποποίησης. Κάθε εγγραφή προσδιορίζεται από ένα \emph{διακεκριμένο όνομα} (distinguished name), όπως \texttt{CN=Maria Papadopoulou,OU=Finance,DC=example,DC=com}.

Οι \textbf{Active Directory Domain Services (AD DS)} της Microsoft συνδυάζουν έναν κατάλογο LDAP με ταυτοποίηση Kerberos, DNS και πολιτικές ομάδας (Group Policy). Οργανώνουν τα αντικείμενα σε \emph{τομείς}, \emph{δέντρα} και \emph{δάση}, ενώ οι \emph{ελεγκτές τομέα} αναπαράγουν τον κατάλογο μεταξύ τους. Επειδή το AD ελέγχει την ταυτοποίηση σχεδόν κάθε συστήματος σε μια επιχείρηση Windows, η παραβίαση ενός λογαριασμού διαχειριστή τομέα σημαίνει συνήθως παραβίαση ολόκληρου του οργανισμού. Γι' αυτό η ασφάλεια του AD οργανώνεται γύρω από την \textbf{κλιμακωτή διαχείριση} (tiered administration): οι λογαριασμοί υψηλών προνομίων χρησιμοποιούνται μόνο σε αποκλειστικά, θωρακισμένα συστήματα και ποτέ σε συνηθισμένους σταθμούς εργασίας, όπου τα διαπιστευτήριά τους θα μπορούσαν να υποκλαπούν.

\section{Το Kerberos και οι επιθέσεις εναντίον του}\label{sec:8.5}
Το \textbf{Kerberos} επιτρέπει στους χρήστες να ταυτοποιούνται μία φορά και στη συνέχεια να έχουν πρόσβαση σε πολλές υπηρεσίες χωρίς να αποστέλλουν τον κωδικό τους μέσω του δικτύου. Ένα αξιόπιστο \textbf{Κέντρο Διανομής Κλειδιών} (Key Distribution Center, KDC) —στο AD, κάθε ελεγκτής τομέα— περιλαμβάνει δύο λογικές υπηρεσίες. Κατά τη σύνδεση, η \textbf{Υπηρεσία Ταυτοποίησης} (Authentication Service, AS) επαληθεύει τον χρήστη και εκδίδει ένα \textbf{Εισιτήριο Χορήγησης Εισιτηρίων} (Ticket-Granting Ticket, TGT), κρυπτογραφημένο με το κλειδί του ειδικού λογαριασμού \texttt{krbtgt}. Όταν ο χρήστης θέλει να προσπελάσει μια υπηρεσία, ο πελάτης παρουσιάζει το TGT στην \textbf{Υπηρεσία Χορήγησης Εισιτηρίων} (Ticket-Granting Service, TGS), η οποία εκδίδει ένα \emph{εισιτήριο υπηρεσίας} κρυπτογραφημένο με το κλειδί του λογαριασμού που εκτελεί την υπηρεσία. Η υπηρεσία αποκρυπτογραφεί το εισιτήριο και παραχωρεί την πρόσβαση. Τα εισιτήρια έχουν περιορισμένη διάρκεια ισχύος και το πρωτόκολλο προϋποθέτει συγχρονισμένα ρολόγια.

Ο σχεδιασμός αυτός επιτρέπει αρκετές γνωστές επιθέσεις. Στο \textbf{Kerberoasting}, οποιοσδήποτε χρήστης του τομέα ζητά εισιτήρια υπηρεσίας για λογαριασμούς που διαθέτουν \emph{όνομα κύριας υπηρεσίας} (Service Principal Name, SPN) και τα «σπάει» εκτός σύνδεσης· οι αδύναμοι κωδικοί λογαριασμών υπηρεσιών υποχωρούν γρήγορα. Στο \textbf{pass-the-ticket}, κλεμμένα εισιτήρια επαναχρησιμοποιούνται από άλλο μηχάνημα. Ένα \textbf{golden ticket} είναι ένα πλαστό TGT που δημιουργείται με το κλεμμένο κλειδί του \texttt{krbtgt} και παρέχει απεριόριστη πρόσβαση στον τομέα. Τα αντίμετρα περιλαμβάνουν μεγάλους τυχαίους κωδικούς ή διαχειριζόμενους λογαριασμούς υπηρεσιών ομάδας (gMSA), χρήση αποκλειστικά κρυπτογράφησης AES, προστασία των ελεγκτών τομέα ως στοιχείων της ανώτατης βαθμίδας (Tier 0), διπλή επαναφορά του κωδικού του \texttt{krbtgt} μετά από παραβίαση και παρακολούθηση ασυνήθιστων αιτημάτων εισιτηρίων.

\section{Ομοσπονδιακή ταυτότητα και ενιαία σύνδεση: SAML, OAuth 2.0 και OpenID Connect}\label{sec:8.6}
Η \textbf{ομοσπονδία ταυτοτήτων} (federation) επεκτείνει την ενιαία σύνδεση πέρα από έναν οργανισμό: ο χρήστης ταυτοποιείται στον δικό του \textbf{πάροχο ταυτότητας} (Identity Provider, IdP) και πολλοί \textbf{πάροχοι υπηρεσιών} αποδέχονται αυτή την ταυτοποίηση. Το \textbf{SAML 2.0} ανταλλάσσει ψηφιακά υπογεγραμμένους \emph{ισχυρισμούς} (assertions) σε XML και χρησιμοποιείται ευρέως σε επιχειρησιακές εφαρμογές ιστού. Το \textbf{OAuth 2.0} δεν είναι, αυστηρά μιλώντας, πρωτόκολλο ταυτοποίησης αλλά πλαίσιο \emph{εξουσιοδότησης}: επιτρέπει σε έναν χρήστη να παραχωρήσει σε μια εφαρμογή περιορισμένη πρόσβαση στους πόρους του (για παράδειγμα «ανάγνωση του ημερολογίου μου») μέσω έκδοσης ενός \emph{διακριτικού πρόσβασης} (access token), χωρίς να κοινοποιήσει τον κωδικό του. Το \textbf{OpenID Connect (OIDC)} προσθέτει ένα επίπεδο ταυτότητας πάνω από το OAuth 2.0, μέσω ενός υπογεγραμμένου \emph{διακριτικού ταυτότητας} (ID token) που δηλώνει ποιος είναι ο χρήστης.

Η ομοσπονδία συγκεντρώνει την εμπιστοσύνη και, επομένως, τον κίνδυνο. Αν ένας επιτιθέμενος υποκλέψει το κλειδί υπογραφής διακριτικών του παρόχου ταυτότητας, μπορεί να πλαστογραφήσει πρόσβαση σε κάθε συνδεδεμένη υπηρεσία —όπως συνέβη με την τεχνική \emph{Golden SAML} που χρησιμοποιήθηκε κατά την εκστρατεία SolarWinds. Άλλες συνήθεις αδυναμίες είναι οι εφαρμογές που δεν επικυρώνουν την υπογραφή, τον αποδέκτη ή τον χρόνο λήξης των διακριτικών· τα κλεμμένα διακριτικά συνεδρίας που επαναχρησιμοποιούνται από άλλη συσκευή· και το \emph{ψάρεμα συγκατάθεσης} (consent phishing), κατά το οποίο οι χρήστες εξαπατώνται ώστε να παραχωρήσουν ευρεία δικαιώματα OAuth σε κακόβουλη εφαρμογή. Οι άμυνες περιλαμβάνουν την προστασία των κλειδιών υπογραφής σε HSM, την αυστηρή επικύρωση των διακριτικών, τη μικρή διάρκεια ισχύος διακριτικών δεσμευμένων σε συσκευή και την έγκριση από διαχειριστή για συγκαταθέσεις OAuth υψηλού κινδύνου.

\section*{Βασικοί όροι}
\begin{description}[style=nextline,leftmargin=1.2em]
  \item[Ταυτοποίηση έναντι εξουσιοδότησης] Επαλήθευση του ποιος είναι ένα υποκείμενο έναντι απόφασης για το τι επιτρέπεται να κάνει.
  \item[FIDO2 / passkey] Ταυτοποίηση δημόσιου κλειδιού δεσμευμένη στην προέλευση, ανθεκτική στο ψάρεμα και χωρίς κοινό μυστικό.
  \item[ABAC] Έλεγχος πρόσβασης με βάση χαρακτηριστικά του υποκειμένου, του πόρου, της ενέργειας και του περιβάλλοντος.
  \item[TGT] Εισιτήριο χορήγησης εισιτηρίων του Kerberos, για τη λήψη εισιτηρίων υπηρεσίας χωρίς εκ νέου εισαγωγή διαπιστευτηρίων.
  \item[Kerberoasting] Αίτηση εισιτηρίων υπηρεσίας και «σπάσιμο» των κωδικών λογαριασμών υπηρεσιών εκτός σύνδεσης.
  \item[OIDC] OpenID Connect: επίπεδο ταυτότητας πάνω από το OAuth 2.0 που εκδίδει υπογεγραμμένα διακριτικά ταυτότητας.
\end{description}

\section*{Σύνοψη κεφαλαίου}
\begin{itemize}
  \item Η ταυτότητα είναι η σύγχρονη περίμετρος· η IAM διακρίνει την αναγνώριση, την ταυτοποίηση, την εξουσιοδότηση και την καταγραφή και διαχειρίζεται ολόκληρο τον κύκλο ζωής της ταυτότητας.
  \item Το MFA συνδυάζει παράγοντες από διαφορετικές κατηγορίες· το FIDO2 αντιστέκεται στο ψάρεμα επειδή οι υπογραφές δεσμεύονται στην προέλευση του ιστοτόπου.
  \item Το RBAC προσφέρει απλότητα, το ABAC ακρίβεια με βάση το πλαίσιο και ο έλεγχος βάσει κανόνων επιβάλλει γενικούς περιορισμούς.
  \item Το Active Directory και το Kerberos είναι στόχοι υψηλής αξίας· η κλιμακωτή διαχείριση και η σωστή διαχείριση των λογαριασμών υπηρεσιών είναι απαραίτητες.
  \item Τα πρωτόκολλα ομοσπονδίας επιτρέπουν την ενιαία σύνδεση, αλλά συγκεντρώνουν τον κίνδυνο στα κλειδιά υπογραφής και στην επικύρωση των διακριτικών.
\end{itemize}

\section*{Ερωτήσεις ανασκόπησης}
\begin{enumerate}
  \item Γιατί ένας κωδικός μίας χρήσης μέσω SMS είναι πιο αδύναμος από ένα κλειδί ασφαλείας FIDO2 απέναντι σε έναν διακομιστή ψαρέματος πραγματικού χρόνου;
  \item Σχεδιάστε μια πολιτική πρόσβασης για ένα σύστημα ιατρικών φακέλων με RBAC και, στη συνέχεια, εξειδικεύστε την με χαρακτηριστικά ABAC.
  \item Εξηγήστε γιατί το Kerberoasting είναι εφικτό για οποιονδήποτε ταυτοποιημένο χρήστη του τομέα και πώς το περιορίζουν οι λογαριασμοί gMSA.
  \item Γιατί το OAuth 2.0 από μόνο του δεν επαρκεί για την ταυτοποίηση χρηστών και τι προσθέτει το OIDC;
\end{enumerate}

% ============================================================
% Chapter 9: Secure Software Design and Engineering
% Language: Greek — edit freely; regenerate from src/content if preferred.
% ============================================================
\chapter{Ασφαλής Σχεδίαση και Ανάπτυξη Λογισμικού}\label{ch:9}
\chaptermeta{Shift-left και DevSecOps · Έγχυση κώδικα · Cross-site scripting · CSRF και clickjacking · Ασφάλεια API, SSRF και σπασμένος έλεγχος πρόσβασης · Μυστικά και εφοδιαστική αλυσίδα λογισμικού}{Επίπεδο: Μεσαίο επίπεδο · Ανάπτυξη λογισμικού}{Χρόνος μελέτης: 8–10 ώρες}

Οι περισσότερες ευπάθειες δεν δημιουργούνται από τους επιτιθέμενους· γράφονται, ακούσια, από τους προγραμματιστές. Ένας μόνο έλεγχος επικύρωσης που λείπει μπορεί να εκθέσει εκατομμύρια εγγραφές, ενώ η διόρθωση ενός σφάλματος μετά την κυκλοφορία κοστίζει πολύ περισσότερο από την πρόληψή του κατά τον σχεδιασμό. Η ασφαλής μηχανική λογισμικού στοχεύει, επομένως, να καταστήσει την ασφάλεια φυσιολογική ιδιότητα της καθημερινής αναπτυξιακής εργασίας και όχι έναν τελικό έλεγχο πριν από την κυκλοφορία.

Το κεφάλαιο παρουσιάζει πρώτα τη φιλοσοφία \emph{shift-left} και τις αρχές μηχανικής που καθοδηγούν τον ασφαλή σχεδιασμό. Στη συνέχεια εξηγεί, για καθεμιά από τις σημαντικότερες κατηγορίες ευπαθειών ιστού —έγχυση, cross-site scripting, πλαστογράφηση αιτημάτων μεταξύ ιστοτόπων, πλαστογράφηση αιτημάτων από την πλευρά του διακομιστή και σπασμένο έλεγχο πρόσβασης— γιατί προκύπτει το σφάλμα και ποια άμυνα αντιμετωπίζει τη βασική του αιτία. Ολοκληρώνεται με τη διαχείριση μυστικών, την καταγραφή και την προστασία της εφοδιαστικής αλυσίδας λογισμικού. Οι κατηγορίες ευπαθειών που εξετάζονται αντιστοιχούν σε μεγάλο βαθμό στο OWASP Top 10.

\begin{outcomesbox}{Μαθησιακά αποτελέσματα}
Μετά τη μελέτη του κεφαλαίου θα πρέπει να μπορείτε:
\begin{itemize}
  \item Να εξηγείτε την ασφάλεια shift-left και πώς ενσωματώνονται οι δραστηριότητες ασφάλειας σε μια ροή DevSecOps.
  \item Να εξηγείτε τη βασική αιτία της έγχυσης και να εφαρμόζετε παραμετροποιημένα ερωτήματα.
  \item Να διακρίνετε το αποθηκευμένο, το ανακλώμενο και το βασισμένο στο DOM XSS και να εφαρμόζετε κωδικοποίηση εξόδου ανάλογα με το πλαίσιο και CSP.
  \item Να περιγράφετε τις επιθέσεις CSRF, clickjacking, SSRF και τον σπασμένο έλεγχο πρόσβασης, μαζί με τις άμυνές τους.
  \item Να περιγράφετε πώς η ανάλυση σύνθεσης λογισμικού, τα SBOM και το SLSA προστατεύουν την εφοδιαστική αλυσίδα λογισμικού.
\end{itemize}
\end{outcomesbox}

\section{Ασφάλεια shift-left, DevSecOps και βασικές αρχές μηχανικής}\label{sec:9.1}
Ο όρος \textbf{shift-left} σημαίνει τη μετακίνηση των δραστηριοτήτων ασφάλειας νωρίτερα —«προς τα αριστερά»— στο χρονοδιάγραμμα της ανάπτυξης. Οι απαιτήσεις ασφάλειας ορίζονται μαζί με τις λειτουργικές· η \textbf{μοντελοποίηση απειλών} (για παράδειγμα με τη μέθοδο STRIDE: πλαστοπροσωπία, αλλοίωση, αποποίηση, αποκάλυψη πληροφοριών, άρνηση υπηρεσίας, κλιμάκωση προνομίων) εντοπίζει τους κινδύνους όσο ο σχεδιασμός μπορεί ακόμη να αλλάξει με χαμηλό κόστος· και αυτοματοποιημένοι έλεγχοι εκτελούνται σε κάθε αλλαγή του κώδικα. Το \textbf{DevSecOps} ενσωματώνει αυτούς τους ελέγχους στη ροή συνεχούς ενοποίησης και παράδοσης, ώστε μια υποβολή κώδικα που εισάγει μια βιβλιοθήκη με γνωστή ευπάθεια ή έναν ενσωματωμένο κωδικό να αποτυγχάνει στη μεταγλώττιση, όπως ακριβώς και ένας αποτυχημένος έλεγχος μονάδας.

Ορισμένες αρχές μηχανικής καθοδηγούν τη συγγραφή ασφαλούς κώδικα. \textbf{Επικύρωση της εισόδου} στην πλευρά του διακομιστή, με βάση λίστα όσων αναμένονται, επειδή κάθε είσοδος από έξω από το όριο εμπιστοσύνης μπορεί να είναι εχθρική. \textbf{Κωδικοποίηση της εξόδου} σύμφωνα με το πλαίσιο στο οποίο χρησιμοποιείται. \textbf{Ελαχιστοποίηση της επιφάνειας επίθεσης} με την αφαίρεση αχρησιμοποίητων λειτουργιών και σημείων πρόσβασης. \textbf{Ασφαλής αποτυχία}, ώστε τα σφάλματα να απορρίπτουν την πρόσβαση και να μη διαρρέουν εσωτερικές λεπτομέρειες. Και προτίμηση των δοκιμασμένων \textbf{λειτουργιών ασφάλειας των πλαισίων ανάπτυξης} έναντι αυτοσχέδιων μηχανισμών, στο πνεύμα της αρχής της οικονομίας μηχανισμού από το Κεφάλαιο 1.

\section{Έγχυση SQL και NoSQL: παραμετροποίηση έναντι συνένωσης συμβολοσειρών}\label{sec:9.2}
Η \textbf{έγχυση} (injection) συμβαίνει όταν μη αξιόπιστα δεδομένα ερμηνεύονται ως μέρος μιας εντολής. Η κλασική περίπτωση είναι ένα ερώτημα SQL που κατασκευάζεται με συνένωση συμβολοσειρών, όπως \texttt{''SELECT * FROM users WHERE name = ''' + input + '''''}. Αν ένας επιτιθέμενος εισαγάγει \texttt{' OR '1'='1}, η συνθήκη γίνεται πάντα αληθής και το ερώτημα επιστρέφει όλους τους χρήστες· πιο σύνθετα φορτία μπορούν να διαβάσουν άλλους πίνακες, να τροποποιήσουν δεδομένα ή, σε ορισμένες διαμορφώσεις, να εκτελέσουν εντολές του λειτουργικού συστήματος. Η βασική αιτία είναι η \emph{ανάμειξη κώδικα και δεδομένων} σε μία συμβολοσειρά.

Η οριστική άμυνα είναι το \textbf{παραμετροποιημένο ερώτημα} (prepared statement), στο οποίο η δομή του SQL αποστέλλεται στη βάση δεδομένων χωριστά από τις τιμές: \texttt{SELECT * FROM users WHERE name = ?}. Η βάση δεδομένων αντιμετωπίζει τότε την είσοδο αυστηρά ως δεδομένα, όποιους χαρακτήρες κι αν περιέχει. Οι βιβλιοθήκες αντιστοίχισης αντικειμένων–σχέσεων (ORM) παρέχουν την ίδια προστασία όταν χρησιμοποιούνται σωστά. Η επικύρωση εισόδου και οι λογαριασμοί βάσης δεδομένων με ελάχιστα προνόμια είναι πολύτιμα πρόσθετα επίπεδα, όμως η χειροκίνητη διαφυγή ειδικών χαρακτήρων είναι επιρρεπής σε σφάλματα και δεν πρέπει να αποτελεί τη βασική άμυνα. Οι βάσεις NoSQL δεν είναι απρόσβλητες: η αποδοχή αντικειμένων JSON όπως \texttt{\{''\$ne'': null\}} σε ένα ερώτημα MongoDB μπορεί να παρακάμψει την ταυτοποίηση με ακριβώς τον ίδιο τρόπο.

\section{Cross-site scripting: αποθηκευμένο, ανακλώμενο και βασισμένο στο DOM}\label{sec:9.3}
Το \textbf{cross-site scripting (XSS)} είναι έγχυση στον φυλλομετρητή ιστού. Ο επιτιθέμενος κάνει έναν ιστότοπο να παραδώσει κώδικα JavaScript που εκτελείται στον φυλλομετρητή του θύματος με όλα τα προνόμια του ιστοτόπου —διαβάζοντας δεδομένα συνεδρίας, εκτελώντας ενέργειες εκ μέρους του χρήστη ή αλλοιώνοντας τη σελίδα. Στο \textbf{αποθηκευμένο XSS}, το κακόβουλο σενάριο αποθηκεύεται στον διακομιστή (για παράδειγμα σε ένα σχόλιο) και σερβίρεται σε κάθε επισκέπτη. Στο \textbf{ανακλώμενο XSS}, το σενάριο είναι μέρος ενός ειδικά κατασκευασμένου συνδέσμου και επιστρέφεται αμέσως στην απόκριση. Στο \textbf{XSS βασισμένο στο DOM}, η ευπάθεια βρίσκεται εξ ολοκλήρου σε κώδικα της πλευράς του πελάτη που γράφει μη αξιόπιστα δεδομένα στη σελίδα, για παράδειγμα μέσω της ιδιότητας \texttt{innerHTML}.

Η κύρια άμυνα είναι η \textbf{κωδικοποίηση εξόδου ανάλογα με το πλαίσιο}: δεδομένα που εισάγονται σε HTML, σε χαρακτηριστικό HTML, σε JavaScript ή σε URL πρέπει να κωδικοποιούνται σύμφωνα με τους κανόνες του αντίστοιχου πλαισίου. Τα σύγχρονα πλαίσια ανάπτυξης, όπως τα React και Angular, κωδικοποιούν εξ ορισμού, και οι προγραμματιστές πρέπει να αποφεύγουν παρακάμψεις όπως το \texttt{dangerouslySetInnerHTML}, εκτός αν το περιεχόμενο έχει εξυγιανθεί με μια βιβλιοθήκη όπως η DOMPurify. Μια \textbf{Πολιτική Ασφάλειας Περιεχομένου} (Content Security Policy, CSP) προσθέτει άμυνα σε βάθος, δίνοντας εντολή στον φυλλομετρητή να εκτελεί μόνο σενάρια από εγκεκριμένες πηγές, ενώ η σημαία \texttt{HttpOnly} στα cookies εμποδίζει τα σενάρια να διαβάζουν τα cookies συνεδρίας.

\section{Πλαστογράφηση αιτημάτων μεταξύ ιστοτόπων και clickjacking}\label{sec:9.4}
Στην \textbf{πλαστογράφηση αιτημάτων μεταξύ ιστοτόπων} (Cross-Site Request Forgery, CSRF), ένας κακόβουλος ιστότοπος κάνει τον φυλλομετρητή του θύματος να στείλει αίτημα σε άλλον ιστότοπο στον οποίο το θύμα είναι συνδεδεμένο. Επειδή οι φυλλομετρητές επισυνάπτουν αυτόματα τα cookies, ο ιστότοπος-στόχος βλέπει ένα ταυτοποιημένο αίτημα —για παράδειγμα αλλαγής της διεύθυνσης email του λογαριασμού ή μεταφοράς χρημάτων— που ο χρήστης δεν σκόπευε ποτέ να στείλει. Οι άμυνες περιλαμβάνουν τα \textbf{διακριτικά anti-CSRF} (απρόβλεπτες τιμές ενσωματωμένες στις φόρμες, τις οποίες επαληθεύει ο διακομιστής), το χαρακτηριστικό \textbf{SameSite} των cookies, το οποίο εμποδίζει την αποστολή τους με τα περισσότερα αιτήματα μεταξύ ιστοτόπων, και την επανεπαλήθευση ταυτότητας για ευαίσθητες ενέργειες.

Το \textbf{clickjacking} ενσωματώνει αόρατα τον ιστότοπο-στόχο μέσα σε ένα πλαίσιο (frame) στη σελίδα του επιτιθέμενου και εξαπατά τον χρήστη ώστε να πατήσει κουμπιά που δεν βλέπει. Αποτρέπεται με την απαγόρευση της ενσωμάτωσης σε πλαίσια μέσω της οδηγίας \texttt{frame-ancestors} της CSP ή της παλαιότερης κεφαλίδας \texttt{X-Frame-Options}.

\section{Ασφάλεια API: σπασμένος έλεγχος πρόσβασης, SSRF και XXE}\label{sec:9.5}
Ο \textbf{σπασμένος έλεγχος πρόσβασης} είναι η πιο διαδεδομένη κατηγορία του OWASP Top 10. Η συνηθέστερη μορφή του στις διεπαφές API είναι η \emph{μη ασφαλής άμεση αναφορά σε αντικείμενο} (IDOR, γνωστή και ως BOLA): ένα αίτημα όπως \texttt{GET /api/invoices/1043} επιστρέφει το τιμολόγιο ακόμη κι όταν ανήκει σε άλλον πελάτη, επειδή ο διακομιστής ελέγχει \emph{ότι} ο χρήστης είναι συνδεδεμένος αλλά όχι \emph{αν} του ανήκει το αντικείμενο. Κάθε αίτημα πρέπει να εξουσιοδοτείται στον διακομιστή, για κάθε αντικείμενο, σύμφωνα με την αρχή της πλήρους διαμεσολάβησης. Οι διεπαφές API πρέπει επίσης να επιβάλλουν ταυτοποίηση με σωστά επικυρωμένα διακριτικά, να εφαρμόζουν όρια ρυθμού αιτημάτων και να επιστρέφουν μόνο τα πεδία που χρειάζεται ο πελάτης.

Η \textbf{πλαστογράφηση αιτημάτων από την πλευρά του διακομιστή} (Server-Side Request Forgery, SSRF) εξαπατά έναν διακομιστή ώστε να πραγματοποιήσει αιτήματα για λογαριασμό του επιτιθέμενου —για παράδειγμα, μια λειτουργία «λήψη εικόνας από URL» που κατευθύνεται σε εσωτερικές υπηρεσίες ή στο σημείο μεταδεδομένων του νέφους \texttt{169.254.169.254}, το οποίο μπορεί να αποκαλύψει προσωρινά διαπιστευτήρια. Οι άμυνες περιλαμβάνουν λίστες επιτρεπόμενων προορισμών, αποκλεισμό εσωτερικών περιοχών διευθύνσεων και χρήση της θωρακισμένης υπηρεσίας μεταδεδομένων (IMDSv2). Οι επιθέσεις \textbf{εξωτερικής οντότητας XML} (XML External Entity, XXE) εκμεταλλεύονται αναλυτές XML που επιλύουν εξωτερικές αναφορές, επιτρέποντας την ανάγνωση αρχείων από τον διακομιστή· η θεραπεία είναι η απενεργοποίηση των εξωτερικών οντοτήτων και των ορισμών τύπου εγγράφου στη ρύθμιση του αναλυτή.

\section{Μυστικά, καταγραφή και εφοδιαστική αλυσίδα λογισμικού}\label{sec:9.6}
Τα \textbf{μυστικά} —κωδικοί, κλειδιά API, ιδιωτικά κλειδιά— δεν πρέπει ποτέ να γράφονται σε πηγαίο κώδικα ή σε αρχεία ρυθμίσεων που καταχωρίζονται στο σύστημα ελέγχου εκδόσεων, επειδή τα αποθετήρια αντιγράφονται, κοινοποιούνται και ενίοτε δημοσιοποιούνται. Φυλάσσονται σε ειδικό διαχειριστή μυστικών (όπως το HashiCorp Vault ή ένα θησαυροφυλάκιο κλειδιών στο νέφος), εισάγονται κατά την εκτέλεση και εναλλάσσονται τακτικά· εργαλεία σάρωσης μυστικών στη ροή ανάπτυξης εντοπίζουν τις ακούσιες καταχωρίσεις. Η \textbf{καταγραφή ασφάλειας} πρέπει να καταγράφει συμβάντα ταυτοποίησης, αποτυχίες ελέγχου πρόσβασης και διαχειριστικές ενέργειες με επαρκές πλαίσιο για τη διερεύνηση, αλλά δεν πρέπει ποτέ να καταγράφει κωδικούς, διακριτικά ή περιττά προσωπικά δεδομένα.

Οι σύγχρονες εφαρμογές αποτελούνται κυρίως από κώδικα τρίτων: ένα τυπικό έργο μπορεί να εξαρτάται από εκατοντάδες πακέτα ανοιχτού κώδικα. Η \textbf{ανάλυση σύνθεσης λογισμικού} (Software Composition Analysis, SCA) εντοπίζει αυτές τις εξαρτήσεις και επισημαίνει όσες έχουν γνωστές ευπάθειες, όπως στο περιστατικό Log4Shell του 2021. Ένας \textbf{κατάλογος υλικών λογισμικού} (Software Bill of Materials, SBOM), σε μορφές όπως SPDX ή CycloneDX, απαριθμεί κάθε στοιχείο, ώστε τα επηρεαζόμενα συστήματα να εντοπίζονται μέσα σε λίγα λεπτά όταν ανακοινώνεται μια νέα ευπάθεια. Το πλαίσιο \textbf{SLSA} (\emph{Supply-chain Levels for Software Artifacts}) ορίζει αυξανόμενα επίπεδα διασφάλισης για τη διαδικασία κατασκευής λογισμικού —όπως υπογεγραμμένες, αναπαραγώγιμες εκδόσεις με επαληθεύσιμη προέλευση— ώστε επιθέσεις όπως η παραβίαση της διαδικασίας κατασκευής της SolarWinds να γίνονται ανιχνεύσιμες.

\section*{Βασικοί όροι}
\begin{description}[style=nextline,leftmargin=1.2em]
  \item[Shift-left] Μετακίνηση των δραστηριοτήτων ασφάλειας σε πρωιμότερα στάδια του κύκλου ανάπτυξης.
  \item[Μοντελοποίηση απειλών] Συστηματικός εντοπισμός απειλών κατά ενός σχεδιασμού, π.χ. με τη μέθοδο STRIDE.
  \item[Παραμετροποιημένο ερώτημα] Ερώτημα με σταθερή δομή, του οποίου οι τιμές διαβιβάζονται χωριστά ως δεδομένα.
  \item[XSS] Cross-site scripting: έγχυση σεναρίου που εκτελείται στον φυλλομετρητή του θύματος υπό την προέλευση του ιστοτόπου.
  \item[IDOR / BOLA] Πρόσβαση σε αντικείμενα μέσω αναγνωριστικού χωρίς έλεγχο αν ο αιτών είναι εξουσιοδοτημένος γι' αυτά.
  \item[SBOM] Κατάλογος υλικών λογισμικού: απογραφή όλων των στοιχείων ενός λογισμικού.
\end{description}

\section*{Σύνοψη κεφαλαίου}
\begin{itemize}
  \item Η ασφάλεια κοστίζει λιγότερο όταν ενσωματώνεται νωρίς: η μοντελοποίηση απειλών και οι αυτοματοποιημένοι έλεγχοι στη ροή ανάπτυξης αποτελούν τον πυρήνα του DevSecOps.
  \item Η έγχυση προκύπτει από την ανάμειξη κώδικα και δεδομένων· τα παραμετροποιημένα ερωτήματα εξαλείφουν τη βασική αιτία.
  \item Το XSS αποτρέπεται με κωδικοποίηση εξόδου ανάλογα με το πλαίσιο, ασφαλείς προεπιλογές των πλαισίων ανάπτυξης και CSP.
  \item Τα CSRF, clickjacking, SSRF και XXE έχουν συγκεκριμένες, καθιερωμένες άμυνες· ο σπασμένος έλεγχος πρόσβασης απαιτεί εξουσιοδότηση κάθε αντικειμένου στην πλευρά του διακομιστή.
  \item Οι διαχειριστές μυστικών, η προσεκτική καταγραφή, η SCA, τα SBOM και το SLSA προστατεύουν την εφαρμογή και την εφοδιαστική της αλυσίδα.
\end{itemize}

\section*{Ερωτήσεις ανασκόπησης}
\begin{enumerate}
  \item Εξηγήστε γιατί η διαφυγή των εισαγωγικών είναι ασθενέστερη άμυνα κατά της έγχυσης SQL σε σύγκριση με τα παραμετροποιημένα ερωτήματα.
  \item Κατατάξτε ως αποθηκευμένο, ανακλώμενο ή βασισμένο στο DOM XSS την εξής περίπτωση: μια σελίδα αναζήτησης που εμφανίζει το ερώτημα από το URL χρησιμοποιώντας innerHTML.
  \item Γιατί το χαρακτηριστικό SameSite των cookies περιορίζει το CSRF αλλά όχι το XSS;
  \item Όταν ανακοινώθηκε το Log4Shell, πώς θα βοηθούσε ένα SBOM έναν οργανισμό να ανταποκριθεί;
\end{enumerate}

% ============================================================
% Chapter 10: Security Testing: SAST, DAST, Vulnerability Assessment and Penetration Testing
% Language: Greek — edit freely; regenerate from src/content if preferred.
% ============================================================
\chapter{Έλεγχος Ασφάλειας: SAST, DAST, Αξιολόγηση Ευπαθειών και Δοκιμές Διείσδυσης}\label{ch:10}
\chaptermeta{Στατική ανάλυση · Δυναμική ανάλυση · Κύκλος διαχείρισης ευπαθειών · Μεθοδολογία και δεοντολογία δοκιμών διείσδυσης · CVSS και επαγγελματική αναφορά}{Επίπεδο: Μεσαίο–Προχωρημένο επίπεδο · Διασφάλιση}{Χρόνος μελέτης: 7–9 ώρες}

Το Κεφάλαιο 9 περιέγραψε πώς αναπτύσσεται λογισμικό με ασφάλεια· το παρόν κεφάλαιο εξετάζει πώς μπορούμε να \emph{επαληθεύσουμε} ότι είναι ασφαλές. Καμία μεμονωμένη τεχνική δεν εντοπίζει κάθε αδυναμία. Η στατική ανάλυση διαβάζει τον κώδικα χωρίς να τον εκτελεί, η δυναμική ανάλυση εξετάζει μια εφαρμογή σε λειτουργία από έξω, η αξιολόγηση ευπαθειών ελέγχει ολόκληρο ένα περιβάλλον για γνωστές αδυναμίες και η δοκιμή διείσδυσης αποδεικνύει τι θα μπορούσε πράγματι να επιτύχει ένας ικανός αντίπαλος. Καθεμιά απαντά σε διαφορετικό ερώτημα, με διαφορετικό κόστος και σε διαφορετικό στάδιο.

Επειδή ορισμένες από αυτές τις δραστηριότητες χρησιμοποιούν πραγματικά επιθετικές τεχνικές, το κεφάλαιο δίνει ίση βαρύτητα στη μέθοδο και στη δεοντολογία. Εξηγεί τα νομικά θεμέλια του εξουσιοδοτημένου ελέγχου, τα όρια της δραστηριότητας μετά την εκμετάλλευση και τον χειρισμό ευαίσθητων αποδεικτικών στοιχείων, και δείχνει πώς πρέπει να βαθμολογούνται και να αναφέρονται τα ευρήματα, ώστε να οδηγούν σε αποκατάσταση και όχι σε ένα έγγραφο που δεν διαβάζει κανείς.

\begin{outcomesbox}{Μαθησιακά αποτελέσματα}
Μετά τη μελέτη του κεφαλαίου θα πρέπει να μπορείτε:
\begin{itemize}
  \item Να συγκρίνετε τα SAST, DAST, την αξιολόγηση ευπαθειών και τη δοκιμή διείσδυσης ως προς το εύρος, τα πλεονεκτήματα και τους περιορισμούς τους.
  \item Να περιγράφετε τον κύκλο διαχείρισης ευπαθειών και να ιεραρχείτε τα ευρήματα με βάση το CVSS και την πιθανότητα εκμετάλλευσης.
  \item Να εξηγείτε τα μοντέλα, τις φάσεις και τις αναγνωρισμένες μεθοδολογίες των δοκιμών διείσδυσης.
  \item Να διατυπώνετε τις νομικές και δεοντολογικές απαιτήσεις του εξουσιοδοτημένου ελέγχου, συμπεριλαμβανομένων των κανόνων εμπλοκής.
  \item Να συντάσσετε ένα εύρημα που συνδυάζει τεκμηρίωση, αξιολόγηση κινδύνου και εφαρμόσιμη πρόταση αποκατάστασης.
\end{itemize}
\end{outcomesbox}

\section{Στατικός έλεγχος ασφάλειας εφαρμογών (SAST)}\label{sec:10.1}
Ο \textbf{στατικός έλεγχος ασφάλειας εφαρμογών} (Static Application Security Testing, SAST) αναλύει πηγαίο κώδικα, ενδιάμεσο κώδικα ή εκτελέσιμα \emph{χωρίς να τα εκτελεί}. Είναι τεχνική «λευκού κουτιού»: το εργαλείο έχει πρόσβαση σε ολόκληρο τον κώδικα. Τα σύγχρονα εργαλεία SAST κατασκευάζουν ένα μοντέλο του προγράμματος και εκτελούν \textbf{ανάλυση μόλυνσης} (taint analysis), παρακολουθώντας τα δεδομένα από τις \emph{πηγές} (όπως οι παράμετροι HTTP) έως τα επικίνδυνα \emph{σημεία κατάληξης} (όπως η εκτέλεση SQL ή η έξοδος HTML). Αν μολυσμένα δεδομένα φτάσουν σε σημείο κατάληξης χωρίς να περάσουν από αναγνωρισμένο \emph{μηχανισμό εξυγίανσης}, το εργαλείο αναφέρει πιθανή ευπάθεια μαζί με την ακριβή γραμμή του κώδικα.

Το μεγάλο πλεονέκτημα του SAST είναι ο χρόνος εφαρμογής του: μπορεί να εκτελείται στον επεξεργαστή κώδικα του προγραμματιστή και σε κάθε υποβολή, πολύ πριν εγκατασταθεί η εφαρμογή. Η κύρια αδυναμία του είναι η ακρίβεια. Επειδή δεν μπορεί να παρατηρήσει τη συμπεριφορά κατά την εκτέλεση, παράγει \textbf{ψευδώς θετικά} αποτελέσματα (αναφερόμενα ζητήματα που δεν είναι εκμεταλλεύσιμα) και παραλείπει σφάλματα που εξαρτώνται από τη διαμόρφωση ή από την επιχειρησιακή λογική. Τα επιτυχημένα προγράμματα προσαρμόζουν επομένως τους κανόνες στα πλαίσια ανάπτυξης που χρησιμοποιούν, αξιολογούν τα αποτελέσματα και αντιμετωπίζουν το SAST ως ένα από πολλά επίπεδα.

\section{Δυναμικός έλεγχος ασφάλειας εφαρμογών (DAST)}\label{sec:10.2}
Ο \textbf{δυναμικός έλεγχος ασφάλειας εφαρμογών} (Dynamic Application Security Testing, DAST) εξετάζει μια εφαρμογή \emph{σε λειτουργία} από έξω, όπως θα έκανε ένας επιτιθέμενος, χωρίς πρόσβαση στον πηγαίο κώδικα. Ένας σαρωτής DAST, όπως το OWASP ZAP, πρώτα \emph{ανιχνεύει} την εφαρμογή για να ανακαλύψει σελίδες, παραμέτρους και σημεία πρόσβασης API και στη συνέχεια αποστέλλει ειδικά διαμορφωμένες εισόδους και αναλύει τις αποκρίσεις για ενδείξεις ευπάθειας —ένα μήνυμα σφάλματος SQL, ένα ανακλώμενο σενάριο, μια κεφαλίδα ασφάλειας που λείπει. Επειδή παρατηρεί την πραγματική συμπεριφορά, ένα εύρημα DAST είναι συνήθως πράγματι εκμεταλλεύσιμο, και μπορεί να εντοπίσει προβλήματα διαμόρφωσης και διακομιστή που το SAST δεν βλέπει.

Το DAST έχει επίσης όρια. Μπορεί να ελέγξει μόνο όσα καταφέρνει να προσεγγίσει, οπότε σελίδες πίσω από σύνθετη ταυτοποίηση ή ροές πολλών βημάτων ενδέχεται να παραλειφθούν, και δεν μπορεί να υποδείξει την υπεύθυνη γραμμή κώδικα. Εκτελείται σε όψιμο στάδιο του κύκλου ζωής, σε εγκατεστημένο περιβάλλον δοκιμών. Ο \textbf{διαδραστικός έλεγχος} (Interactive Application Security Testing, IAST) συνδυάζει τις δύο προσεγγίσεις, εξοπλίζοντας την εφαρμογή με αισθητήρες ενώ εκτελούνται οι δυναμικές δοκιμές, ώστε κάθε παρατηρούμενο ζήτημα να συνδέεται με τη θέση του στον κώδικα.

\section{Ο κύκλος διαχείρισης ευπαθειών}\label{sec:10.3}
Η \textbf{αξιολόγηση ευπαθειών} εντοπίζει συστηματικά γνωστές αδυναμίες —ενημερώσεις που λείπουν, μη ασφαλείς διαμορφώσεις, εργοστασιακά διαπιστευτήρια— σε δίκτυα, υπολογιστές και εφαρμογές, συνήθως με αυτοματοποιημένους σαρωτές όπως τα Nessus, OpenVAS ή τα σενάρια του Nmap. Είναι εκτενής αλλά ρηχή: αναφέρει τι \emph{ενδέχεται} να είναι εκμεταλλεύσιμο χωρίς να επιχειρεί εκμετάλλευση. Η αξιολόγηση αποτελεί ένα στάδιο του συνεχούς \textbf{κύκλου διαχείρισης ευπαθειών}: \emph{ανακάλυψη περιουσιακών στοιχείων} (δεν μπορείς να προστατεύσεις ό,τι δεν γνωρίζεις ότι έχεις), \emph{σάρωση}, \emph{ανάλυση και ιεράρχηση}, \emph{αποκατάσταση}, \emph{επαλήθευση} ότι η διόρθωση λειτουργεί και \emph{αναφορά} των τάσεων.

Η ιεράρχηση είναι το σημείο όπου δυσκολεύονται τα περισσότερα προγράμματα, επειδή οι σαρωτές αναφέρουν συνήθως χιλιάδες ευρήματα. Η σοβαρότητα από μόνη της δεν αρκεί. Μια τεκμηριωμένη προσέγγιση συνδυάζει την τεχνική σοβαρότητα (CVSS), την πιθανότητα εκμετάλλευσης —για παράδειγμα τη βαθμολογία \textbf{EPSS} ή την καταχώριση στον κατάλογο \emph{Known Exploited Vulnerabilities} της CISA— και την επιχειρησιακή σημασία και έκθεση του επηρεαζόμενου στοιχείου. Ένας διακομιστής εκτεθειμένος στο διαδίκτυο με ευπάθεια που εκμεταλλεύονται ενεργά οι επιτιθέμενοι χρειάζεται προσοχή σήμερα, ακόμη κι αν ένα «κρίσιμο» εύρημα σε ένα απομονωμένο μηχάνημα δοκιμών περιμένει.

\section{Δοκιμές διείσδυσης: μοντέλα, φάσεις, δεοντολογία και πλαίσια}\label{sec:10.4}
Μια \textbf{δοκιμή διείσδυσης} είναι μια εξουσιοδοτημένη, προσομοιωμένη επίθεση που επιχειρεί να εκμεταλλευτεί ευπάθειες, ώστε να αποδείξει την πραγματική τους επίπτωση. Οι δοκιμές κατατάσσονται με βάση τη γνώση που δίνεται στον ελεγκτή: \textbf{μαύρου κουτιού} (χωρίς προηγούμενη πληροφόρηση, όπως ένας εξωτερικός επιτιθέμενος), \textbf{γκρίζου κουτιού} (ορισμένες πληροφορίες, όπως ένας λογαριασμός χρήστη) και \textbf{λευκού κουτιού} (πλήρης τεκμηρίωση και πηγαίος κώδικας, που προσφέρει την πληρέστερη κάλυψη στον διαθέσιμο χρόνο). Μια τυπική ανάθεση εξελίσσεται σε φάσεις: \emph{προετοιμασία} και καθορισμός πεδίου, \emph{αναγνώριση}, \emph{σάρωση και απαρίθμηση}, \emph{εκμετάλλευση}, \emph{δραστηριότητα μετά την εκμετάλλευση} και \emph{αναφορά}. Αναγνωρισμένες μεθοδολογίες, όπως οι PTES, OWASP Web Security Testing Guide, OSSTMM και NIST SP 800-115, δομούν αυτές τις φάσεις και καθιστούν τα αποτελέσματα επαναλήψιμα.

Η νομιμότητα στηρίζεται αποκλειστικά στην \textbf{εξουσιοδότηση}. Πριν ξεκινήσει οποιαδήποτε δοκιμή, ο πελάτης πρέπει να υπογράψει σύμβαση και \textbf{κανόνες εμπλοκής} (rules of engagement) που ορίζουν το πεδίο (ποια συστήματα, διευθύνσεις και εφαρμογές), τις επιτρεπόμενες τεχνικές, τα χρονικά παράθυρα, τα πρόσωπα επικοινωνίας έκτακτης ανάγκης και τον τρόπο χειρισμού ευαίσθητων δεδομένων. Ο έλεγχος ενός συστήματος εκτός του συμφωνημένου πεδίου —ακόμη κι ενός που φαίνεται προφανώς σχετικό— αποτελεί μη εξουσιοδοτημένη πρόσβαση. Οι επαγγελματίες ελεγκτές αποφεύγουν επίσης τις καταστροφικές ενέργειες και σταματούν αμέσως αν εντοπίσουν ενδείξεις πραγματικής, εν εξελίξει παραβίασης.

\section{Όρια της φάσης μετά την εκμετάλλευση και χειρισμός αποδεικτικών στοιχείων}\label{sec:10.5}
Η φάση \textbf{μετά την εκμετάλλευση} (post-exploitation) αποδεικνύει την επιχειρησιακή επίπτωση μιας παραβίασης: ποια δεδομένα θα μπορούσαν να προσπελαστούν, αν ήταν δυνατή η κλιμάκωση προνομίων και πόσο μακριά θα μπορούσε να κινηθεί ο ελεγκτής μέσα στο δίκτυο. Είναι επίσης η φάση με τον μεγαλύτερο δεοντολογικό κίνδυνο. Ο κατευθυντήριος κανόνας είναι η \textbf{ελάχιστη αναγκαία ενέργεια}: αποδεικνύεις την πρόσβαση, δεν την εκμεταλλεύεσαι. Ένας ελεγκτής που φτάνει σε μια βάση δεδομένων πελατών καταγράφει ένα στιγμιότυπο της δομής του πίνακα και τον αριθμό των εγγραφών, όχι τις ίδιες τις εγγραφές. Κάθε μηχανισμός μόνιμης παρουσίας, λογαριασμός ή αρχείο που δημιουργείται κατά τη δοκιμή πρέπει να τεκμηριώνεται και να αφαιρείται στη συνέχεια.

Τα αποδεικτικά στοιχεία που συλλέγονται κατά τη δοκιμή —στιγμιότυπα, υποκλαπέντα διαπιστευτήρια, αποσπάσματα ευαίσθητων δεδομένων— είναι και τα ίδια ευαίσθητα. Πρέπει να αποθηκεύονται κρυπτογραφημένα, να κοινοποιούνται μόνο σε εξουσιοδοτημένους παραλήπτες, να διατηρούνται μόνο για όσο διάστημα απαιτεί η σύμβαση και στη συνέχεια να καταστρέφονται με ασφάλεια. Όταν εμπλέκονται προσωπικά δεδομένα, η επεξεργασία πρέπει επιπλέον να συμμορφώνεται με τη νομοθεσία προστασίας δεδομένων, όπως ο ΓΚΠΔ (Κεφάλαιο 13).

\section{Βαθμολόγηση και αναφορά: CVSS, τεκμηρίωση και αποκατάσταση}\label{sec:10.6}
Το \textbf{Common Vulnerability Scoring System (CVSS)} εκφράζει την τεχνική σοβαρότητα μιας ευπάθειας με βαθμολογία από 0 έως 10. Οι \emph{βασικές μετρικές} του περιγράφουν πώς μπορεί να γίνει εκμετάλλευση της ευπάθειας (διάνυσμα επίθεσης, πολυπλοκότητα, απαιτούμενα προνόμια, αλληλεπίδραση χρήστη) και ποια είναι η επίπτωσή της στην εμπιστευτικότητα, την ακεραιότητα και τη διαθεσιμότητα. Η βαθμολογία συνοδεύεται από μια \emph{συμβολοσειρά διανύσματος}, όπως \texttt{CVSS:3.1/AV:N/AC:L/PR:N/UI:N/S:U/C:H/I:H/A:H} (βαθμολογία 9,8), η οποία καθιστά διαφανή τη συλλογιστική. Το CVSS μετρά τη σοβαρότητα και όχι τον κίνδυνο· το πλαίσιο του οργανισμού πρέπει πάντοτε να λαμβάνεται υπόψη.

Μια επαγγελματική αναφορά απευθύνεται σε δύο κοινά. Η \textbf{σύνοψη για τη διοίκηση} εξηγεί, σε μη τεχνική γλώσσα, τον συνολικό κίνδυνο και τις σημαντικότερες συστάσεις. Κάθε \textbf{τεχνικό εύρημα} περιλαμβάνει σαφή τίτλο, τα επηρεαζόμενα στοιχεία, περιγραφή, τεκμηρίωση \emph{αναπαραγωγής} βήμα προς βήμα, τη σοβαρότητα με το διάνυσμα CVSS, την επιχειρησιακή επίπτωση και μια συγκεκριμένη, εφαρμόσιμη \textbf{πρόταση αποκατάστασης} —όχι «βελτιώστε την ασφάλεια», αλλά «αντικαταστήστε τη συνένωση συμβολοσειρών στη γραμμή 42 του \texttt{search.php} με παραμετροποιημένο ερώτημα και εφαρμόστε έναν κανόνα WAF ως προσωρινό μέτρο». Ένας επανέλεγχος επαληθεύει στη συνέχεια ότι οι διορθώσεις είναι αποτελεσματικές.

\begin{table}[htbp]
  \centering\small
  \caption{Σύγκριση προσεγγίσεων ελέγχου}
  \begin{tabularx}{\linewidth}{>{\raggedright\arraybackslash}X >{\raggedright\arraybackslash}X >{\raggedright\arraybackslash}X >{\raggedright\arraybackslash}X >{\raggedright\arraybackslash}X}
    \toprule
    \textbf{Προσέγγιση} & \textbf{Πρόσβαση} & \textbf{Πότε} & \textbf{Πλεονέκτημα} & \textbf{Περιορισμός} \\
    \midrule
    SAST & Πηγαίος κώδικας (λευκό κουτί) & Κατά την ανάπτυξη & Πρώιμο, υποδεικνύει την ακριβή γραμμή & Ψευδώς θετικά, καμία εικόνα εκτέλεσης \\
    DAST & Εφαρμογή σε λειτουργία (μαύρο κουτί) & Περιβάλλον δοκιμών & Επιβεβαιώνει την πραγματική συμπεριφορά & Περιορισμένη κάλυψη, όψιμο \\
    Αξιολόγηση ευπαθειών & Δίκτυο, υπολογιστές, εφαρμογές & Συνεχώς & Ευρεία κάλυψη & Καμία απόδειξη εκμεταλλευσιμότητας \\
    Δοκιμή διείσδυσης & Συμφωνημένο πεδίο & Περιοδικά / μετά από σημαντική αλλαγή & Αποδεικνύει την πραγματική επίπτωση & Χρονικά περιορισμένη, δαπανηρή, στιγμιαία εικόνα \\
    \bottomrule
  \end{tabularx}
\end{table}
\section*{Βασικοί όροι}
\begin{description}[style=nextline,leftmargin=1.2em]
  \item[Ανάλυση μόλυνσης] Παρακολούθηση μη αξιόπιστων δεδομένων από τις πηγές έως τα επικίνδυνα σημεία κατάληξης στον κώδικα.
  \item[Ψευδώς θετικό] Αναφερόμενο ζήτημα που στην πραγματικότητα δεν είναι εκμεταλλεύσιμο.
  \item[EPSS] Σύστημα πρόβλεψης εκμετάλλευσης: η εκτιμώμενη πιθανότητα να γίνει εκμετάλλευση μιας ευπάθειας.
  \item[Κανόνες εμπλοκής] Υπογεγραμμένο έγγραφο που ορίζει το πεδίο, τις μεθόδους, τον χρόνο και τα πρόσωπα επικοινωνίας μιας δοκιμής.
  \item[Μετά την εκμετάλλευση] Δραστηριότητες μετά την αρχική πρόσβαση που αποδεικνύουν την επίπτωση, με όριο την ελάχιστη αναγκαία ενέργεια.
  \item[Διάνυσμα CVSS] Συμβολοσειρά που καταγράφει κάθε μετρική που χρησιμοποιήθηκε για τον υπολογισμό μιας βαθμολογίας CVSS.
\end{description}

\section*{Σύνοψη κεφαλαίου}
\begin{itemize}
  \item Τα SAST, DAST, η αξιολόγηση ευπαθειών και η δοκιμή διείσδυσης απαντούν σε διαφορετικά ερωτήματα και αλληλοσυμπληρώνονται.
  \item Η διαχείριση ευπαθειών είναι συνεχής κύκλος· η ιεράρχηση πρέπει να συνδυάζει σοβαρότητα, πιθανότητα εκμετάλλευσης και πλαίσιο του στοιχείου.
  \item Οι δοκιμές διείσδυσης ακολουθούν αναγνωρισμένες φάσεις και μεθοδολογίες και είναι νόμιμες μόνο εντός υπογεγραμμένου πεδίου και κανόνων εμπλοκής.
  \item Η φάση μετά την εκμετάλλευση πρέπει να αποδεικνύει την πρόσβαση με την ελάχιστη αναγκαία ενέργεια, ενώ τα αποδεικτικά στοιχεία της δοκιμής προστατεύονται όπως τα δεδομένα παραγωγής.
  \item Οι καλές αναφορές συνδυάζουν διαφανή βαθμολόγηση CVSS με αναπαραγώγιμη τεκμηρίωση και συγκεκριμένη αποκατάσταση.
\end{itemize}

\section*{Ερωτήσεις ανασκόπησης}
\begin{enumerate}
  \item Γιατί το SAST μπορεί να αναφέρει μια ευπάθεια που το DAST δεν μπορεί να επιβεβαιώσει, και αντιστρόφως;
  \item Δύο ευρήματα: CVSS 9,1 σε έναν απομονωμένο εργαστηριακό διακομιστή και CVSS 7,5 σε έναν διακομιστή εκτεθειμένο στο διαδίκτυο που περιλαμβάνεται στον κατάλογο CISA KEV. Ποιο διορθώνετε πρώτο και γιατί;
  \item Κατά τη διάρκεια μιας δοκιμής ανακαλύπτετε ότι ένα σύστημα λίγο έξω από το πεδίο σας είναι ευάλωτο. Τι πρέπει να κάνετε;
  \item Αναδιατυπώστε τη σύσταση «Διορθώστε το πρόβλημα XSS» ώστε να γίνει εφαρμόσιμη πρόταση αποκατάστασης.
\end{enumerate}

\part{Λειτουργίες, Διερεύνηση και Διακυβέρνηση}
% ============================================================
% Chapter 11: Detection, Threat Hunting, Incident Response and Lab Operations
% Language: Greek — edit freely; regenerate from src/content if preferred.
% ============================================================
\chapter{Ανίχνευση, Προληπτική Αναζήτηση Απειλών, Απόκριση σε Περιστατικά και Εργαστηριακές Λειτουργίες}\label{ch:11}
\chaptermeta{IDS/IPS · Υπογραφές και ανίχνευση ανωμαλιών · Τεχνικές αποφυγής · Προληπτική αναζήτηση απειλών · YARA και Sigma · Απομονωμένα εργαστήρια · Απόκριση σε περιστατικά κατά NIST · Λειτουργίες SOC}{Επίπεδο: Μεσαίο–Προχωρημένο επίπεδο · Λειτουργίες}{Χρόνος μελέτης: 10–12 ώρες}

Η πρόληψη αργά ή γρήγορα αποτυγχάνει. Το Κεφάλαιο 1 εισήγαγε την άμυνα σε βάθος ακριβώς επειδή κάθε προληπτικό μέτρο έχει κενά, και τα προηγούμενα κεφάλαια έδειξαν πώς τα εντοπίζουν οι ικανοί αντίπαλοι. Το παρόν κεφάλαιο αφορά, επομένως, όσα ακολουθούν: πώς ένας οργανισμός \emph{ανιχνεύει} κακόβουλη δραστηριότητα, \emph{αναζητά} προληπτικά απειλές που διέφυγαν από τους συναγερμούς του και \emph{αποκρίνεται} με πειθαρχημένο τρόπο που περιορίζει τη ζημιά και διαφυλάσσει τα αποδεικτικά στοιχεία.

Ξεκινάμε με τα συστήματα ανίχνευσης και πρόληψης εισβολών, τα δύο κύρια παραδείγματα επιθεώρησης και τις τεχνικές με τις οποίες οι επιτιθέμενοι τα αποφεύγουν. Στη συνέχεια στρεφόμαστε στην προληπτική αναζήτηση απειλών, στις γλώσσες κανόνων YARA και Sigma και στον σχεδιασμό απομονωμένων εργαστηρίων, όπου το κακόβουλο λογισμικό μπορεί να μελετηθεί με ασφάλεια. Το κεφάλαιο ολοκληρώνεται με τον κύκλο απόκρισης σε περιστατικά όπως τον ορίζει το NIST και με την οργάνωση ενός Κέντρου Επιχειρήσεων Ασφάλειας (SOC).

\begin{outcomesbox}{Μαθησιακά αποτελέσματα}
Μετά τη μελέτη του κεφαλαίου θα πρέπει να μπορείτε:
\begin{itemize}
  \item Να συγκρίνετε τα δικτυακά και τα τοπικά IDS/IPS, καθώς και την επιθεώρηση βάσει υπογραφών και βάσει ανωμαλιών.
  \item Να εξηγείτε συνήθεις τεχνικές αποφυγής των IDS και τον ρόλο της κανονικοποίησης της κίνησης.
  \item Να σχεδιάζετε μια προληπτική αναζήτηση απειλών βασισμένη σε υποθέσεις, αξιοποιώντας τηλεμετρία τερματικών και μνήμης.
  \item Να συντάσσετε βασικούς κανόνες YARA και Sigma και να διαχειρίζεστε τις ανιχνεύσεις ως κώδικα.
  \item Να εφαρμόζετε τον κύκλο απόκρισης σε περιστατικά του NIST και να περιγράφετε τις βαθμίδες ενός SOC και τη ροή ενός SIEM.
\end{itemize}
\end{outcomesbox}

\section{Αρχιτεκτονικές IDS και IPS και παραδείγματα επιθεώρησης}\label{sec:11.1}
Ένα \textbf{σύστημα ανίχνευσης εισβολών} (Intrusion Detection System, IDS) παρακολουθεί τη δραστηριότητα και εκδίδει ειδοποιήσεις· ένα \textbf{σύστημα πρόληψης εισβολών} (Intrusion Prevention System, IPS) βρίσκεται στη διαδρομή της κίνησης και μπορεί να αποκλείει αυτόματα την κακόβουλη δραστηριότητα. Ένα \textbf{δικτυακό IDS} (NIDS), όπως τα Suricata, Snort ή Zeek, αναλύει αντίγραφα της δικτυακής κίνησης από θύρα κατοπτρισμού ενός μεταγωγέα ή από δικτυακό διακλαδωτή (tap) και βλέπει την επικοινωνία μεταξύ πολλών υπολογιστών. Ένα \textbf{τοπικό IDS} (HIDS), όπως τα Wazuh ή OSSEC, εκτελείται σε μεμονωμένα συστήματα και παρατηρεί αλλαγές αρχείων, αρχεία καταγραφής και διεργασίες —συμπεριλαμβανομένης δραστηριότητας μέσα σε κρυπτογραφημένες συνδέσεις που ένας δικτυακός αισθητήρας δεν μπορεί να διαβάσει. Οι δύο οπτικές είναι συμπληρωματικές.

Η λογική της ανίχνευσης ακολουθεί δύο παραδείγματα. Η ανίχνευση \textbf{βάσει υπογραφών} αναζητά γνωστά μοτίβα, όπως μια συγκεκριμένη συμβολοσειρά εκμετάλλευσης· είναι ακριβής και εύκολα εξηγήσιμη, αλλά τυφλή απέναντι σε νέες επιθέσεις. Η ανίχνευση \textbf{βάσει ανωμαλιών ή συμπεριφοράς} μοντελοποιεί το φυσιολογικό και επισημαίνει αποκλίσεις, όπως έναν σταθμό εργασίας που ξαφνικά μεταφέρει gigabytes τη νύχτα· μπορεί να εντοπίσει άγνωστες επιθέσεις, αλλά παράγει περισσότερα ψευδώς θετικά αποτελέσματα και απαιτεί προσεκτικά καθορισμένη συμπεριφορά αναφοράς. Τα ώριμα προγράμματα συνδυάζουν και τα δύο και τα βελτιστοποιούν συνεχώς, επειδή μια ειδοποίηση που οι αναλυτές έχουν μάθει να αγνοούν είναι χειρότερη από την απουσία ειδοποίησης.

\section{Τεχνικές αποφυγής και αυτοματοποιημένη αντιμετώπιση}\label{sec:11.2}
Οι επιτιθέμενοι προσπαθούν να κάνουν το IDS να «βλέπει» κάτι διαφορετικό από αυτό που λαμβάνει πράγματι ο στόχος. Ο \textbf{κατακερματισμός} διασπά ένα κακόβουλο φορτίο σε πολλά μικρά πακέτα ή σε επικαλυπτόμενα τμήματα· η \textbf{κωδικοποίηση} και η συσκότιση μεταμφιέζουν γνωστές συμβολοσειρές· η χειραγώγηση του χρονισμού ή των τιμών \emph{TTL} μπορεί να κάνει τον αισθητήρα και τον προορισμό να ανασυνθέτουν μια ροή με διαφορετικό τρόπο· και, όλο και συχνότερα, οι επιθέσεις απλώς κρύβονται μέσα σε \textbf{κρυπτογραφημένη} κίνηση TLS. Τα αντίμετρα περιλαμβάνουν την \textbf{κανονικοποίηση της κίνησης}, η οποία ανασυνθέτει τις ροές με τον ίδιο τρόπο όπως ο τελικός υπολογιστής, την επιθεώρηση TLS σε ελεγχόμενα σημεία όπου αυτό είναι νόμιμο και αναλογικό, καθώς και τη μεγαλύτερη αξιοποίηση της τοπικής τηλεμετρίας και των μεταδεδομένων, όπως ο χρονισμός των συνδέσεων.

Όταν ένα IPS ή ένα αυτοματοποιημένο σχέδιο ενεργειών αποκρίνεται, μπορεί να \textbf{απορρίψει} τα επίμαχα πακέτα, να \textbf{επαναφέρει} τη σύνδεση TCP, να \textbf{αποκλείσει} τη διεύθυνση προέλευσης για ένα διάστημα, να \textbf{αναδιαμορφώσει} ένα τείχος προστασίας ή να \textbf{απομονώσει} ένα τερματικό από το δίκτυο. Ο αυτοματισμός μειώνει δραματικά τον χρόνο απόκρισης, όμως ένας λανθασμένα ρυθμισμένος κανόνας μπορεί επίσης να διακόψει νόμιμη επιχειρησιακή δραστηριότητα. Γι' αυτό ο αυτόματος αποκλεισμός επιφυλάσσεται για ανιχνεύσεις υψηλής βεβαιότητας, ενώ οι ειδοποιήσεις χαμηλότερης βεβαιότητας διοχετεύονται στους αναλυτές.

\section{Προληπτική αναζήτηση απειλών και τηλεμετρία τερματικών}\label{sec:11.3}
Η \textbf{προληπτική αναζήτηση απειλών} (threat hunting) ξεκινά από την παραδοχή ότι ένας αντίπαλος μπορεί να βρίσκεται ήδη στο περιβάλλον χωρίς να έχει ενεργοποιήσει κανέναν συναγερμό. Η αναζήτηση καθοδηγείται από μια \textbf{υπόθεση}, που συνήθως προέρχεται από πληροφορίες απειλών ή από το ATT\&CK —για παράδειγμα «ένας επιτιθέμενος χρησιμοποιεί προγραμματισμένες εργασίες για μόνιμη παρουσία στους διακομιστές μας». Ο αναλυτής υποβάλλει στη συνέχεια ερωτήματα στη διαθέσιμη τηλεμετρία, διερευνά τις ανωμαλίες και καταλήγει σε συμπέρασμα. Όποιο κι αν είναι το αποτέλεσμα, μια καλή αναζήτηση ολοκληρώνεται με τη μετατροπή της λογικής της σε αυτοματοποιημένη ανίχνευση, ώστε η ίδια συμπεριφορά να εντοπίζεται στο μέλλον.

Η αναζήτηση απειλών εξαρτάται από πλούσια \textbf{τηλεμετρία τερματικών}. Εργαλεία όπως το Sysmon της Microsoft καταγράφουν τη δημιουργία διεργασιών με την πλήρη γραμμή εντολών, τις δικτυακές συνδέσεις, τη φόρτωση οδηγών και τις αλλαγές στο μητρώο· οι πλατφόρμες EDR συλλέγουν παρόμοια δεδομένα σε μεγάλη κλίμακα· και στο Linux, το \texttt{auditd} και οι αισθητήρες που βασίζονται στο eBPF παρέχουν ισοδύναμη ορατότητα. Η \textbf{ανάλυση μνήμης} με πλαίσια όπως το Volatility μπορεί να αποκαλύψει εγχυμένο κώδικα, κρυφές διεργασίες και αποκρυπτογραφημένες ρυθμίσεις κακόβουλου λογισμικού που δεν γράφτηκαν ποτέ στον δίσκο —ικανότητα απαραίτητη απέναντι στις τεχνικές χωρίς αρχεία που περιγράφηκαν στο Κεφάλαιο 4.

\section{YARA, Sigma και ανιχνεύσεις ως κώδικας}\label{sec:11.4}
Η \textbf{YARA} είναι μια γλώσσα αντιστοίχισης προτύπων για την ταξινόμηση αρχείων και περιοχών μνήμης. Ένας κανόνας έχει τρία μέρη: \emph{meta} (περιγραφικές πληροφορίες), \emph{strings} (πρότυπα κειμένου, δεκαεξαδικά ή κανονικές εκφράσεις) και μια \emph{condition} που τα συνδυάζει, για παράδειγμα \texttt{uint16(0) == 0x5A4D and 2 of (\$s*)}, που σημαίνει «εκτελέσιμο Windows που περιέχει τουλάχιστον δύο από τις δηλωμένες συμβολοσειρές». Οι καλοί κανόνες στοχεύουν χαρακτηριστικά που ο δημιουργός του κακόβουλου λογισμικού δεν μπορεί να αλλάξει εύκολα —μοναδικές ακολουθίες κώδικα, δομές ρυθμίσεων— και όχι μία μόνο συμβολοσειρά, και δοκιμάζονται σε συλλογή νόμιμων αρχείων για την αποφυγή ψευδώς θετικών αποτελεσμάτων.

Ενώ η YARA επιθεωρεί αρχεία, η \textbf{Sigma} περιγράφει ύποπτα μοτίβα σε \emph{συμβάντα αρχείων καταγραφής}, σε μορφή ανεξάρτητη από προμηθευτή, η οποία μπορεί να μετατραπεί σε ερωτήματα για διαφορετικές πλατφόρμες SIEM. Η προσέγγιση των \textbf{ανιχνεύσεων ως κώδικα} (detection-as-code) εφαρμόζει και στις δύο την πειθαρχία της μηχανικής λογισμικού: οι κανόνες αποθηκεύονται σε σύστημα ελέγχου εκδόσεων, ελέγχονται από ομοτίμους, δοκιμάζονται αυτόματα σε δείγματα δεδομένων, αντιστοιχίζονται σε τεχνικές του ATT\&CK και εγκαθίστανται μέσω μιας ροής ανάπτυξης. Οι πλατφόρμες \textbf{SOAR} (ενορχήστρωση, αυτοματοποίηση και απόκριση ασφάλειας) εκτελούν στη συνέχεια σχέδια ενεργειών όταν ενεργοποιείται μια ανίχνευση —εμπλουτίζοντας την ειδοποίηση, ανοίγοντας ένα δελτίο και, αν η βεβαιότητα είναι υψηλή, απομονώνοντας τον υπολογιστή.

\section{Σχεδιασμός απομονωμένων εργαστηρίων ανάλυσης}\label{sec:11.5}
Η μηχανική ανιχνεύσεων και η ανάλυση κακόβουλου λογισμικού απαιτούν έναν χώρο όπου ο εχθρικός κώδικας μπορεί να εκτελεστεί χωρίς να θέτει σε κίνδυνο πραγματικά συστήματα. Ένα ασφαλές εργαστήριο βασίζεται στην \textbf{εικονικοποίηση}: τα μηχανήματα ανάλυσης λειτουργούν ως εικονικές μηχανές σε αποκλειστικό φυσικό υπολογιστή, με \textbf{στιγμιότυπα} (snapshots) που επιτρέπουν την επαναφορά τους σε καθαρή κατάσταση μετά από κάθε πείραμα. Η δικτύωση ορίζεται ως \emph{host-only} ή ως εσωτερικό δίκτυο χωρίς διαδρομή προς το διαδίκτυο ή το περιβάλλον παραγωγής· όπου πρέπει να παρατηρηθεί η δικτυακή συμπεριφορά, προσομοιωμένες υπηρεσίες όπως τα INetSim ή FakeNet απαντούν στα αιτήματα του κακόβουλου λογισμικού.

Η απομόνωση πρέπει να είναι και διαδικαστική. Οι κοινόχρηστοι φάκελοι και ο συγχρονισμός του προχείρου μεταξύ φυσικού και εικονικού μηχανήματος πρέπει να απενεργοποιούνται, τα δείγματα να αποθηκεύονται σε αρχεία συμπίεσης με κωδικό και με σαφή σήμανση, ενώ οι αναλυτές πρέπει να θυμούνται ότι ορισμένα κακόβουλα προγράμματα ανιχνεύουν τις εικονικές μηχανές και αλλάζουν συμπεριφορά. Οι προφυλάξεις αυτές εφαρμόζουν την ίδια λογική περιορισμού που διέπει την απόκριση σε περιστατικά.

\section{Ο κύκλος απόκρισης σε περιστατικά (NIST SP 800-61)}\label{sec:11.6}
Η Ειδική Έκδοση 800-61 του NIST περιγράφει την απόκριση σε περιστατικά ως κύκλο τεσσάρων φάσεων. Πρώτη έρχεται η \textbf{προετοιμασία}, η οποία καθορίζει όλα τα υπόλοιπα: εγκεκριμένο σχέδιο, καθορισμένοι ρόλοι, κατάλογοι επαφών, πρότυπα επικοινωνίας, εργαλεία, καταγραφή και τακτικές ασκήσεις. Η \textbf{ανίχνευση και ανάλυση} επιβεβαιώνει ότι σημειώθηκε περιστατικό, προσδιορίζει το εύρος και τη σοβαρότητά του και τεκμηριώνει κάθε βήμα. Ο \textbf{περιορισμός, η εξάλειψη και η ανάκαμψη} περιορίζουν τη ζημιά —βραχυπρόθεσμα με την απομόνωση συστημάτων και την απενεργοποίηση παραβιασμένων λογαριασμών, μακροπρόθεσμα με την ανακατασκευή τους— αφαιρούν τα ερείσματα του επιτιθέμενου και αποκαθιστούν την κανονική λειτουργία από αξιόπιστες πηγές. Η \textbf{δραστηριότητα μετά το περιστατικό} περιλαμβάνει μια ανασκόπηση διδαγμάτων χωρίς απόδοση ευθυνών μέσα σε λίγες εβδομάδες και επιστρέφει τις βελτιώσεις στη φάση της προετοιμασίας.

Δύο εντάσεις διατρέχουν κάθε περιστατικό. Η πρώτη είναι μεταξύ \textbf{ταχύτητας και αποδεικτικών στοιχείων}: η απενεργοποίηση ενός μηχανήματος σταματά την επίθεση, αλλά καταστρέφει τα περιεχόμενα της μνήμης· γι' αυτό τα αποδεικτικά στοιχεία πρέπει, όπου είναι δυνατόν, να συλλέγονται πριν από τις ενέργειες που διαταράσσουν το σύστημα (Κεφάλαιο 12). Η δεύτερη είναι μεταξύ \textbf{περιορισμού και συλλογής πληροφοριών}: η πρόωρη απομάκρυνση του επιτιθέμενου, πριν γίνει γνωστό ολόκληρο το αποτύπωμά του, μπορεί απλώς να τον προειδοποιήσει και να τον ωθήσει να χρησιμοποιήσει άλλη κερκόπορτα. Η αναθεωρημένη καθοδήγηση (Rev. 3, 2025) ευθυγραμμίζει αυτές τις δραστηριότητες με τις λειτουργίες του NIST Cybersecurity Framework 2.0.

\begin{notebox}{Νομικές και κανονιστικές προθεσμίες}
Πολλά περιστατικά ενεργοποιούν υποχρεωτικές γνωστοποιήσεις —για παράδειγμα, εντός 72 ωρών στην εποπτική αρχή βάσει του ΓΚΠΔ για παραβιάσεις προσωπικών δεδομένων και έγκαιρη προειδοποίηση εντός 24 ωρών βάσει της NIS2 για σημαντικά περιστατικά (Κεφάλαιο 13). Το σχέδιο απόκρισης πρέπει να προσδιορίζει ποιος αποφασίζει και ποιος προβαίνει στη γνωστοποίηση.
\end{notebox}

\section{Λειτουργίες SOC: βαθμίδες, η ροή του SIEM και σχεδιασμός σεναρίων χρήσης}\label{sec:11.7}
Ένα \textbf{Κέντρο Επιχειρήσεων Ασφάλειας} (Security Operations Centre, SOC) είναι η ομάδα και η υποδομή που παρακολουθούν συνεχώς έναν οργανισμό. Η εργασία κατανέμεται παραδοσιακά σε βαθμίδες: οι αναλυτές της \textbf{Βαθμίδας 1} διαλογούν τις εισερχόμενες ειδοποιήσεις, η \textbf{Βαθμίδα 2} διερευνά σε βάθος τα επιβεβαιωμένα περιστατικά και οι ειδικοί της \textbf{Βαθμίδας 3} αναζητούν απειλές, αναλύουν κακόβουλο λογισμικό και σχεδιάζουν ανιχνεύσεις. Το κεντρικό εργαλείο είναι η πλατφόρμα \textbf{SIEM} (διαχείριση πληροφοριών και συμβάντων ασφάλειας), της οποίας η ροή \emph{συλλέγει} αρχεία καταγραφής από πολλές πηγές, τα \emph{κανονικοποιεί} σε ένα κοινό σχήμα, τα \emph{εμπλουτίζει} με πλαίσιο, όπως οι υπεύθυνοι των στοιχείων και οι πληροφορίες απειλών, \emph{συσχετίζει} συμβάντα σε ειδοποιήσεις και τα \emph{αποθηκεύει} για διερεύνηση και κανονιστική συμμόρφωση.

Ένα SIEM είναι τόσο καλό όσο τα \textbf{σενάρια χρήσης} του —οι συγκεκριμένες απειλές που έχει ρυθμιστεί να ανιχνεύει. Κάθε σενάριο χρήσης πρέπει να δηλώνει την απειλή που αντιμετωπίζει (ιδανικά ως τεχνική του ATT\&CK), τις πηγές δεδομένων που απαιτεί, τη λογική ανίχνευσης, το αναμενόμενο ποσοστό ψευδώς θετικών και τη διαδικασία απόκρισης. Η απόδοση ενός SOC μετράται συνήθως μέσω του \textbf{μέσου χρόνου ανίχνευσης (MTTD)} και του \textbf{μέσου χρόνου απόκρισης (MTTR)}, σε συνδυασμό με το ποσοστό των ειδοποιήσεων που αποδεικνύονται αξιοποιήσιμες —βασικό δείκτη του φόρτου εργασίας και της εξουθένωσης των αναλυτών.

\section*{Βασικοί όροι}
\begin{description}[style=nextline,leftmargin=1.2em]
  \item[NIDS / HIDS] Δικτυακά και τοπικά συστήματα ανίχνευσης εισβολών.
  \item[Κανονικοποίηση κίνησης] Ανασύνθεση της κίνησης όπως θα την ανασυνέθετε ο προορισμός, για την ακύρωση τεχνικών αποφυγής.
  \item[Προληπτική αναζήτηση απειλών] Προληπτική αναζήτηση αντιπάλων που διέφυγαν της ανίχνευσης, καθοδηγούμενη από υποθέσεις.
  \item[YARA / Sigma] Γλώσσες κανόνων για την αντιστοίχιση προτύπων σε αρχεία (YARA) και σε συμβάντα καταγραφής (Sigma).
  \item[Περιορισμός] Ενέργειες απόκρισης που περιορίζουν τη διάδοση και την επίπτωση μιας επίθεσης.
  \item[SIEM] Πλατφόρμα που συλλέγει, κανονικοποιεί, συσχετίζει και αποθηκεύει συμβάντα ασφάλειας.
\end{description}

\section*{Σύνοψη κεφαλαίου}
\begin{itemize}
  \item Οι δικτυακοί και οι τοπικοί αισθητήρες παρέχουν συμπληρωματική ορατότητα· η ανίχνευση βάσει υπογραφών και βάσει ανωμαλιών έχει η καθεμιά πλεονεκτήματα και αδυναμίες.
  \item Οι επιτιθέμενοι αποφεύγουν την επιθεώρηση μέσω κατακερματισμού, κωδικοποίησης και κρυπτογράφησης· η κανονικοποίηση και η τοπική τηλεμετρία τούς αντιμετωπίζουν.
  \item Η προληπτική αναζήτηση απειλών καθοδηγείται από υποθέσεις και πρέπει να καταλήγει στη δημιουργία νέων αυτοματοποιημένων ανιχνεύσεων.
  \item Οι κανόνες YARA και Sigma, όταν διαχειρίζονται ως κώδικας και συνδέονται με SOAR, μετατρέπουν τη γνώση σε επαναλήψιμη ανίχνευση και απόκριση.
  \item Η απόκριση σε περιστατικά ακολουθεί τον κύκλο του NIST, εξισορροπώντας την ταχύτητα με τη διαφύλαξη των αποδεικτικών στοιχείων, και το SOC την υλοποιεί μέσω βαθμίδων, ενός SIEM και μετρήσιμων σεναρίων χρήσης.
\end{itemize}

\section*{Ερωτήσεις ανασκόπησης}
\begin{enumerate}
  \item Γιατί ένας τοπικός αισθητήρας μπορεί να ανιχνεύσει μια επίθεση μέσα σε κίνηση HTTPS που διαφεύγει από ένα δικτυακό IDS;
  \item Διατυπώστε μια υπόθεση αναζήτησης για την υποκλοπή διαπιστευτηρίων από τη διεργασία LSASS και αναφέρετε την τηλεμετρία που θα χρειαζόσασταν.
  \item Γιατί ένας κανόνας YARA δεν πρέπει να βασίζεται σε μία μόνο συμβολοσειρά, όπως ένα όνομα τομέα C2;
  \item Κατά τη διάρκεια ενός περιστατικού, ένας προϊστάμενος θέλει να απενεργοποιήσει αμέσως τον επηρεαζόμενο διακομιστή. Συζητήστε τα αντισταθμιστικά οφέλη και κόστη.
\end{enumerate}

% ============================================================
% Chapter 12: Malware Analysis, Digital Forensics and OSINT Reconnaissance
% Language: Greek — edit freely; regenerate from src/content if preferred.
% ============================================================
\chapter{Ανάλυση Κακόβουλου Λογισμικού, Ψηφιακή Εγκληματολογία και Αναγνώριση μέσω OSINT}\label{ch:12}
\chaptermeta{Στατική αρχική αξιολόγηση · Δυναμική ανάλυση και ανάλυση συμπεριφοράς · Παθητική και ενεργητική αναγνώριση · Απαρίθμηση DNS · Δεοντολογία OSINT · Σειρά πτητικότητας · Χρονολόγια και ανάκτηση δεδομένων}{Επίπεδο: Προχωρημένο επίπεδο · Διερεύνηση}{Χρόνος μελέτης: 10–12 ώρες}

Όταν σημειώνεται ένα περιστατικό, προκύπτουν τρία ερευνητικά ερωτήματα. Τι κάνει ο κακόβουλος κώδικας; Τι μπορεί να αποκαλύψει σε έναν αντίπαλο η έκθεση του οργανισμού; Και τι ακριβώς συνέβη στα επηρεαζόμενα συστήματα; Το κεφάλαιο εισάγει τους κλάδους που απαντούν σε αυτά: την ανάλυση κακόβουλου λογισμικού, την αναγνώριση μέσω πληροφοριών ανοιχτών πηγών (OSINT) και την ψηφιακή εγκληματολογία.

Οι τρεις κλάδοι μοιράζονται μια κοινή δεοντολογία. Οι αναλυτές εργάζονται σε ελεγχόμενα περιβάλλοντα, συλλέγουν πληροφορίες νόμιμα και χειρίζονται τα αποδεικτικά στοιχεία με τρόπο που διαφυλάσσει την ακεραιότητα και το παραδεκτό τους. Το κεφάλαιο κινείται από τη στατική και δυναμική ανάλυση ενός ύποπτου αρχείου, στη χαρτογράφηση της εξωτερικής επιφάνειας επίθεσης ενός οργανισμού και, τέλος, στις αρχές της εγκληματολογικής απόκτησης δεδομένων και στην ανασύνθεση των γεγονότων από τα ίχνη του συστήματος αρχείων.

\begin{outcomesbox}{Μαθησιακά αποτελέσματα}
Μετά τη μελέτη του κεφαλαίου θα πρέπει να μπορείτε:
\begin{itemize}
  \item Να πραγματοποιείτε στατική αρχική αξιολόγηση ενός ύποπτου αρχείου με συνόψεις, κεφαλίδες, εντροπία και συμβολοσειρές.
  \item Να σχεδιάζετε μια ασφαλή δυναμική ανάλυση και να ερμηνεύετε τα ίχνη συμπεριφοράς.
  \item Να διακρίνετε την παθητική από την ενεργητική αναγνώριση και να χαρτογραφείτε την επιφάνεια επίθεσης ενός οργανισμού μέσω τεχνικών DNS και μηχανών αναζήτησης.
  \item Να εφαρμόζετε νομικούς, δεοντολογικούς κανόνες και κανόνες επιχειρησιακής ασφάλειας στην εργασία OSINT.
  \item Να εφαρμόζετε τη σειρά πτητικότητας και την αλυσίδα επιμέλειας και να κατασκευάζετε χρονολόγιο από τα ίχνη του συστήματος αρχείων.
\end{itemize}
\end{outcomesbox}

\section{Στατική αρχική αξιολόγηση: συνόψεις, κεφαλίδες, εντροπία, συμβολοσειρές και packers}\label{sec:12.1}
Η \textbf{στατική ανάλυση} εξετάζει ένα αρχείο χωρίς να το εκτελεί και αποτελεί πάντοτε το πρώτο και ασφαλέστερο βήμα. Ο αναλυτής ξεκινά υπολογίζοντας \textbf{κρυπτογραφικές συνόψεις} (SHA-256), ώστε να ταυτοποιήσει μοναδικά το δείγμα και να αναζητήσει προηγούμενες αναφορές σε βάσεις πληροφοριών απειλών. Στη συνέχεια, η \textbf{κεφαλίδα} του αρχείου αποκαλύπτει τον πραγματικό του τύπο ανεξάρτητα από την επέκταση: τα εκτελέσιμα των Windows ξεκινούν με τα bytes \texttt{MZ} και η δομή τους \emph{Portable Executable (PE)} περιλαμβάνει χρονοσφραγίδες μεταγλώττισης, ενότητες και εισαγόμενες συναρτήσεις. Εισαγωγές όπως οι \texttt{VirtualAllocEx} και \texttt{CreateRemoteThread} υποδηλώνουν έγχυση σε διεργασίες· η \texttt{CryptEncrypt} μαζί με συναρτήσεις απαρίθμησης αρχείων μπορεί να παραπέμπει σε λυτρισμικό.

Η \textbf{εντροπία} μετρά την τυχαιότητα σε κλίμακα από 0 έως 8 bit ανά byte. Ο συνηθισμένος κώδικας έχει μέτρια εντροπία, ενώ ενότητες με εντροπία κοντά στο 8 είναι πιθανότατα συμπιεσμένες ή κρυπτογραφημένες —τυπική ένδειξη ενός \textbf{packer}, δηλαδή ενός εργαλείου που περιτυλίγει το πραγματικό φορτίο ώστε αυτό να αποκαλύπτεται μόνο στη μνήμη. Η εξαγωγή αναγνώσιμων \textbf{συμβολοσειρών} μπορεί να αποκαλύψει URL, διευθύνσεις IP, διαδρομές μητρώου ή μηνύματα σφάλματος, αν και οι δημιουργοί κακόβουλου λογισμικού συχνά τις συσκοτίζουν. Όταν απαιτείται βαθύτερη κατανόηση, οι \textbf{αποσυμβολομεταφραστές και αποσυμπιλητές} (disassemblers, decompilers), όπως τα Ghidra ή IDA, μεταφράζουν τον κώδικα μηχανής σε συμβολική γλώσσα ή ψευδο-C, επιτρέποντας στον αναλυτή να παρακολουθήσει τη λογική του προγράμματος.

\section{Δυναμική ανάλυση και ανάλυση συμπεριφοράς}\label{sec:12.2}
Η \textbf{δυναμική ανάλυση} εκτελεί το δείγμα στο απομονωμένο εργαστήριο που περιγράφηκε στο Κεφάλαιο 11 και παρατηρεί τι κάνει. Αυτοματοποιημένα \textbf{περιβάλλοντα απομονωμένης εκτέλεσης} (sandboxes), όπως το CAPE ή αντίστοιχα εμπορικά, εκτελούν το αρχείο, καταγράφουν τη δραστηριότητά του και παράγουν αναφορά μέσα σε λίγα λεπτά. Η χειροκίνητη ανάλυση προσφέρει περισσότερο έλεγχο: ο αναλυτής λαμβάνει στιγμιότυπο, ενεργοποιεί εργαλεία παρακολούθησης, εκτελεί το δείγμα, αλληλεπιδρά μαζί του αν χρειάζεται και στη συνέχεια συγκρίνει την κατάσταση του συστήματος πριν και μετά την εκτέλεση.

Ο \textbf{έλεγχος συμπεριφοράς} εστιάζει στα ίχνη που αφήνει το δείγμα. Στο σύστημα αρχείων, οι αναλυτές αναζητούν αρχεία που δημιουργήθηκαν και έγγραφα που τροποποιήθηκαν· στο μητρώο, νέα κλειδιά \emph{Run} ή υπηρεσίες· στο δέντρο διεργασιών, θυγατρικές διεργασίες και έγχυση σε νόμιμες διεργασίες· και στο δίκτυο, ερωτήματα DNS και συνδέσεις με διακομιστές διοίκησης και ελέγχου. Εργαλεία όπως τα Process Monitor, Process Explorer, Regshot και Wireshark καταγράφουν αυτά τα συμβάντα, ενώ η παρακολούθηση κλήσεων API αποκαλύπτει ποιες λειτουργίες του συστήματος καλεί το κακόβουλο λογισμικό. Οι αναλυτές πρέπει να θυμούνται ότι ένα εξελιγμένο κακόβουλο πρόγραμμα μπορεί να αναγνωρίσει το sandbox —ελέγχοντας για εικονικό υλικό, μικρό χρόνο λειτουργίας ή απουσία δραστηριότητας χρήστη— και να παραμείνει αδρανές· επομένως, μια «καθαρή» αναφορά sandbox δεν αποτελεί ποτέ απόδειξη αθωότητας.

\section{Αναγνώριση: παθητικές και ενεργητικές τεχνικές, dorking και DNS}\label{sec:12.3}
Η αναγνώριση είναι η πρώτη φάση του Kill Chain, και την πραγματοποιούν επίσης οι αμυνόμενοι, ώστε να βλέπουν τον οργανισμό τους όπως θα τον έβλεπε ένας επιτιθέμενος. Η \textbf{παθητική αναγνώριση} συλλέγει πληροφορίες χωρίς άμεση αλληλεπίδραση με τα συστήματα του στόχου: δημόσιους ιστοτόπους, μέσα κοινωνικής δικτύωσης, αγγελίες εργασίας που αποκαλύπτουν τις τεχνολογίες που χρησιμοποιούνται, εγγραφές WHOIS, αρχεία διαφάνειας πιστοποιητικών και βάσεις δεδομένων σαρώσεων ολόκληρου του διαδικτύου, όπως τα Shodan ή Censys. Η \textbf{ενεργητική αναγνώριση} αλληλεπιδρά με τον στόχο —σάρωση θυρών, λήψη εμβλημάτων υπηρεσιών (banner grabbing), δοκιμή καταλόγων με ωμή βία— και είναι επομένως ανιχνεύσιμη και νομικά επιτρεπτή μόνο με εξουσιοδότηση.

Οι μηχανές αναζήτησης είναι ισχυρά εργαλεία αναγνώρισης. Το \textbf{Google dorking} χρησιμοποιεί προηγμένους τελεστές, όπως οι \texttt{site:}, \texttt{filetype:} και \texttt{intitle:}, για τον εντοπισμό εκτεθειμένων εγγράφων, σελίδων σύνδεσης, αρχείων ρυθμίσεων ή λιστών καταλόγων που δεν προορίζονταν ποτέ για δημόσια χρήση. Η \textbf{απαρίθμηση DNS} χαρτογραφεί την παρουσία ενός οργανισμού στο διαδίκτυο: υποβολή ερωτημάτων για τύπους εγγραφών (A, MX, TXT, NS), ανακάλυψη υποτομέων μέσω αρχείων διαφάνειας πιστοποιητικών και βάσεων παθητικού DNS, και έλεγχος αν λανθασμένα ρυθμισμένοι διακομιστές ονομάτων επιτρέπουν \emph{μεταφορές ζώνης} που αποκαλύπτουν κάθε εγγραφή. Ξεχασμένοι υποτομείς που εξακολουθούν να παραπέμπουν σε καταργημένους πόρους νέφους μπορούν ακόμη και να \emph{καταληφθούν} από έναν επιτιθέμενο —κίνδυνος που οι τακτικές επισκοπήσεις της επιφάνειας επίθεσης έχουν σχεδιαστεί να εντοπίζουν.

\section{Δεοντολογία OSINT, επιχειρησιακή ασφάλεια και αξιοποιήσιμη αναφορά}\label{sec:12.4}
Το ότι μια πληροφορία είναι δημόσια προσβάσιμη δεν σημαίνει ότι μπορεί να συλλεχθεί και να χρησιμοποιηθεί χωρίς όρια. Η εργασία OSINT πρέπει να σέβεται τον νόμο —οι κανόνες προστασίας δεδομένων, όπως ο ΓΚΠΔ, εφαρμόζονται στα προσωπικά δεδομένα ακόμη κι όταν αυτά είναι δημόσια— καθώς και τους όρους χρήσης των πλατφορμών και το συμφωνημένο πεδίο της ανάθεσης. Οι αναλυτές πρέπει να συλλέγουν μόνο ό,τι είναι αναγκαίο για τον δηλωμένο σκοπό και να μην επιχειρούν ποτέ πρόσβαση σε λογαριασμούς ή παράκαμψη ελέγχων πρόσβασης, κάτι που θα μετέτρεπε την αναγνώριση σε εισβολή.

Η \textbf{επιχειρησιακή ασφάλεια} (OPSEC) προστατεύει την ίδια την έρευνα: οι αναλυτές χρησιμοποιούν αποκλειστικά ερευνητικά περιβάλλοντα και λογαριασμούς, αποφεύγουν να αποκαλύψουν στον στόχο την ταυτότητα ή τον οργανισμό τους και τεκμηριώνουν τις πηγές τους με χρονοσφραγίδες και αρχειοθετημένα αντίγραφα, ώστε τα ευρήματα να είναι αναπαραγώγιμα. Η αξία της αναγνώρισης βρίσκεται σε όσα ακολουθούν. Μια \textbf{αναφορά επιφάνειας επίθεσης} πρέπει να μετατρέπει τα ευρήματα σε ιεραρχημένα δελτία αποκατάστασης με συγκεκριμένο υπεύθυνο —«κατάργηση του ορφανού υποτομέα \texttt{old-shop.example.com}», «αφαίρεση του εκτεθειμένου αρχείου αντιγράφου ασφαλείας»— αντί να παρουσιάζει έναν αδιαφοροποίητο κατάλογο ανακαλύψεων.

\section{Αρχές εγκληματολογικής ανάλυσης, σειρά πτητικότητας και αλυσίδα επιμέλειας}\label{sec:12.5}
Η \textbf{ψηφιακή εγκληματολογία} είναι η επιστημονική συλλογή, διαφύλαξη, ανάλυση και παρουσίαση ψηφιακών αποδεικτικών στοιχείων με τρόπο που μπορεί να αντέξει τον νομικό έλεγχο. Η θεμελιώδης αρχή της, που απορρέει από την αρχή ανταλλαγής του Locard, είναι ότι κάθε ενέργεια αφήνει ίχνος —συμπεριλαμβανομένων των ενεργειών του ίδιου του ερευνητή. Τα αποδεικτικά στοιχεία πρέπει, επομένως, να αποκτώνται με τρόπο που τα μεταβάλλει όσο το δυνατόν λιγότερο, και κάθε βήμα πρέπει να τεκμηριώνεται. Οι δίσκοι αντιγράφονται bit προς bit σε εγκληματολογικά είδωλα με τη χρήση \textbf{αναστολέων εγγραφής} (write blockers), και κάθε είδωλο επαληθεύεται συγκρίνοντας τη σύνοψή του με εκείνη του πρωτοτύπου.

Επειδή ορισμένα δεδομένα εξαφανίζονται ταχύτερα από άλλα, τα αποδεικτικά στοιχεία συλλέγονται σύμφωνα με τη \textbf{σειρά πτητικότητας} (RFC 3227): πρώτα οι καταχωρητές και η κρυφή μνήμη του επεξεργαστή, έπειτα η κύρια μνήμη (RAM), συμπεριλαμβανομένων των διεργασιών σε εκτέλεση και των δικτυακών συνδέσεων· στη συνέχεια τα προσωρινά αρχεία και ο χώρος εναλλαγής· έπειτα τα περιεχόμενα του δίσκου· και, τέλος, τα απομακρυσμένα αρχεία καταγραφής και τα αρχειακά μέσα. Αν αποσυνδεθεί πρώτα το καλώδιο ρεύματος, καταστρέφονται τα πιο πτητικά —και συχνά τα πιο αποκαλυπτικά— αποδεικτικά στοιχεία. Κάθε στοιχείο καταχωρίζεται σε ένα αρχείο \textbf{αλυσίδας επιμέλειας} (chain of custody), το οποίο δηλώνει ποιος το συνέλεξε, πότε, πώς, πού φυλάχθηκε και σε ποιον παραδόθηκε. Ένα κενό σε αυτή την αλυσίδα μπορεί να καταστήσει μη παραδεκτά αποδεικτικά στοιχεία που κατά τα λοιπά είναι έγκυρα.

\section{Ίχνη του συστήματος αρχείων, ανάκτηση δεδομένων και χρονολόγια}\label{sec:12.6}
Τα συστήματα αρχείων καταγράφουν πολύ περισσότερα από τα περιεχόμενα των αρχείων. Στα Windows, ο \textbf{Κύριος Πίνακας Αρχείων} (Master File Table, MFT) του NTFS αποθηκεύει μεταδεδομένα για κάθε αρχείο, συμπεριλαμβανομένων τεσσάρων χρονοσφραγίδων —δημιουργίας, τροποποίησης, πρόσβασης και αλλαγής της εγγραφής MFT— σε δύο χωριστά χαρακτηριστικά, κάτι που βοηθά τους ερευνητές να εντοπίζουν το \emph{timestomping}, δηλαδή τη σκόπιμη παραποίηση χρονοσφραγίδων. Άλλα πολύτιμα ίχνη είναι το ημερολόγιο αλλαγών \textbf{\$UsnJrnl}, τα αρχεία \textbf{Prefetch} που αποδεικνύουν ότι ένα πρόγραμμα εκτελέστηκε, τα αρχεία συντομεύσεων \textbf{LNK} και οι \textbf{Jump Lists} που αποκαλύπτουν πρόσφατα ανοιγμένα έγγραφα, τα \textbf{ShellBags} που καταγράφουν την πρόσβαση σε φακέλους και τα αρχεία συμβάντων των Windows. Η διαγραφή ενός αρχείου συνήθως αφαιρεί μόνο την εγγραφή του στον κατάλογο, οπότε τα περιεχόμενά του μπορεί να παραμένουν στον δίσκο έως ότου επαναχρησιμοποιηθεί ο χώρος.

Η \textbf{ανάκτηση αρχείων με βάση το περιεχόμενο} (file carving) ανακτά τέτοια διαγραμμένα δεδομένα αναζητώντας στον μη κατανεμημένο χώρο γνωστές υπογραφές αρχείων —για παράδειγμα την αρχή και το τέλος μιας εικόνας JPEG— χωρίς να βασίζεται στα μεταδεδομένα του συστήματος αρχείων. Τέλος, οι ερευνητές συνδυάζουν τις χρονοσφραγίδες όλων των πηγών σε ένα ενιαίο \textbf{υπερ-χρονολόγιο} (super-timeline), με εργαλεία όπως τα Plaso ή Autopsy. Ένα χρονολόγιο μετατρέπει χιλιάδες μεμονωμένα ίχνη σε αφήγηση: το μήνυμα ηλεκτρονικού ψαρέματος έφτασε στις 09:12, το συνημμένο άνοιξε στις 09:14, το PowerShell εκτελέστηκε στις 09:14:07, μια προγραμματισμένη εργασία δημιουργήθηκε στις 09:15 και τα δεδομένα συμπιέστηκαν και αποστάλθηκαν εκτός οργανισμού στις 02:30 της επόμενης νύχτας.

\section*{Βασικοί όροι}
\begin{description}[style=nextline,leftmargin=1.2em]
  \item[Εντροπία] Μέτρο τυχαιότητας· υψηλές τιμές υποδηλώνουν συμπιεσμένο ή κρυπτογραφημένο (packed) περιεχόμενο.
  \item[Sandbox] Απομονωμένο περιβάλλον που εκτελεί δείγματα και καταγράφει τη συμπεριφορά τους.
  \item[Παθητική αναγνώριση] Συλλογή πληροφοριών χωρίς άμεση αλληλεπίδραση με τα συστήματα του στόχου.
  \item[Σειρά πτητικότητας] Συλλογή αποδεικτικών στοιχείων από την πιο βραχύβια προς την πιο μόνιμη πηγή.
  \item[Αλυσίδα επιμέλειας] Τεκμηριωμένο αρχείο του ποιος χειρίστηκε τα αποδεικτικά στοιχεία, πότε και πώς.
  \item[Υπερ-χρονολόγιο] Ενιαία χρονολογική αλληλουχία γεγονότων που κατασκευάζεται από πολλές πηγές ιχνών.
\end{description}

\section*{Σύνοψη κεφαλαίου}
\begin{itemize}
  \item Η στατική αρχική αξιολόγηση —συνόψεις, κεφαλίδες, εισαγόμενες συναρτήσεις, εντροπία και συμβολοσειρές— είναι το ασφαλές πρώτο βήμα και καθοδηγεί τη βαθύτερη ανάλυση.
  \item Η δυναμική ανάλυση αποκαλύπτει τη συμπεριφορά, όμως το κακόβουλο λογισμικό με τεχνικές αποφυγής σημαίνει ότι ένα «καθαρό» αποτέλεσμα sandbox δεν αποδεικνύει τίποτα.
  \item Η παθητική αναγνώριση, το dorking και η απαρίθμηση DNS επιτρέπουν στους αμυνόμενους να βλέπουν την έκθεσή τους όπως οι επιτιθέμενοι.
  \item Η εργασία OSINT πρέπει να είναι νόμιμη, αναλογική και να λαμβάνει υπόψη την επιχειρησιακή ασφάλεια, ενώ τα αποτελέσματά της πρέπει να μετατρέπονται σε ενέργειες αποκατάστασης με συγκεκριμένο υπεύθυνο.
  \item Η εγκληματολογική ανάλυση ακολουθεί τη σειρά πτητικότητας και την αλυσίδα επιμέλειας, ενώ τα χρονολόγια μετατρέπουν τα ίχνη σε αξιόπιστη αφήγηση.
\end{itemize}

\section*{Ερωτήσεις ανασκόπησης}
\begin{enumerate}
  \item Ένα αρχείο PE έχει μια ενότητα με εντροπία 7,9 και ελάχιστες εισαγόμενες συναρτήσεις. Τι συμπεραίνετε και ποιο είναι το επόμενο βήμα σας;
  \item Γιατί ένα αποτέλεσμα sandbox «καμία κακόβουλη δραστηριότητα» δεν επαρκεί για να θεωρηθεί ασφαλές ένα ύποπτο αρχείο;
  \item Κατατάξτε ως παθητικές ή ενεργητικές τις εξής ενέργειες: αναζήτηση σε αρχεία διαφάνειας πιστοποιητικών, σάρωση με Nmap, αναζήτηση στο Shodan, απόπειρα μεταφοράς ζώνης.
  \item Γιατί πρέπει να συλλέγεται η μνήμη πριν από το είδωλο του δίσκου κατά την απόκριση σε σύστημα σε λειτουργία;
\end{enumerate}

% ============================================================
% Chapter 13: Cyber Resilience, Governance and Future Trends
% Language: Greek — edit freely; regenerate from src/content if preferred.
% ============================================================
\chapter{Κυβερνοανθεκτικότητα, Διακυβέρνηση και Μελλοντικές Τάσεις}\label{ch:13}
\chaptermeta{Αντίγραφα ασφαλείας και ανάκαμψη · BCP, RPO και RTO · Διαχείριση κινδύνου · NIST CSF 2.0, ISO/IEC 27001, CIS · ΓΚΠΔ και NIS2 · Δείκτες · Ασφάλεια OT · Μηδενική εμπιστοσύνη και νέφος}{Επίπεδο: Προχωρημένο επίπεδο · Στρατηγική και διακυβέρνηση}{Χρόνος μελέτης: 10–12 ώρες}

Το τελευταίο κεφάλαιο επιστρέφει στο ερώτημα που τέθηκε στο Κεφάλαιο 1: \emph{μπορεί ο οργανισμός να συνεχίσει να λειτουργεί και να ανακάμψει όταν τα πράγματα πάνε στραβά;} Τα τεχνικά μέτρα είναι αναγκαία αλλά όχι επαρκή. Η ανθεκτικότητα εξαρτάται επίσης από αντίγραφα ασφαλείας που λειτουργούν όταν χρειάζονται, από σχέδια που έχουν δοκιμαστεί, από αποφάσεις για τον κίνδυνο που λαμβάνονται συνειδητά από υπόλογα πρόσωπα και από τη συμμόρφωση με ένα διαρκώς διευρυνόμενο νομικό πλαίσιο.

Το κεφάλαιο καλύπτει πρώτα τη μηχανική ανθεκτικότητας, τις αρχιτεκτονικές αντιγράφων ασφαλείας και την επιχειρησιακή συνέχεια. Στη συνέχεια στρέφεται στη διακυβέρνηση: τη διαχείριση κινδύνου στην πράξη, τα κυριότερα πλαίσια και κανονιστικά κείμενα και τους δείκτες που δείχνουν στη διοίκηση αν η ασφάλεια βελτιώνεται. Ολοκληρώνεται εξετάζοντας περιβάλλοντα και τάσεις που θα διαμορφώσουν τα επόμενα χρόνια —την επιχειρησιακή τεχνολογία και τις κρίσιμες υποδομές, την αρχιτεκτονική μηδενικής εμπιστοσύνης και την ασφάλεια του νέφους και της εφοδιαστικής αλυσίδας.

\begin{outcomesbox}{Μαθησιακά αποτελέσματα}
Μετά τη μελέτη του κεφαλαίου θα πρέπει να μπορείτε:
\begin{itemize}
  \item Να σχεδιάζετε στρατηγική αντιγράφων ασφαλείας που αντέχει σε λυτρισμικό και να ορίζετε RPO και RTO για κρίσιμες υπηρεσίες.
  \item Να καταρτίζετε μητρώο κινδύνων και να επιλέγετε κατάλληλες επιλογές αντιμετώπισης.
  \item Να συγκρίνετε τα NIST CSF 2.0, ISO/IEC 27001 και CIS Controls και να συνοψίζετε τις υποχρεώσεις του ΓΚΠΔ και της NIS2.
  \item Να διακρίνετε τους ουσιαστικούς δείκτες ασφάλειας από τους δείκτες «βιτρίνας».
  \item Να εξηγείτε τους ιδιαίτερους περιορισμούς της ασφάλειας OT, καθώς και τις αρχές της μηδενικής εμπιστοσύνης και της κοινής ευθύνης στο νέφος.
\end{itemize}
\end{outcomesbox}

\section{Μηχανική ανθεκτικότητας και αρχιτεκτονικές αντιγράφων ασφαλείας}\label{sec:13.1}
Η \textbf{μηχανική ανθεκτικότητας} αποδέχεται ότι ορισμένες επιθέσεις και αστοχίες θα επιτύχουν και σχεδιάζει τα συστήματα ώστε να υποβαθμίζονται ομαλά αντί να καταρρέουν. Τα εργαλεία της περιλαμβάνουν πλεονασμό χωρίς κοινά μεμονωμένα σημεία αστοχίας, δυνατότητα απομόνωσης των στοιχείων που έχουν υποστεί ζημιά και προσχεδιασμένους τρόπους περιορισμένης λειτουργίας —για παράδειγμα, ένα νοσοκομείο που μπορεί να συνεχίσει να δέχεται ασθενείς σε χαρτί όταν το ηλεκτρονικό του σύστημα δεν είναι διαθέσιμο.

Τα αντίγραφα ασφαλείας αποτελούν το θεμέλιο της ανάκαμψης, και το σύγχρονο λυτρισμικό τα στοχεύει σκόπιμα. Ο κλασικός \textbf{κανόνας 3-2-1} —τρία αντίγραφα των δεδομένων, σε δύο διαφορετικούς τύπους μέσων, ένα από τα οποία εκτός εγκατάστασης— έχει γι' αυτό επεκταθεί σε \textbf{3-2-1-1-0}: τουλάχιστον ένα αντίγραφο πρέπει να είναι \textbf{εκτός σύνδεσης ή αμετάβλητο} (αποθήκευση εφάπαξ εγγραφής ή κλείδωμα αντικειμένων που ούτε ένας διαχειριστής δεν μπορεί να διαγράψει κατά την περίοδο διατήρησης), και οι επαναφορές πρέπει να δοκιμάζονται, ώστε να υπάρχουν \textbf{μηδέν} ανεπιβεβαίωτα σφάλματα. Τα συστήματα αντιγράφων ασφαλείας χρειάζονται δικά τους, ξεχωριστά διαπιστευτήρια και MFA, διότι διαφορετικά ένας επιτιθέμενος με δικαιώματα διαχειριστή τομέα απλώς θα τα διαγράψει πριν ξεκινήσει την κρυπτογράφηση.

\section{Επιχειρησιακή συνέχεια και ανάκαμψη από καταστροφές: RPO και RTO}\label{sec:13.2}
Ο \textbf{σχεδιασμός επιχειρησιακής συνέχειας} (Business Continuity Planning, BCP) διασφαλίζει ότι οι κρίσιμες επιχειρησιακές λειτουργίες συνεχίζονται κατά τη διάρκεια μιας διαταραχής, ενώ η \textbf{ανάκαμψη από καταστροφές} (Disaster Recovery, DR) εστιάζει στην αποκατάσταση των συστημάτων πληροφορικής που τις υποστηρίζουν. Και οι δύο ξεκινούν με μια \textbf{ανάλυση επιχειρησιακών επιπτώσεων}, η οποία εντοπίζει τις κρίσιμες διεργασίες και τις συνέπειες της διακοπής τους σε βάθος χρόνου. Δύο παράμετροι καθοδηγούν στη συνέχεια τον τεχνικό σχεδιασμό. Ο \textbf{στόχος σημείου ανάκαμψης} (Recovery Point Objective, RPO) είναι η μέγιστη αποδεκτή απώλεια δεδομένων, μετρούμενη σε χρόνο: RPO μίας ώρας απαιτεί αντίγραφα ή αναπαραγωγή τουλάχιστον ανά ώρα. Ο \textbf{στόχος χρόνου ανάκαμψης} (Recovery Time Objective, RTO) είναι ο μέγιστος αποδεκτός χρόνος αποκατάστασης της υπηρεσίας. Οι αυστηρότεροι στόχοι κοστίζουν περισσότερο και πρέπει, επομένως, να δικαιολογούνται από τις επιχειρησιακές επιπτώσεις.

Τα σχέδια είναι αξιόπιστα μόνο εφόσον δοκιμάζονται. Τα \textbf{εγχειρίδια ανάκαμψης} (DR playbooks) περιγράφουν βήμα προς βήμα πώς αποκαθίσταται κάθε κρίσιμο σύστημα, με ποια σειρά (οι υπηρεσίες ταυτότητας και η δικτύωση προηγούνται συνήθως) και ποιος είναι υπεύθυνος. Επικυρώνονται μέσω ασκήσεων επί χάρτου (tabletop), μερικών δοκιμών επαναφοράς και πλήρων ασκήσεων μετάπτωσης, και οι μετρούμενοι χρόνοι ανάκαμψης συγκρίνονται με τον RTO. Ένα σχέδιο που δεν έχει δοκιμαστεί ποτέ πρέπει να θεωρείται ότι δεν λειτουργεί.

\section{Διαχείριση επιχειρησιακού κινδύνου στην πράξη}\label{sec:13.3}
Το Κεφάλαιο 1 όρισε τον κίνδυνο· η διακυβέρνηση μετατρέπει αυτόν τον ορισμό σε διαδικασία διαχείρισης. Ένα \textbf{μητρώο κινδύνων} καταγράφει κάθε κίνδυνο με σαφή διατύπωση (\emph{απειλή} που εκμεταλλεύεται \emph{ευπάθεια} προκαλώντας \emph{επίπτωση} σε \emph{περιουσιακό στοιχείο}), υπεύθυνο, υφιστάμενα μέτρα, εκτιμήσεις πιθανότητας και επίπτωσης, την επιλεγμένη αντιμετώπιση και ημερομηνία επανεξέτασης. Ένας \textbf{πίνακας κινδύνων} —συνήθως πέντε επί πέντε— βοηθά στην οπτικοποίηση και σύγκριση των κινδύνων, αν και οι ποιοτικές κλίμακες πρέπει να ορίζονται προσεκτικά, ώστε το «πιθανό» να σημαίνει το ίδιο για όλους. Ποσοτικές μέθοδοι όπως η FAIR εκφράζουν τον κίνδυνο σε χρηματικούς όρους και μπορούν να υποστηρίξουν αποφάσεις επενδύσεων.

Κάθε κίνδυνος πρέπει να λαμβάνει ρητή \textbf{απόφαση αντιμετώπισης}: \emph{μείωση} με πρόσθετα μέτρα, \emph{μεταφορά} μέσω ασφάλισης ή συμβάσεων, \emph{αποφυγή} με διακοπή της δραστηριότητας ή επίσημη \emph{αποδοχή}. Η αποδοχή είναι θεμιτή μόνο όταν αποφασίζεται από πρόσωπο με την αρμοδιότητα να φέρει τις συνέπειες, τεκμηριώνεται και επανεξετάζεται περιοδικά. Οι τεχνικές βαθμολογίες, όπως το CVSS (Κεφάλαιο 10), αποτελούν δεδομένα εισόδου αυτής της διαδικασίας και όχι υποκατάστατό της.

\section{Πλαίσια και ρυθμίσεις: NIST CSF 2.0, ISO/IEC 27001, CIS, ΓΚΠΔ και NIS2}\label{sec:13.4}
Τα πλαίσια δίνουν δομή σε ένα πρόγραμμα ασφάλειας. Το \textbf{NIST Cybersecurity Framework 2.0} (2024) οργανώνει τα επιδιωκόμενα αποτελέσματα σε έξι λειτουργίες: \emph{Διακυβέρνηση}, \emph{Αναγνώριση}, \emph{Προστασία}, \emph{Ανίχνευση}, \emph{Απόκριση} και \emph{Ανάκαμψη}· η νέα λειτουργία της Διακυβέρνησης υπογραμμίζει ότι η κυβερνοασφάλεια είναι ευθύνη της ηγεσίας. Το \textbf{ISO/IEC 27001} ορίζει τις απαιτήσεις ενός \emph{συστήματος διαχείρισης ασφάλειας πληροφοριών} (ΣΔΑΠ/ISMS) —ενός κύκλου αξιολόγησης κινδύνου, επιλογής μέτρων, παρακολούθησης και βελτίωσης— και αποτελεί τη βάση για ανεξάρτητη πιστοποίηση· το Παράρτημα Α περιλαμβάνει 93 μέτρα αναφοράς. Τα \textbf{CIS Critical Security Controls} είναι ένα ιεραρχημένο σύνολο συγκεκριμένων τεχνικών μέτρων, ομαδοποιημένων σε ομάδες υλοποίησης για οργανισμούς διαφορετικής ωριμότητας.

Το κανονιστικό πλαίσιο καθιστά ορισμένες πρακτικές υποχρεωτικές. Ο \textbf{Γενικός Κανονισμός για την Προστασία Δεδομένων} (ΓΚΠΔ/GDPR) της ΕΕ απαιτεί κατάλληλη ασφάλεια για τα προσωπικά δεδομένα, προστασία δεδομένων ήδη από τον σχεδιασμό και εξ ορισμού, καθώς και γνωστοποίηση των παραβιάσεων προσωπικών δεδομένων στην εποπτική αρχή εντός 72 ωρών. Η \textbf{Οδηγία NIS2} επεκτείνει τις υποχρεώσεις κυβερνοασφάλειας σε ευρύ φάσμα βασικών και σημαντικών οντοτήτων —ενέργεια, υγεία, μεταφορές, ψηφιακές υποδομές κ.ά.— απαιτώντας μέτρα διαχείρισης κινδύνου, ασφάλεια της εφοδιαστικής αλυσίδας, αναφορά περιστατικών (έγκαιρη προειδοποίηση εντός 24 ωρών) και προσωπική λογοδοσία των διοικητικών οργάνων. Στις Ηνωμένες Πολιτείες, ο νόμος \textbf{HIPAA} επιβάλλει ανάλογες διασφαλίσεις για τις πληροφορίες υγείας.

\begin{table}[htbp]
  \centering\small
  \caption{Πλαίσια και κανονιστικά κείμενα με μια ματιά}
  \begin{tabularx}{\linewidth}{>{\raggedright\arraybackslash}X >{\raggedright\arraybackslash}X >{\raggedright\arraybackslash}X >{\raggedright\arraybackslash}X}
    \toprule
    \textbf{Κείμενο} & \textbf{Φύση} & \textbf{Κύρια εστίαση} & \textbf{Τυπική χρήση} \\
    \midrule
    NIST CSF 2.0 & Προαιρετικό πλαίσιο & Αποτελέσματα σε έξι λειτουργίες & Δομή προγράμματος, αποτύπωση ωριμότητας \\
    ISO/IEC 27001 & Πιστοποιήσιμο πρότυπο & Σύστημα διαχείρισης (ΣΔΑΠ) & Πιστοποίηση, διαβεβαίωση πελατών \\
    CIS Controls v8 & Ιεραρχημένο σύνολο μέτρων & Συγκεκριμένα τεχνικά μέτρα & Οδικός χάρτης υλοποίησης \\
    ΓΚΠΔ & Κανονισμός ΕΕ (δεσμευτικός) & Προστασία προσωπικών δεδομένων & Νομική συμμόρφωση, γνωστοποίηση παραβιάσεων \\
    NIS2 & Οδηγία ΕΕ (δεσμευτική) & Ασφάλεια βασικών/σημαντικών οντοτήτων & Διαχείριση κινδύνου, αναφορά περιστατικών \\
    \bottomrule
  \end{tabularx}
\end{table}
\section{Δείκτες ασφάλειας: ουσιαστική πληροφορία έναντι βιτρίνας}\label{sec:13.5}
Η διοίκηση χρειάζεται τεκμήρια ότι η επένδυση στην ασφάλεια μειώνει τον κίνδυνο. Πολλοί συνήθως αναφερόμενοι αριθμοί είναι \textbf{δείκτες βιτρίνας}: τα «εκατομμύρια επιθέσεων που αποκλείστηκαν» μετρούν κυρίως τον αυτοματοποιημένο «θόρυβο» του διαδικτύου, ενώ ο «αριθμός των εργαλείων που έχουν εγκατασταθεί» δεν λέει τίποτα για την αποτελεσματικότητά τους. Οι \textbf{ουσιαστικοί δείκτες} συνδέονται με αποφάσεις και αποτελέσματα: το ποσοστό των κρίσιμων περιουσιακών στοιχείων που καλύπτονται από EDR και καταγραφή· ο χρόνος εφαρμογής ενημερώσεων για ευπάθειες σε συστήματα εκτεθειμένα στο διαδίκτυο που είναι γνωστό ότι αποτελούν αντικείμενο εκμετάλλευσης· οι MTTD και MTTR (Κεφάλαιο 11)· το ποσοστό αναφοράς μηνυμάτων ψαρέματος (Κεφάλαιο 5)· το ποσοστό των προνομιούχων λογαριασμών που προστατεύονται με MFA ανθεκτικό στο ψάρεμα· και το ποσοστό επιτυχίας και η διάρκεια των δοκιμών επαναφοράς αντιγράφων ασφαλείας σε σύγκριση με τον RTO.

Οι καλοί δείκτες έχουν καθορισμένη πηγή δεδομένων, στόχο και τάση στον χρόνο, και ο καθένας πρέπει να υποδεικνύει μια ενέργεια αν κινηθεί προς τη λάθος κατεύθυνση. Η συνεπής αναφορά λίγων τέτοιων δεικτών είναι πολύ πιο χρήσιμη για ένα διοικητικό συμβούλιο από έναν πίνακα με εκατοντάδες ανεξήγητα νούμερα.

\section{Κρίσιμες υποδομές και επιχειρησιακή τεχνολογία}\label{sec:13.6}
Η \textbf{επιχειρησιακή τεχνολογία} (Operational Technology, OT) —βιομηχανικά συστήματα ελέγχου (ICS), συστήματα SCADA και προγραμματιζόμενοι λογικοί ελεγκτές που λειτουργούν δίκτυα ηλεκτρικής ενέργειας, εγκαταστάσεις επεξεργασίας νερού, παραγωγικές μονάδες και ιατροτεχνολογικά προϊόντα— έχει προτεραιότητες διαφορετικές από την πληροφορική γραφείου. Η ασφάλεια των ανθρώπων και η διαθεσιμότητα προηγούνται, τα συστήματα μπορεί να λειτουργούν για είκοσι χρόνια ή και περισσότερο, και πολλά βιομηχανικά πρωτόκολλα, όπως το Modbus, σχεδιάστηκαν χωρίς έλεγχο ταυτότητας. Η εφαρμογή ενημερώσεων μπορεί να απαιτεί διακοπή μιας γραμμής παραγωγής, ενώ μια επιθετική δικτυακή σάρωση μπορεί να θέσει εκτός λειτουργίας έναν ευαίσθητο ελεγκτή. Όπως έδειξαν το Stuxnet και μεταγενέστερες επιθέσεις σε ενεργειακά δίκτυα, οι συνέπειες μπορεί να είναι φυσικές.

Η ασφάλεια OT βασίζεται επομένως σε μεγάλο βαθμό στην αρχιτεκτονική. Το \textbf{μοντέλο Purdue} διαιρεί το περιβάλλον σε επίπεδα, από τις φυσικές διεργασίες (επίπεδο 0) και τους ελεγκτές (επίπεδο 1) έως την επιχειρησιακή πληροφορική (επίπεδα 4–5). Το πρότυπο \textbf{ISA/IEC 62443} ομαδοποιεί τα στοιχεία σε \emph{ζώνες} και επιτρέπει την επικοινωνία μόνο μέσω ελεγχόμενων \emph{αγωγών} (conduits), με μια βιομηχανική αποστρατιωτικοποιημένη ζώνη μεταξύ IT και OT. Η παθητική παρακολούθηση του δικτύου, σχεδιασμένη για βιομηχανικά πρωτόκολλα, παρέχει ορατότητα χωρίς να διαταράσσει τις ευαίσθητες συσκευές, ενώ ο αυστηρός έλεγχος της απομακρυσμένης πρόσβασης των προμηθευτών κλείνει ένα από τα σημεία εισόδου που γίνονται συχνότερα αντικείμενο κατάχρησης.

\section{Μηδενική εμπιστοσύνη, ασφάλεια νέφους και αναδυόμενες απειλές}\label{sec:13.7}
Η \textbf{αρχιτεκτονική μηδενικής εμπιστοσύνης} (Zero Trust Architecture, NIST SP 800-207) εγκαταλείπει την ιδέα ότι οτιδήποτε βρίσκεται μέσα στην περίμετρο του δικτύου είναι αξιόπιστο. Κάθε αίτημα πρόσβασης αξιολογείται ρητά, με βάση την ταυτότητα του χρήστη, την κατάσταση της συσκευής και την ευαισθησία του πόρου· η πρόσβαση παραχωρείται με ελάχιστο προνόμιο, για περιορισμένη συνεδρία· και ο σχεδιασμός προϋποθέτει ότι μια παραβίαση έχει ήδη συμβεί. Ένα \emph{σημείο λήψης απόφασης πολιτικής} αξιολογεί κάθε αίτημα και ένα \emph{σημείο επιβολής πολιτικής} μπροστά από κάθε πόρο εφαρμόζει την απόφαση. Η μηδενική εμπιστοσύνη είναι πορεία και όχι προϊόν: οι οργανισμοί συνήθως προχωρούν από την ισχυρή ταυτότητα και το MFA, στη συμμόρφωση των συσκευών και στη μικροκατάτμηση και, τέλος, σε συνεχείς αποφάσεις πρόσβασης προσαρμοσμένες στον κίνδυνο.

Στο \textbf{νέφος}, η ασφάλεια ακολουθεί ένα \textbf{μοντέλο κοινής ευθύνης}: ο πάροχος προστατεύει την υποκείμενη υποδομή, ενώ ο πελάτης παραμένει υπεύθυνος για τις ταυτότητες, τη διαμόρφωση και τα δεδομένα. Οι περισσότερες παραβιάσεις στο νέφος οφείλονται σε λανθασμένες ρυθμίσεις από την πλευρά του πελάτη —δημόσια αναγνώσιμους χώρους αποθήκευσης, υπερβολικά δικαιώματα, εκτεθειμένα κλειδιά— τις οποίες έχουν σχεδιαστεί να εντοπίζουν τα εργαλεία \emph{διαχείρισης της στάσης ασφάλειας στο νέφος} (CSPM). Τα κοντέινερ προσθέτουν τις δικές τους απαιτήσεις: ελάχιστα και ελεγμένα είδωλα, εκτέλεση χωρίς δικαιώματα root και πολιτικές δικτύου Kubernetes με άρνηση εξ ορισμού. Κοιτάζοντας μπροστά, οι αμυνόμενοι πρέπει να προετοιμαστούν για τη μετάβαση στη \textbf{μετακβαντική κρυπτογραφία}, για το ηλεκτρονικό ψάρεμα με υποβοήθηση τεχνητής νοημοσύνης και τις απάτες με deepfakes, καθώς και για τη συνέχιση των επιθέσεων στην εφοδιαστική αλυσίδα λογισμικού. Οι αρχές αυτού του βιβλίου —ελάχιστο προνόμιο, άμυνα σε βάθος, επαλήθευση και ανθεκτικότητα— παραμένουν ο πιο αξιόπιστος οδηγός μέσα σε αυτές τις αλλαγές.

\section*{Βασικοί όροι}
\begin{description}[style=nextline,leftmargin=1.2em]
  \item[3-2-1-1-0] Κανόνας αντιγράφων: 3 αντίγραφα, 2 μέσα, 1 εκτός εγκατάστασης, 1 εκτός σύνδεσης/αμετάβλητο, 0 ανεπιβεβαίωτα σφάλματα επαναφοράς.
  \item[RPO / RTO] Μέγιστη αποδεκτή απώλεια δεδομένων (RPO) και μέγιστος αποδεκτός χρόνος διακοπής (RTO).
  \item[Μητρώο κινδύνων] Δομημένη καταγραφή κινδύνων, υπευθύνων, μέτρων, εκτιμήσεων και αποφάσεων αντιμετώπισης.
  \item[ΣΔΑΠ (ISMS)] Σύστημα διαχείρισης ασφάλειας πληροφοριών, όπως ορίζεται στο ISO/IEC 27001.
  \item[Μοντέλο Purdue] Πολυεπίπεδη αρχιτεκτονική αναφοράς που διαχωρίζει τα επίπεδα βιομηχανικού ελέγχου από την επιχειρησιακή πληροφορική.
  \item[Κοινή ευθύνη] Κατανομή των καθηκόντων ασφάλειας μεταξύ ενός παρόχου νέφους και του πελάτη του.
\end{description}

\section*{Σύνοψη κεφαλαίου}
\begin{itemize}
  \item Η ανθεκτικότητα προϋποθέτει την αστοχία: τα αμετάβλητα, δοκιμασμένα αντίγραφα ασφαλείας με ξεχωριστά διαπιστευτήρια αποτελούν τον πυρήνα της ανάκαμψης από λυτρισμικό.
  \item Η επιχειρησιακή συνέχεια και η ανάκαμψη από καταστροφές καθοδηγούνται από τις επιχειρησιακές επιπτώσεις, εκφράζονται μέσω RPO και RTO και είναι αξιόπιστες μόνο όταν δοκιμάζονται.
  \item Η διαχείριση κινδύνου καταγράφει, αξιολογεί και αντιμετωπίζει ρητά κάθε κίνδυνο, με την αποδοχή να αποφασίζεται από υπόλογους υπευθύνους.
  \item Τα NIST CSF 2.0, ISO/IEC 27001 και CIS Controls δομούν τα προγράμματα ασφάλειας· ο ΓΚΠΔ και η NIS2 καθιστούν υποχρεωτικές βασικές πρακτικές.
  \item Η ασφάλεια OT, η μηδενική εμπιστοσύνη και η κοινή ευθύνη στο νέφος επεκτείνουν τις βασικές αρχές του βιβλίου σε νέα περιβάλλοντα.
\end{itemize}

\section*{Ερωτήσεις ανασκόπησης}
\begin{enumerate}
  \item Γιατί ο κανόνας 3-2-1 από μόνος του δεν εγγυάται την ανάκαμψη από λυτρισμικό; Τι προσθέτουν το επιπλέον «1» και το «0»;
  \item Ένα ηλεκτρονικό κατάστημα μπορεί να ανεχθεί απώλεια παραγγελιών 15 λεπτών και διακοπή λειτουργίας 2 ωρών. Μετατρέψτε αυτά σε RPO/RTO και προτείνετε τεχνικό σχεδιασμό.
  \item Ποιος πρέπει να επιτρέπεται να αποδεχθεί έναν υψηλό υπολειπόμενο κίνδυνο και γιατί η αποδοχή πρέπει να τεκμηριώνεται;
  \item Εξηγήστε γιατί η ενεργητική σάρωση ευπαθειών, συνήθης στην πληροφορική, μπορεί να είναι ακατάλληλη σε περιβάλλον OT.
\end{enumerate}

\end{document}
